% Leçon 20 — Expression orale : collectivités locales — Chinois Tle A — Cours signé OnBuch+
% Programme officiel MINESEC. Compiler avec : tectonic ZH20-oral-collectivites-locales.tex
\newcommand{\DOCMATIERE}{Chinois}
\newcommand{\DOCNIVEAU}{Tle A}
\newcommand{\DOCMODULE}{Module 3 : Citoyenneté et ouverture au monde}
\newcommand{\DOCLECON}{ZH20}
\newcommand{\DOCTITRE}{Leçon 20 — Expression orale : collectivités locales}
% OnBuch+ — préambule « cours signés » ENRICHI (compilé avec tectonic / XeLaTeX)
% Système visuel « Pop » d'OnBuch+ : fond crème (jamais blanc), bordures encre
% épaisses, ombres « dures » décalées, titres Archivo Black en capitales.
% Version MATHÉMATIQUES : graphes (pgfplots/tikz), boîtes pédagogiques
% (définition, propriété, méthode, attention, le savais-tu, exercice résolu,
% exercice, TP, bilan), page de garde et signature OnBuch+ ancrée en bas.
%
% Macros attendues dans main.tex AVANT \input{preamble} :
%   \DOCMATIERE  \DOCNIVEAU  \DOCMODULE  \DOCLECON (numéro)  \DOCTITRE
\documentclass[11pt]{article}

\usepackage{fontspec}
\usepackage[french]{babel} % français = langue principale ; le chinois via \zh{..} / textebox (xeCJK)
\usepackage{xeCJK}
\usepackage{pifont} % \ding{51} \ding{55} utilisés par les modèles
\usepackage[a4paper,margin=2cm,top=3.4cm,bottom=2.8cm,headheight=1.4cm,footskip=1.3cm]{geometry}
\usepackage{amsmath,amssymb,amsfonts}
\usepackage{mathtools} % \xmapsto, \coloneqq... souvent utilisés par les modèles
\usepackage{cancel} % simplifications barrées (bilans de forces)
% \abs{z} (module, valeur absolue) et \norm{u} (norme), utilisés par les modèles
\providecommand{\abs}{}\let\abs\relax\DeclarePairedDelimiter{\abs}{\lvert}{\rvert}
\providecommand{\norm}{}\let\norm\relax\DeclarePairedDelimiter{\norm}{\lVert}{\rVert}
\usepackage[table]{xcolor} % [table] : fournit aussi \rowcolors (alternance de lignes)
\usepackage{pagecolor}
\usepackage{tikz}
\usetikzlibrary{arrows.meta,positioning,calc,shapes.geometric,shapes.misc,shapes.arrows,decorations.pathmorphing,decorations.pathreplacing,decorations.markings,patterns,shadows,mindmap,trees,fit,backgrounds,matrix,intersections,angles,quotes}
\usepackage{pgfplots}
\usepackage[version=4]{mhchem} % équations nucléaires \ce{^{235}_{92}U}
\usepackage[european,siunitx]{circuitikz} % schémas électriques
\pgfplotsset{compat=1.18}
\usepgfplotslibrary{fillbetween,groupplots} % aires entre courbes (calcul intégral)
\usepackage{enumitem}
\usepackage{fancyhdr}
% [most] charge toutes les bibliothèques tcolorbox (listings, minted, raster,
% documentation...) et gonfle inutilement la mémoire TeX sur un document long
% (~300 boîtes) : on ne charge que ce qui est réellement utilisé.
\usepackage[skins,breakable]{tcolorbox}
\usepackage{graphicx}
\usepackage{colortbl}
\usepackage{booktabs}
\usepackage{tabularx}
\usepackage{diagbox} % case d'en-tête diagonale des tableaux à double entrée
% Colonnes extensibles centrée / gauche / droite (C, L, R), souvent utilisées par les modèles
\newcolumntype{C}{>{\centering\arraybackslash}X}
\newcolumntype{L}{>{\raggedright\arraybackslash}X}
\newcolumntype{R}{>{\raggedleft\arraybackslash}X}
\usepackage{array}
\usepackage{multicol}
\usepackage{multirow}
\usepackage{siunitx}
% parse-numbers=false : accepte aussi \SI{2,5 \times 10^{-3}}{...}
\sisetup{locale=FR,output-decimal-marker={,},group-minimum-digits=4,parse-numbers=false}
\DeclareSIUnit\ohm{\ensuremath{\symup{\Omega}}} % U+2126 absent de la police
\DeclareSIUnit\division{div} % réglages d'oscilloscope (ms/div, V/div)
\DeclareSIUnit\tr{tr} % tours (vitesse de rotation, ex. \SI{1500}{\tr\per\minute})
\DeclareSIUnit\molar{M} % concentration molaire (ex. \SI{3}{\molar}, \SI{1}{\micro\molar})
\DeclareSIUnit\an{an} % année (le modèle confond parfois avec \year, un compteur TeX) — ex. \SI{5}{\kilo\meter\per\an}
\DeclareSIUnit\FCFA{FCFA} % franc CFA (ex. \SI{500}{\FCFA\per\kilo\gram})
\DeclareSIUnit\degreeBrix{\textdegree Brix} % teneur en sucre d'un sirop/jus (réfractométrie)
\DeclareSIUnit\month{mois} % le modèle utilise parfois \month, un compteur TeX, comme unité de durée
\usepackage{qrcode}
% \degree : commande fréquemment utilisée par les modèles pour un angle ou une
% température en dehors de \SI{}{} (non fournie par les paquets chargés).
\providecommand{\degree}{^{\circ}}
% \qed (fin de démonstration), utilisé par les modèles sans amsthm
\providecommand{\qed}{\hfill$\square$}
% \barre : double barre des valeurs interdites dans un tableau de variations (tkz-tab)
\providecommand{\barre}{\ensuremath{\|}}
% environnement proof (amsthm non chargé) utilisé par les modèles
\ifdefined\proof\else\newenvironment{proof}[1][Démonstration]{\par\noindent\textit{#1.}\ }{\qed\par}\fi
\usepackage{lastpage}
\usepackage{hyperref}
\hypersetup{hidelinks}

% ============================================================
% Palette « Pop » OnBuch+ — mêmes hex que la classe P (plus_ui.dart)
% ============================================================
\definecolor{poporange}{HTML}{F59321}
\definecolor{popdark}{HTML}{E07A0C}
\definecolor{popcream}{HTML}{FFF6E8}
\definecolor{popink}{HTML}{1C1714}
\definecolor{poppurple}{HTML}{7A5AE0}
\definecolor{popgreen}{HTML}{2E9E6A}
\definecolor{poppink}{HTML}{E0457B}
\definecolor{popblue}{HTML}{2F7DE1}
\definecolor{popgold}{HTML}{C98A12}
\definecolor{popmuted}{HTML}{8A7F76}
\definecolor{popline}{HTML}{EADFD2}
\definecolor{popgray}{HTML}{8A7F76}
\colorlet{popgrey}{popgray}
% Alias tolérants (le modèle nomme parfois les couleurs différemment)
\colorlet{popwhite}{white}
\colorlet{popblack}{popink}
\colorlet{popred}{poppink}
\colorlet{popyellow}{popgold}
% Teintes claires (fonds de tableaux/encadrés) — le modèle nomme parfois
% les couleurs avec un suffixe L (ex. popblueL) sans les avoir définies.
\colorlet{poporangeL}{poporange!20}
\colorlet{popdarkL}{popdark!20}
\colorlet{poppurpleL}{poppurple!20}
\colorlet{popgreenL}{popgreen!20}
\colorlet{poppinkL}{poppink!20}
\colorlet{popblueL}{popblue!20}
\colorlet{popgoldL}{popgold!20}
\colorlet{popredL}{popred!20}
\colorlet{popgrayL}{popgray!20}
% Teintes claires pour fonds de tableaux / zones
\colorlet{popblueL}{popblue!10!white}
\colorlet{poporangeL}{poporange!14!white}
\colorlet{popgreenL}{popgreen!12!white}
\colorlet{poppurpleL}{poppurple!12!white}
\colorlet{poppinkL}{poppink!12!white}
\colorlet{popgoldL}{popgold!14!white}

\pagecolor{popcream}

% ============================================================
% Typographie
% ============================================================
\defaultfontfeatures{Ligatures=TeX}
\setmainfont{PlusJakartaSans-Medium.ttf}[Path=../../../fonts/,BoldFont=PlusJakartaSans-Bold.ttf]
% Symboles, API et emojis absents de la police principale : repli sur DejaVu Sans (caractères actifs)
\newfontface\symfont{DejaVuSans.ttf}[Path=../../../fonts/]
\makeatletter
\newcommand\symactive[1]{\catcode#1=13 \begingroup\lccode`\~=#1 \lowercase{\endgroup\def~}{{\symfont\char#1}}}
\newcommand\symblank[1]{\catcode#1=13 \begingroup\lccode`\~=#1 \lowercase{\endgroup\def~}{}}
\newcommand\symhyph[1]{\catcode#1=13 \begingroup\lccode`\~=#1 \lowercase{\endgroup\def~}{\mbox{-}}}
\newcount\symc
\newcommand\symrange[2]{\symc=#1 \@whilenum\symc<#2 \do{\expandafter\symactive\expandafter{\the\symc}\advance\symc by 1 }}
\newcommand\symblankrange[2]{\symc=#1 \@whilenum\symc<#2 \do{\expandafter\symblank\expandafter{\the\symc}\advance\symc by 1 }}
\makeatother
\symrange{"0250}{"02FF}\symactive{"2713}\symactive{"2714}\symactive{"2717}\symactive{"2718}\symactive{"2610}\symactive{"2611}\symactive{"2612}
\symrange{"2600}{"27BF}
\catcode"2705=13 \begingroup\lccode`\~="2705 \lowercase{\endgroup\def~}{{\symfont\char"2713}}\symblank{"26BD}\symblank{"26A1}
\symrange{"2460}{"257F}\symactive{"223C}\symactive{"2248}\symactive{"2260}
\symactive{"01F8}\symactive{"01F9}\symactive{"1E3E}\symactive{"1E3F}
\symhyph{"2011}\symhyph{"2E1A}\symblankrange{"1F300}{"1FAFF}\symblank{"FE0F}
% Caractères chinois : WenQuanYi Zen Hei (sous-ensemble GB2312 + pinyin), gras/italique simulés
\setCJKmainfont{WenQuanYiZenHei-subset.ttf}[Path=../../../fonts/,AutoFakeBold=3,AutoFakeSlant=0.2]
\xeCJKsetup{CJKspace=false}
% Pinyin : voyelles à caron (ǎ ǐ ǒ ǔ ǖ ǘ ǚ ǜ) absentes de la police principale -> caractères actifs
\newfontface\pyfont{WenQuanYiZenHei-subset.ttf}[Path=../../../fonts/]
\newcommand\pyactive[1]{\catcode#1=13 \begingroup\lccode`\~=#1 \lowercase{\endgroup\def~}{{\pyfont\char#1}}}
\pyactive{"01CD}\pyactive{"01CE}\pyactive{"01CF}\pyactive{"01D0}\pyactive{"01D1}\pyactive{"01D2}\pyactive{"01D3}\pyactive{"01D4}\pyactive{"01D5}\pyactive{"01D6}\pyactive{"01D7}\pyactive{"01D8}\pyactive{"01D9}\pyactive{"01DA}\pyactive{"01DB}\pyactive{"01DC}
% Maths en sans-serif (Fira Math) avec chiffres et lettres droites de Plus
% Jakarta, pour que les formules se fondent dans le texte.
\usepackage[math-style=ISO,bold-style=ISO]{unicode-math}
\setmathfont{FiraMath-Regular.otf}[Path=../../../fonts/]
\setmathfont{PlusJakartaSans-Medium.ttf}[Path=../../../fonts/,range={up/num,up/{Latin,latin},\mathup}]
\usepackage{icomma} % virgule décimale sans espace en mode math
% Caractères mathématiques Unicode absents des polices de texte (Plus Jakarta,
% Archivo Black) : rendus par leur équivalent mathématique, en texte comme en
% formule — sinon ils s'affichent en carré vide (ex. « Arithmétique dans ℤ »).
\usepackage{newunicodechar}
\newunicodechar{ℝ}{\ensuremath{\mathbb{R}}}
\newunicodechar{ℤ}{\ensuremath{\mathbb{Z}}}
\newunicodechar{ℂ}{\ensuremath{\mathbb{C}}}
\newunicodechar{ℕ}{\ensuremath{\mathbb{N}}}
\newunicodechar{ℚ}{\ensuremath{\mathbb{Q}}}
\newunicodechar{𝔻}{\ensuremath{\mathbb{D}}}
\newunicodechar{α}{\ensuremath{\alpha}}
\newunicodechar{β}{\ensuremath{\beta}}
\newunicodechar{γ}{\ensuremath{\gamma}}
\newunicodechar{δ}{\ensuremath{\delta}}
\newunicodechar{ε}{\ensuremath{\varepsilon}}
\newunicodechar{θ}{\ensuremath{\theta}}
\newunicodechar{λ}{\ensuremath{\lambda}}
\newunicodechar{μ}{\ensuremath{\mu}}
\newunicodechar{σ}{\ensuremath{\sigma}}
\newunicodechar{τ}{\ensuremath{\tau}}
\newunicodechar{φ}{\ensuremath{\varphi}}
\newunicodechar{ψ}{\ensuremath{\psi}}
\newunicodechar{ω}{\ensuremath{\omega}}
\newunicodechar{Σ}{\ensuremath{\Sigma}}
% majuscules produites par \MakeUppercase dans les titres (α-aminés -> Α)
\newunicodechar{Α}{\ensuremath{\alpha}}
\newunicodechar{Β}{\ensuremath{\beta}}
\newunicodechar{Γ}{\ensuremath{\Gamma}}
\newunicodechar{Δ}{\ensuremath{\Delta}}
\newunicodechar{Ε}{\ensuremath{\varepsilon}}
\newunicodechar{Θ}{\ensuremath{\Theta}}
\newunicodechar{Λ}{\ensuremath{\Lambda}}
\newunicodechar{Μ}{\ensuremath{\mu}}
\newunicodechar{Τ}{\ensuremath{\tau}}
\newunicodechar{Φ}{\ensuremath{\Phi}}
\newunicodechar{Ψ}{\ensuremath{\Psi}}
\newunicodechar{Ω}{\ensuremath{\Omega}}
\newunicodechar{↦}{\ensuremath{\mapsto}}
\newunicodechar{∈}{\ensuremath{\in}}
\newunicodechar{∉}{\ensuremath{\notin}}
\newunicodechar{∧}{\ensuremath{\wedge}}
\newunicodechar{⇔}{\ensuremath{\Leftrightarrow}}
\newunicodechar{⟺}{\ensuremath{\Longleftrightarrow}}
\newunicodechar{⇒}{\ensuremath{\Rightarrow}}
\newunicodechar{⟹}{\ensuremath{\Longrightarrow}}
\newunicodechar{∘}{\ensuremath{\circ}}
\newunicodechar{∪}{\ensuremath{\cup}}
\newunicodechar{∩}{\ensuremath{\cap}}
\newunicodechar{≡}{\ensuremath{\equiv}}
\newunicodechar{⊂}{\ensuremath{\subset}}
\newunicodechar{⊥}{\ensuremath{\perp}}
\newunicodechar{∃}{\ensuremath{\exists}}
\newunicodechar{∀}{\ensuremath{\forall}}
\newunicodechar{∛}{\ensuremath{\sqrt[3]{\,}}}
\newunicodechar{∜}{\ensuremath{\sqrt[4]{\,}}}
\newunicodechar{⃗}{\ensuremath{{}^{\to}}}
\newunicodechar{ⁿ}{\ifmmode^{n}\else\textsuperscript{\ensuremath{n}}\fi}
\newunicodechar{ˣ}{\ifmmode^{x}\else\textsuperscript{\ensuremath{x}}\fi}
\newunicodechar{ᵇ}{\ifmmode^{b}\else\textsuperscript{\ensuremath{b}}\fi}
\newunicodechar{ʳ}{\ifmmode^{r}\else\textsuperscript{\ensuremath{r}}\fi}
\newunicodechar{ᵘ}{\ifmmode^{u}\else\textsuperscript{\ensuremath{u}}\fi}
\newunicodechar{ʸ}{\ifmmode^{y}\else\textsuperscript{\ensuremath{y}}\fi}
\newunicodechar{ᵃ}{\ifmmode^{a}\else\textsuperscript{\ensuremath{a}}\fi}
\newunicodechar{ᵉ}{\ifmmode^{e}\else\textsuperscript{\ensuremath{e}}\fi}
\newunicodechar{ᵏ}{\ifmmode^{k}\else\textsuperscript{\ensuremath{k}}\fi}
\newunicodechar{ᵖ}{\ifmmode^{p}\else\textsuperscript{\ensuremath{p}}\fi}
\newunicodechar{ᴺ}{\ifmmode^{N}\else\textsuperscript{\ensuremath{N}}\fi}
\newunicodechar{ᶜ}{\ifmmode^{c}\else\textsuperscript{\ensuremath{c}}\fi}
\newunicodechar{ʰ}{\ifmmode^{h}\else\textsuperscript{\ensuremath{h}}\fi}
\newunicodechar{ᵠ}{\ifmmode^{q}\else\textsuperscript{\ensuremath{q}}\fi}
\newunicodechar{ₙ}{\ifmmode_{n}\else\textsubscript{\ensuremath{n}}\fi}
\newunicodechar{ₐ}{\ifmmode_{a}\else\textsubscript{\ensuremath{a}}\fi}
\newunicodechar{ᵢ}{\ifmmode_{i}\else\textsubscript{\ensuremath{i}}\fi}
\newunicodechar{ᵣ}{\ifmmode_{r}\else\textsubscript{\ensuremath{r}}\fi}
\newunicodechar{ₖ}{\ifmmode_{k}\else\textsubscript{\ensuremath{k}}\fi}
\newunicodechar{ᵦ}{\ifmmode_{\beta}\else\textsubscript{\ensuremath{\beta}}\fi}
\newunicodechar{ₓ}{\ifmmode_{x}\else\textsubscript{\ensuremath{x}}\fi}
\newfontfamily\popdisplay{ArchivoBlack-Regular.ttf}[Path=../../../fonts/]
\newfontfamily\poplogo{SpaceGrotesk-Bold.ttf}[Path=../../../fonts/]
\newfontfamily\popsemi{PlusJakartaSans-SemiBold.ttf}[Path=../../../fonts/]
\newfontfamily\popxbold{PlusJakartaSans-ExtraBold.ttf}[Path=../../../fonts/]
\newfontfamily\popxboldls{PlusJakartaSans-ExtraBold.ttf}[Path=../../../fonts/,LetterSpace=3.0]

\setlist[enumerate]{leftmargin=1.7em,itemsep=5pt,topsep=4pt}
\setlist[itemize]{leftmargin=1.7em,itemsep=4pt,topsep=4pt,label={\color{poporange}\textbullet}}
\setlength{\parindent}{0pt}
\setlength{\parskip}{5pt}
\renewcommand{\arraystretch}{1.25}

% pgfplots : style Pop par défaut
\pgfplotsset{
  every axis/.append style={axis line style={popink,line width=1pt},tick style={popink},
    grid style={popline,line width=0.5pt},label style={font=\small},tick label style={font=\footnotesize}},
  popcurve/.style={poporange,line width=1.8pt,no markers,smooth},
}
% Même style disponible dans un \draw TikZ simple (hors pgfplots)
\tikzset{popcurve/.style={poporange,line width=1.8pt}}
\tikzset{popgrid/.style={popline,very thin}}
\colorlet{popcurve}{poporange} % le nom du style est parfois utilisé comme couleur

% ============================================================
% Pill / squiggle
% ============================================================
\newcommand{\poppill}[3]{%
  \tikz[baseline=-0.55ex]\node[draw=popink,line width=1.2pt,rounded corners=2.6pt,%
    fill=#2,inner xsep=6.5pt,inner ysep=3pt]{%
    \popxboldls\fontsize{7}{7}\selectfont\color{#3}\MakeUppercase{#1}};%
}
\newcommand{\popsquiggle}[1]{%
  \tikz[baseline=-0.4ex]{
    \draw[#1,line width=2.6pt,line cap=round]
      plot [smooth] coordinates {(0,0.09) (0.34,0.01) (0.68,0.21) (1.02,0.01) (1.36,0.10)};
  }%
}
% Mot-clé surligné (marqueur orange sous le texte)
\newcommand{\cle}[1]{{\popxbold\color{popdark}#1}}

% ============================================================
% En-tête / pied
% ============================================================
\pagestyle{fancy}
\fancyhf{}
\renewcommand{\headrulewidth}{0pt}
\renewcommand{\footrulewidth}{0pt}
\lhead{\raisebox{-0.15ex}{\poplogo\fontsize{13}{13}\selectfont\color{popink}OnBuch\color{poporange}+}}
\rhead{\poppill{\DOCMATIERE\ \textbullet\ \DOCNIVEAU\ \textbullet\ Leçon \DOCLECON}{white}{popink}}
\lfoot{\popxbold\fontsize{7.6}{7.6}\selectfont\color{popmuted}\MakeUppercase{Cours signé OnBuch+}\ \textbullet\ {\popsemi\DOCTITRE}}
\rfoot{\tikz[baseline=-0.6ex]\node[rounded corners=8.5pt,fill=popink,minimum height=17pt,minimum width=17pt,inner xsep=5pt,inner sep=0]{\popxbold\fontsize{8}{8}\selectfont\color{popcream}\thepage\,/\,\pageref*{LastPage}};}

% ============================================================
% Titres
% ============================================================
\newcommand{\chaptitre}[2]{%
  \begin{tcolorbox}[enhanced,colback=poporange,colframe=popink,boxrule=2.6pt,arc=11pt,
      drop shadow=popink,left=15pt,right=15pt,top=12pt,bottom=11pt,before skip=3pt,after skip=18pt]
    \poppill{#2}{popink}{popcream}\par\vspace{9pt}
    {\raggedright\hyphenpenalty=10000\exhyphenpenalty=10000\popdisplay\fontsize{22}{24}\selectfont\color{popcream}\MakeUppercase{#1}\par}
  \end{tcolorbox}
}
% Section numérotée : pastille orange avec le numéro + titre Archivo
\newcounter{popsec}
\newcommand{\coursec}[1]{%
  \stepcounter{popsec}\setcounter{popsub}{0}%
  \par\vspace{16pt}\noindent
  \begin{minipage}{\linewidth}
  \tikz[baseline=-0.7ex]\node[draw=popink,line width=1.6pt,fill=poporange,rounded corners=4pt,minimum size=22pt,inner sep=0]{\popdisplay\fontsize{12}{12}\selectfont\color{popcream}\thepopsec};%
  \hspace{8pt}{\popdisplay\fontsize{15}{17}\selectfont\color{popink}\MakeUppercase{#1}}\par
  \vspace{4pt}\noindent{\color{poporange}\rule{36pt}{3.6pt}}
  \end{minipage}\par\nopagebreak\vspace{8pt}
}
\newcounter{popsub}
\newcommand{\courssub}[1]{%
  \stepcounter{popsub}%
  \par\vspace{9pt}\noindent{\popxbold\fontsize{11.5}{13}\selectfont\color{popink}{\color{poporange}\thepopsec.\thepopsub}\hspace{6pt}#1}\par\nopagebreak\vspace{4pt}
}

% ============================================================
% Boîtes pédagogiques — même recette : fond blanc, bordure épaisse colorée,
% ombre dure encre, badge plein. Titre optionnel [#1].
% ============================================================
\tcbset{popbase/.style={breakable,enhanced,colback=white,boxrule=2.3pt,arc=9pt,
  shadow={3pt}{-3pt}{0pt}{popink},left=14pt,right=14pt,top=11pt,bottom=11pt,before skip=11pt,after skip=15pt,
  fontupper=\popsemi\fontsize{10.6}{14.5}\selectfont\color{popink},
  fontlower=\popsemi\fontsize{10.6}{14.5}\selectfont\color{popink}}}
% Coche dessinée (le glyphe ✓ n'existe pas dans les polices du document)
\newcommand{\popcheck}[1]{\tikz[baseline=-0.5ex]\draw[#1,line width=1.6pt,line cap=round,line join=round](-0.13,0.0)--(-0.04,-0.09)--(0.13,0.1);}
\providecommand{\coche}{\popcheck{popgreen}}
\providecommand{\resultat}[1]{\par\smallskip\noindent\fcolorbox{popgreen}{popgreenL}{\parbox{\dimexpr\linewidth-2\fboxsep-2\fboxrule\relax}{\textbf{Résultat.} #1}}\par\smallskip}
\newcommand{\popboxhead}[3]{\poppill{#1}{#2}{white}\ifx\relax#3\relax\else\hspace{6pt}{\popxbold\fontsize{10.6}{12}\selectfont\color{popink}#3}\fi\par\vspace{7pt}}

\newtcolorbox{aretenir}[1][]{popbase,colframe=popgreen,before upper={\popboxhead{À retenir}{popgreen}{#1}}}
% Texte en chinois (document de lecture, dialogue, extrait) : cadre
% d'étude. Un titre optionnel (source, auteur) s'affiche dans l'en-tête.
\newtcolorbox{textebox}[1][]{popbase,colframe=popblue,before upper={\popboxhead{文本 · Texte}{popblue}{#1}},after upper={}}
% Marques ✓ / ✗ (forme correcte / fautive) souvent utilisées par les modèles
\providecommand{\cross}{\textcolor{poppink}{\ensuremath{\times}}}
\providecommand{\xmark}{\cross}\providecommand{\crossmark}{\cross}\providecommand{\cmark}{\ensuremath{\checkmark}}
% Ligne de réponse / justification à compléter (inventée par les modèles)
\providecommand{\justification}[1]{\par\noindent\textit{Justification~:} \dotfill\par}
% \textipa (package tipa non chargé) : la police ne couvre pas l'API -> simple passage
\providecommand{\textipa}[1]{#1}
% Soulignement (ulem non chargé) : \uline, \ul
\providecommand{\uline}[1]{\underline{#1}}\providecommand{\ul}[1]{\underline{#1}}
% \glossaire{\item ...} inventée par les modèles : liste de glossaire
\providecommand{\glossaire}[1]{\begin{itemize}#1\end{itemize}}
% \alert (beamer) inventée par les modèles : texte mis en évidence
\providecommand{\alert}[1]{\textbf{\textcolor{poppink}{#1}}}
% variantes ASCII d'ordinaux écrites en macro
\providecommand{\iere}{\textsuperscript{re}}\providecommand{\ieme}{\textsuperscript{e}}\providecommand{\ere}{\textsuperscript{re}}\providecommand{\eme}{\textsuperscript{e}}\providecommand{\er}{\textsuperscript{er}}\providecommand{\ie}{\textsuperscript{e}}
% Ordinaux français écrits en macro par les modèles : 1\ière, 3\ième
\newcommand{\ière}{\textsuperscript{re}}\newcommand{\ième}{\textsuperscript{e}}
% \exemple (inventée par les modèles) : étiquette « Exemple : »
\providecommand{\exemple}{\textbf{Exemple~:}\ }
% \square utilisable aussi hors mode math (cases à cocher)
\AtBeginDocument{\let\squareMath\square\renewcommand{\square}{\ensuremath{\squareMath}}}
% Mot ou phrase en chinois à l'intérieur d'un paragraphe français
\newcommand{\zh}[1]{\ifmmode\text{#1}\else{#1}\fi}
\let\eng\zh \let\esp\zh \let\all\zh % alias tolérés
% flèches d'intonation inventées par les modèles
\providecommand{\textsearrow}{}\renewcommand{\textsearrow}{\ensuremath{\searrow}}\providecommand{\textnearrow}{}\renewcommand{\textnearrow}{\ensuremath{\nearrow}}
% \fr{..} : mot français (inventé par les modèles)
\providecommand{\fr}[1]{\textit{#1}}
% zhbox : encadré de texte chinois inventé par les modèles
\newenvironment{zhbox}{\begin{quote}}{\end{quote}}
% \vs : « versus » inventé par les modèles
\providecommand{\vs}{\textit{vs}\ }
% \blank : case à compléter inventée par les modèles
\providecommand{\blank}[1][]{\underline{\hspace{1.6em}}}
\newtcolorbox{exemplebox}[1][]{popbase,colframe=poporange,before upper={\popboxhead{Exemple}{poporange}{#1}}}
\newtcolorbox{definition}[1][]{popbase,colframe=popblue,before upper={\popboxhead{Définition}{popblue}{#1}}}
\newtcolorbox{propriete}[1][]{popbase,colframe=poppurple,before upper={\popboxhead{Propriété}{poppurple}{#1}}}
\newtcolorbox{methode}[1][]{popbase,colframe=popgold,before upper={\popboxhead{Méthode}{popgold}{#1}}}
\newtcolorbox{attention}[1][]{popbase,colframe=poppink,before upper={\popboxhead{Attention}{poppink}{#1}}}
\newtcolorbox{experience}[1][]{popbase,colframe=popblue,colback=popblueL,before upper={\popboxhead{Expérience}{popblue}{#1}}}
% « Le savais-tu ? » : carte violette pleine (contexte, culture, Cameroun)
\newtcolorbox{savaistu}[1][]{popbase,colback=poppurple,colframe=popink,
  fontupper=\popsemi\fontsize{10.4}{14.2}\selectfont\color{white},
  before upper={\poppill{Le savais-tu\,?}{white}{poppurple}\ifx\relax#1\relax\else\hspace{6pt}{\popxbold\color{white}#1}\fi\par\vspace{7pt}}}
% Exercice résolu : énoncé en haut, \tcblower puis la solution sur fond crème
\newcounter{popexo}\newcounter{popexr}
\newtcolorbox{exoresolu}[1][]{popbase,colframe=poporange,colbacklower=popcream,
  segmentation style={popink,line width=1.2pt,dashed},
  before upper={\stepcounter{popexr}\popboxhead{Exercice résolu \thepopexr}{poporange}{#1}},
  before lower={\poppill{Solution}{popink}{popcream}\par\vspace{6pt}}}
% Exercice (non résolu en place) — la correction est dans la partie Corrigés
\newtcolorbox{exercice}[1][]{popbase,colframe=popink,
  before upper={\stepcounter{popexo}\popboxhead{Exercice \thepopexo}{popink}{#1}}}
% Corrigé
\newtcolorbox{corrige}[1][]{popbase,colframe=popgreen,colback=popgreenL,
  before upper={\popboxhead{Corrigé}{popgreen}{#1}}}
% Objectifs / prérequis / bilan
\newtcolorbox{objectifs}{popbase,colframe=popink,colback=poporangeL,
  before upper={\popboxhead{Objectifs de la leçon}{popink}{}}}
\newtcolorbox{prerequis}{popbase,colframe=popink,colback=popblueL,
  before upper={\popboxhead{Prérequis}{popblue}{}}}
\newtcolorbox{bilan}{popbase,colframe=popink,colback=white,boxrule=2.8pt,
  before upper={\popboxhead{Fiche bilan}{poporange}{L'essentiel à connaître pour l'examen}}}
% Figure encadrée avec légende
\newtcolorbox{popfigure}[1][]{enhanced,breakable=false,colback=white,colframe=popline,boxrule=1.4pt,arc=8pt,
  left=10pt,right=10pt,top=10pt,bottom=8pt,before skip=10pt,after skip=12pt,halign=center,
  lower separated=false,halign lower=center,
  fontlower=\popsemi\fontsize{9}{11.5}\selectfont\color{popmuted},
  #1}
\newcommand{\legende}[1]{\par\vspace{6pt}{\popsemi\fontsize{9}{11.5}\selectfont\color{popmuted}#1}}

% ============================================================
% Signature OnBuch+
% ============================================================
\newcommand{\poplogoinline}[1]{{\poplogo\fontsize{#1}{#1}\selectfont\color{popink}OnBuch\color{poporange}+}}
\newcommand{\signatureonbuch}{%
  \begin{tcolorbox}[enhanced,colback=popink,colframe=popink,boxrule=2.2pt,arc=10pt,
      drop shadow=poporange,left=14pt,right=14pt,top=10pt,bottom=10pt,before skip=0pt,after skip=0pt]
    \tikz[baseline=-0.7ex]\node[circle,fill=popgreen,draw=popcream,line width=1.3pt,minimum size=22pt,inner sep=0]{\popcheck{white}};%
    \hspace{9pt}\begin{minipage}[c]{0.62\linewidth}
      {\popxbold\fontsize{12}{13}\selectfont\color{popcream}Signé\ \textbullet\ {\poplogo OnBuch\color{poporange}+}}\par\vspace{2pt}
      {\raggedright\popsemi\fontsize{8}{10.5}\selectfont\color{popline}\MakeUppercase{Équipe pédagogique OnBuch+}\\\MakeUppercase{Conforme au programme officiel MINESEC}\par}
    \end{minipage}\hfill
    {\poplogo\fontsize{22}{22}\selectfont\color{popcream}OnBuch\color{poporange}+}
  \end{tcolorbox}%
}

% ============================================================
% Page de garde
% #1 titre · #2 matière · #3 niveau · #4 module · #5 n° leçon · #6 durée ·
% #7 description / objectifs (texte ou itemize)
% ============================================================
\newcommand{\pagegarde}[7]{%
  \thispagestyle{empty}
  \newgeometry{a4paper,margin=1.7cm,top=1.6cm,bottom=1.4cm}%
  \begin{tikzpicture}[remember picture,overlay]
    \fill[poporange!18] ([xshift=-3.2cm,yshift=-3.6cm]current page.north east) circle (4.6cm);
    \fill[poppurple!14] ([xshift=2.4cm,yshift=7.6cm]current page.south west) circle (3.8cm);
  \end{tikzpicture}%
  {\poplogo\fontsize{30}{30}\selectfont\color{popink}OnBuch\color{poporange}+}\hfill
  \raisebox{0.6ex}{\poppill{Cours signé \textbullet\ Premium}{popink}{popcream}}\par
  \vspace{1pt}\popsquiggle{poppurple}\par
  \vspace{8pt}
  \begin{tcolorbox}[enhanced,colback=poporange,colframe=popink,boxrule=2.8pt,arc=13pt,
      drop shadow=popink,left=18pt,right=18pt,top=12pt,bottom=13pt,after skip=10pt]
    \poppill{#2 \textbullet\ #3}{popink}{popcream}\hspace{5pt}\poppill{#4}{white}{popink}\par\vspace{8pt}
    {\popdisplay\fontsize{14}{14}\selectfont\color{popink}LEÇON #5}\par\vspace{4pt}
    {\raggedright\hyphenpenalty=10000\exhyphenpenalty=10000\popdisplay\fontsize{25}{27}\selectfont\color{popcream}\MakeUppercase{#1}\par}
    \vspace{8pt}
    {\popxbold\fontsize{9.5}{9.5}\selectfont\color{popink}\MakeUppercase{Durée conseillée : #6}}
  \end{tcolorbox}
  \begin{tcolorbox}[enhanced,colback=white,colframe=popink,boxrule=2.2pt,arc=10pt,
      drop shadow=popink,left=16pt,right=16pt,top=10pt,bottom=10pt,after skip=8pt]
    \poppill{Dans cette leçon}{poppurple}{white}\par\vspace{5pt}
    \popsemi\fontsize{9.3}{12.6}\selectfont\color{popink}#7
  \end{tcolorbox}
  \noindent\begin{minipage}[t]{0.487\linewidth}
    \begin{tcolorbox}[enhanced,colback=popgreen,colframe=popink,boxrule=2.2pt,arc=10pt,
        drop shadow=popink,left=11pt,right=11pt,top=10pt,bottom=10pt,height=3.1cm,valign=center]
      \tcbox[colback=white,colframe=popink,boxrule=1.3pt,arc=5pt,boxsep=0pt,nobeforeafter,
        left=3pt,right=3pt,top=3pt,bottom=3pt,valign=center]{\qrcode[height=1.9cm]{https://play.google.com/store/apps/details?id=com.luvvix.onbuch&hl=fr}}%
      \hspace{8pt}\begin{minipage}[c]{0.5\linewidth}
        {\popdisplay\fontsize{11}{12.5}\selectfont\color{white}\MakeUppercase{Télécharge OnBuch}}\par\vspace{4pt}
        {\raggedright\popsemi\fontsize{8.4}{11}\selectfont\color{white}Gratuit \textbullet\ Google Play\\Tuteur IA, annales, concours, résultats\par}
      \end{minipage}
    \end{tcolorbox}
  \end{minipage}\hfill
  \begin{minipage}[t]{0.487\linewidth}
    \begin{tcolorbox}[enhanced,colback=poppurple,colframe=popink,boxrule=2.2pt,arc=10pt,
        drop shadow=popink,left=11pt,right=11pt,top=10pt,bottom=10pt,height=3.1cm,valign=center]
      \tikz[baseline=-0.7ex]\node[circle,fill=white,draw=popink,line width=1.3pt,minimum size=1.6cm,inner sep=0]{\popdisplay\fontsize{24}{24}\selectfont\color{poppurple}+};%
      \hspace{8pt}\begin{minipage}[c]{0.6\linewidth}
        {\popdisplay\fontsize{11}{12.5}\selectfont\color{white}\MakeUppercase{Passe à OnBuch+}}\par\vspace{4pt}
        {\raggedright\popsemi\fontsize{8.4}{11}\selectfont\color{white}Tous les cours signés, Tuteur IA illimité, fiches PDF — onglet OnBuch+ de l'app\par}
      \end{minipage}
    \end{tcolorbox}
  \end{minipage}
  \vspace{8pt}
  \popcontenu
  \vfill
  \signatureonbuch
  \restoregeometry
  \newpage
}

% Rangée « Ce cours contient » (page de garde)
\newcommand{\poptile}[3]{% #1 couleur #2 gros texte #3 légende
  \begin{tcolorbox}[enhanced,colback=#1,colframe=popink,boxrule=2pt,arc=9pt,shadow={2.5pt}{-2.5pt}{0pt}{popink},
      left=6pt,right=6pt,top=9pt,bottom=9pt,halign=center,valign=center,height=1.75cm,width=\linewidth,nobeforeafter]
    {\popdisplay\fontsize{12.5}{13}\selectfont\color{white}\MakeUppercase{#2}}\par\vspace{3pt}
    {\centering\popsemi\fontsize{7.8}{9.5}\selectfont\color{white}#3\par}
  \end{tcolorbox}}
\newcommand{\popcontenu}{%
  \par\noindent{\popxboldls\fontsize{7.5}{7.5}\selectfont\color{popmuted}CE COURS SIGNÉ CONTIENT}\par\vspace{7pt}
  \noindent\begin{minipage}[t]{0.235\linewidth}\poptile{popblue}{Cours}{complet, illustré, pas à pas}\end{minipage}\hfill
  \begin{minipage}[t]{0.235\linewidth}\poptile{poporange}{Méthodes}{et exercices résolus}\end{minipage}\hfill
  \begin{minipage}[t]{0.235\linewidth}\poptile{poppink}{Exercices}{type examen + corrigés}\end{minipage}\hfill
  \begin{minipage}[t]{0.235\linewidth}\poptile{popgreen}{Fiche bilan}{l'essentiel à retenir}\end{minipage}\par}

% Sommaire
\newtcolorbox{sommaire}{enhanced,colback=white,colframe=popink,boxrule=2.2pt,arc=10pt,
  drop shadow=popink,left=18pt,right=18pt,top=13pt,bottom=13pt,before skip=6pt,after skip=20pt,
  before upper={\poppill{Sommaire}{poppurple}{white}\par\vspace{9pt}\popsemi\fontsize{10.6}{14.5}\selectfont\color{popink}}}

% Signature de fin de cours — poussée tout en bas de la dernière page
\newcommand{\findecours}{%
  \par\vfill\nopagebreak
  \begin{center}\popsquiggle{poporange}\end{center}\vspace{4pt}
  \signatureonbuch
}
\AtBeginDocument{\def\llbracket{\mathopen{[\mkern-3.5mu[}}\def\rrbracket{\mathclose{]\mkern-3.5mu]}}} % intervalles d'entiers (glyphe ⟦ absent de la police)
% Glyphes absents de Fira Math : redéfinis à partir de glyphes disponibles
\AtBeginDocument{%
  \def\setminus{\mathbin{\backslash}}\let\smallsetminus\setminus
  \def\vdots{\mathord{\vbox{\baselineskip3.2pt\lineskiplimit0pt\kern4pt\hbox{.}\hbox{.}\hbox{.}}}}%
  \def\ddots{\mathinner{\mkern1mu\raise6.5pt\hbox{.}\mkern2mu\raise3.6pt\hbox{.}\mkern2mu\raise0.7pt\hbox{.}\mkern1mu}}%
  \def\triangle{\mathord{\scalebox{0.9}{$\symup{\Delta}$}}}%
  \def\nabla{\mathord{\raisebox{\depth}{\scalebox{1}[-1]{$\symup{\Delta}$}}}}%
  \def\checkmark{\popcheck{popdark}}%
  \def\star{\mathbin{\ast}}\def\bigstar{\mathord{\ast}}%
  \def\diamond{\mathbin{\scalebox{0.6}{$\Diamond$}}}%
  \def\bigcup{\mathop{\mathchoice{\raisebox{-0.3em}{\scalebox{1.7}{$\cup$}}}{\scalebox{1.25}{$\cup$}}{\cup}{\cup}}\displaylimits}%
  \def\bigcap{\mathop{\mathchoice{\raisebox{-0.3em}{\scalebox{1.7}{$\cap$}}}{\scalebox{1.25}{$\cap$}}{\cap}{\cap}}\displaylimits}%
  \def\lesseqqgtr{\mathrel{\lessgtr}}\def\bigoplus{\mathop{\oplus}\displaylimits}%
}
\providecommand{\ssi}{si et seulement si} % abréviation fréquente
\AtBeginDocument{\providecommand{\wideparen}{\overparen}} % arc de cercle

%% FIGFIX-BEGIN v4 — correctifs globaux des figures (docs/onbuch-generation/tools/figfix)
\usepackage{adjustbox}\usepackage{etoolbox}
% toute figure TikZ/pgfplots plus large que la ligne est réduite à la largeur disponible
\BeforeBeginEnvironment{tikzpicture}{\begin{adjustbox}{max width=\linewidth}}
\AfterEndEnvironment{tikzpicture}{\end{adjustbox}}
% légendes pgfplots lisibles par-dessus les courbes
\pgfplotsset{every axis/.append style={legend style={fill=white,fill opacity=0.92,text opacity=1,draw=black!15,inner sep=3pt}}}
% étiquettes d'arêtes (syntaxe edge["..."]) avec fond blanc
\tikzset{every edge quotes/.append style={fill=white,inner sep=1.2pt}}
%% FIGFIX-END
\begin{document}

\pagegarde{Leçon 20 — Expression orale : collectivités locales}{Chinois}{Tle A}{Module 3 : Citoyenneté et ouverture au monde}{ZH20}{2 h}{Cette leçon d'expression orale entraîne les élèves de Terminale A à présenter un exposé simple et à participer à un débat guidé en chinois sur le thème des collectivités locales (communes, conseils municipaux, services publics). Elle fournit le lexique essentiel, les formules fonctionnelles pour présenter, opiner et relancer, ainsi qu'un travail ciblé sur les tons.
\vspace{4pt}
\begin{multicols}{2}\small
\begin{itemize}
\item À la fin de cette leçon, je sais utiliser le vocabulaire des collectivités locales (commune, maire, conseil, services, projets)
\item À la fin de cette leçon, je sais présenter oralement un exposé simple et structuré en chinois (introduction, développement, conclusion)
\item À la fin de cette leçon, je sais donner mon avis avec les formules 我觉得…, 我认为…, 在我看来…
\item À la fin de cette leçon, je sais exprimer l'accord et le désaccord poliment (我同意 / 我不太同意，因为…)
\item À la fin de cette leçon, je sais relancer un débat et poser des questions (你觉得呢？你有什么看法？)
\item À la fin de cette leçon, je sais prononcer correctement les tons des mots clés du thème, notamment les suites de 3e tons
\end{itemize}\end{multicols}}

\begin{sommaire}
\begin{multicols}{2}
\begin{enumerate}
\item Lexique des collectivités locales
\item Dialogue modèle 1 : présenter sa commune
\item Formules fonctionnelles : opiner, réagir, relancer
\item Dialogue modèle 2 : débat sur un projet municipal
\item Travail phonétique : les tons du thème
\item Jeux de rôle et débat guidé
\item Activité d'intégration
\item Méthodes et exercices résolus
\item Exercices
\item Corrigés des exercices
\item Fiche bilan
\end{enumerate}
\end{multicols}
\end{sommaire}

\begin{objectifs}
\begin{itemize}
\item À la fin de cette leçon, je sais utiliser le vocabulaire des collectivités locales (commune, maire, conseil, services, projets)
\item À la fin de cette leçon, je sais présenter oralement un exposé simple et structuré en chinois (introduction, développement, conclusion)
\item À la fin de cette leçon, je sais donner mon avis avec les formules 我觉得…, 我认为…, 在我看来…
\item À la fin de cette leçon, je sais exprimer l'accord et le désaccord poliment (我同意 / 我不太同意，因为…)
\item À la fin de cette leçon, je sais relancer un débat et poser des questions (你觉得呢？你有什么看法？)
\item À la fin de cette leçon, je sais prononcer correctement les tons des mots clés du thème, notamment les suites de 3e tons
\item À la fin de cette leçon, je sais comparer brièvement les collectivités locales du Cameroun et de la Chine
\item À la fin de cette leçon, je sais gérer un jeu de rôle (maire, journaliste, citoyen) dans un dialogue simulé
\end{itemize}
\end{objectifs}

\begin{prerequis}
\begin{itemize}
\item Structures de base de Première : phrase sujet-verbe-complément, 是 / 有 / 在
\item Lexique de la ville et des lieux publics (学校, 医院, 市场, 路)
\item Formules de présentation personnelle et de goût (我叫…, 我喜欢…)
\item Nombres, mesures et expressions de quantité pour citer des chiffres simples
\item Maîtrise des quatre tons et du ton neutre acquise en Seconde/Première
\end{itemize}
\end{prerequis}

\newpage

\coursec{Leçon 20 — Expression orale : collectivités locales}

\textbf{Situation d'accroche.} Imagine : un délégué chinois de la ville de Suzhou visite ta commune — la mairie de Bafoussam, de Yaoundé ou de Bamenda. Le maire te choisit, toi l'élève de Terminale, pour présenter ta collectivité locale en chinois : ses services, ses projets, ses problèmes. Sauras-tu dire « notre commune construit un marché » ou « je pense que le conseil municipal doit d'abord réparer les routes » ? C'est exactement la compétence exigée à l'épreuve orale du baccalauréat.

\textbf{Problématique.} Pour parler de sa collectivité locale en chinois, il faut d'abord maîtriser un lexique précis : les institutions, les services publics, les actions de la collectivité et les habitants. Comment organiser ce vocabulaire pour le retenir et l'employer correctement à l'oral ?

\courssub{Lexique des collectivités locales}

\courssub{Observer : un premier texte sur la collectivité locale}

Avant toute liste de mots, lis ce court texte. Souligne mentalement les mots qui désignent des lieux, des personnes et des actions.

\begin{textebox}[Notre quartier]
我们的社区有学校、医院和市场。\\
Wǒmen de shèqū yǒu xuéxiào, yīyuàn hé shìchǎng.\\
« Notre quartier a une école, un hôpital et un marché. »\\
市政府计划修一条新路。\\
Shì zhèngfǔ jìhuà xiū yì tiáo xīn lù.\\
« La mairie prévoit de construire une nouvelle route. »\\
居民们很高兴，因为交通会更方便。\\
Jūmínmen hěn gāoxìng, yīnwèi jiāotōng huì gèng fāngbiàn.\\
« Les habitants sont contents, car la circulation sera plus pratique. »
\end{textebox}

\textbf{Observations guidées :}
\begin{itemize}
\item \zh{社区} désigne le lieu où l'on vit : le \cle{quartier}, la \cle{communauté}.
\item \zh{学校}, \zh{医院}, \zh{市场} sont des \cle{services} offerts aux habitants.
\item \zh{市政府} est l'\cle{institution} qui décide ; \zh{计划} et \zh{修} sont ses \cle{actions}.
\item \zh{居民们} : le suffixe \zh{们} marque le pluriel des personnes, comme dans \zh{我们} (nous).
\end{itemize}

\begin{definition}[社区 : le mot central du thème]
\zh{社区} (\emph{shèqū}) signifie « communauté, quartier ». En Chine, c'est l'unité administrative de base en ville, comparable au quartier ou à la commune au Cameroun. Le caractère \zh{社} évoque le groupe social (on le retrouve dans \zh{社会}, la société) et \zh{区} signifie « zone, district ».
\end{definition}

\courssub{Les institutions : qui gouverne la collectivité ?}

Voici le premier groupe de mots : ceux qui nomment la ville et ses dirigeants.

\begin{center}
\begin{tabularx}{\linewidth}{|X|X|X|X|}
\hline
\rowcolor{popblueL} Caractères & Pinyin & Nature & Traduction \\
\hline
城市 & chéngshì & nom & la ville \\
\hline
区 & qū & nom & le district, l'arrondissement \\
\hline
社区 & shèqū & nom & la communauté, le quartier \\
\hline
市政府 & shì zhèngfǔ & nom & la mairie, le gouvernement municipal \\
\hline
市长 & shìzhǎng & nom & le maire \\
\hline
地方 & dìfang & nom / adjectif & le lieu ; local \\
\hline
\end{tabularx}
\end{center}

\textbf{Lecture des caractères.} Le caractère \zh{市} (\emph{shì}, 4\up{e} ton) est la clé de ce groupe : il signifie « ville, municipal ». On le retrouve dans \zh{城市} (ville), \zh{市政府} (gouvernement de la ville = mairie), \zh{市长} (chef de la ville = maire). Le caractère \zh{长} se prononce ici \emph{zhǎng} (3\up{e} ton) et signifie « chef, aîné » ; attention, dans \zh{很长} (très long), il se prononce \emph{cháng} (2\up{e} ton) : c'est un caractère à deux lectures (\cle{多音字}).

\textbf{Exemples :}

\zh{雅温得是一个大城市。} \emph{Yǎwēndé shì yí ge dà chéngshì.} « Yaoundé est une grande ville. »

\zh{我们区的市长很年轻。} \emph{Wǒmen qū de shìzhǎng hěn niánqīng.} « Le maire de notre arrondissement est jeune. »

\zh{这是地方政府的问题。} \emph{Zhè shì dìfang zhèngfǔ de wèntí.} « C'est un problème des autorités locales. »

\courssub{Les services publics : que trouve-t-on dans la commune ?}

\begin{center}
\begin{tabularx}{\linewidth}{|X|X|X|X|}
\hline
\rowcolor{popblueL} Caractères & Pinyin & Nature & Traduction \\
\hline
服务 & fúwù & nom / verbe & le service ; servir \\
\hline
医院 & yīyuàn & nom & l'hôpital \\
\hline
学校 & xuéxiào & nom & l'école \\
\hline
市场 & shìchǎng & nom & le marché \\
\hline
公园 & gōngyuán & nom & le parc \\
\hline
路 & lù & nom & la route, le chemin \\
\hline
\end{tabularx}
\end{center}

\textbf{Astuces de mémorisation.} \zh{医院} contient \zh{医} (médecine, comme dans \zh{医生}, médecin) ; \zh{学校} contient \zh{学} (étudier, comme dans \zh{学生}, élève) ; \zh{公园} contient \zh{公} (public) et \zh{园} (jardin) : littéralement « jardin public ». Pour dire « une route », on utilise le classificateur \zh{条} : \zh{一条路} (\emph{yì tiáo lù}), car \zh{条} sert pour les objets longs et étroits (routes, rivières, pantalons).

\textbf{Exemples :}

\zh{社区里有一个新公园。} \emph{Shèqū lǐ yǒu yí ge xīn gōngyuán.} « Il y a un nouveau parc dans le quartier. »

\zh{这个市场很大，也很热闹。} \emph{Zhège shìchǎng hěn dà, yě hěn rènao.} « Ce marché est grand et très animé. »

\begin{attention}[城市 ou 市场 ?]
Ne confonds pas \zh{城市} (\emph{chéngshì}, la ville) et \zh{市场} (\emph{shìchǎng}, le marché) : ce sont les mêmes caractères, dans l'ordre inversé, avec des sens différents ! Les tons changent aussi : \emph{chéngshì} (2\up{e}+4\up{e}) contre \emph{shìchǎng} (4\up{e}+3\up{e}). À l'oral du baccalauréat, cette inversion est une erreur classique des francophones, qui retiennent les caractères « en bloc » sans fixer leur ordre.
\end{attention}

\courssub{Les actions de la collectivité : que fait la mairie ?}

\begin{center}
\begin{tabularx}{\linewidth}{|X|X|X|X|}
\hline
\rowcolor{popblueL} Caractères & Pinyin & Nature & Traduction \\
\hline
建设 & jiànshè & verbe & construire, développer \\
\hline
修 & xiū & verbe & réparer, construire (route) \\
\hline
管理 & guǎnlǐ & verbe & gérer, administrer \\
\hline
解决 & jiějué & verbe & résoudre \\
\hline
问题 & wèntí & nom & le problème, la question \\
\hline
计划 & jìhuà & nom / verbe & le plan, le projet ; prévoir \\
\hline
\end{tabularx}
\end{center}

\textbf{Collocations à retenir.} En chinois, un verbe s'associe à des noms précis ; apprends les paires, pas les mots isolés :
\begin{itemize}
\item \zh{解决问题} (\emph{jiějué wèntí}) : résoudre un problème ;
\item \zh{修路} (\emph{xiū lù}) : réparer ou construire une route ;
\item \zh{建设城市} (\emph{jiànshè chéngshì}) : développer la ville ;
\item \zh{管理市场} (\emph{guǎnlǐ shìchǎng}) : gérer le marché ;
\item \zh{有计划} (\emph{yǒu jìhuà}) : avoir un projet ; \zh{市政府计划……} : la mairie prévoit de……
\end{itemize}

\textbf{Exemples :}

\zh{市政府计划建设一个新医院。} \emph{Shì zhèngfǔ jìhuà jiànshè yí ge xīn yīyuàn.} « La mairie prévoit de construire un nouvel hôpital. »

\zh{我们要先解决水的问题。} \emph{Wǒmen yào xiān jiějué shuǐ de wèntí.} « Nous devons d'abord résoudre le problème de l'eau. »

\zh{谁管理这个公园？} \emph{Shéi guǎnlǐ zhège gōngyuán?} « Qui gère ce parc ? »

\courssub{Les habitants : pour qui travaille la collectivité ?}

\begin{center}
\begin{tabularx}{\linewidth}{|X|X|X|X|}
\hline
\rowcolor{popblueL} Caractères & Pinyin & Nature & Traduction \\
\hline
居民 & jūmín & nom & les habitants, les résidents \\
\hline
市民 & shìmín & nom & les citoyens, les citadins \\
\hline
生活 & shēnghuó & nom / verbe & la vie ; vivre \\
\hline
环境 & huánjìng & nom & l'environnement, le cadre de vie \\
\hline
\end{tabularx}
\end{center}

\textbf{Nuance utile à l'oral.} \zh{居民} désigne les personnes qui \emph{résident} dans un lieu (un quartier, un immeuble) ; \zh{市民} désigne les habitants de la \emph{ville} en tant que citoyens. Pour parler de ton quartier, dis \zh{居民} ; pour parler des habitants de Douala en général, dis \zh{市民}.

\textbf{Exemples :}

\zh{居民们希望有更好的环境。} \emph{Jūmínmen xīwàng yǒu gèng hǎo de huánjìng.} « Les habitants espèrent un meilleur environnement. »

\zh{市民的生活越来越方便。} \emph{Shìmín de shēnghuó yuè lái yuè fāngbiàn.} « La vie des citadins devient de plus en plus pratique. »

\courssub{Carte mentale : organiser le lexique en quatre familles}

\begin{popfigure}
\begin{tikzpicture}[mindmap, grow cyclic, every node/.style={concept}, concept color=popblue!60, text=white,
level 1/.append style={level distance=3.2cm, sibling angle=90},
level 2/.append style={level distance=2.2cm}]
\node[align=center]{社区\\ collectivité locale}
child[concept color=poporange!80] { node[align=center]{机构\\ institutions} }
child[concept color=popgreen!70] { node[align=center]{服务\\ services} }
child[concept color=poppurple!70] { node[align=center]{行动\\ actions} }
child[concept color=poppink!80] { node[align=center]{人\\ habitants} };
\end{tikzpicture}
\legende{Carte mentale du lexique : quatre familles autour du mot central 社区. Branche 机构 : 城市、区、市政府、市长 ; branche 服务 : 医院、学校、市场、公园、路 ; branche 行动 : 建设、修、管理、解决、计划 ; branche 人 : 居民、市民、生活、环境.}
\end{popfigure}

\begin{methode}[Classer pour mémoriser]
Pour retenir durablement ce lexique, procède ainsi :
\begin{enumerate}
\item Prends chaque mot nouveau et classe-le dans l'une des quatre catégories : \cle{institution}, \cle{service}, \cle{action}, \cle{personne}.
\item Pour chaque catégorie, construis une phrase simple qui te concerne : \zh{我们社区有一个市场。}
\item Récite les quatre familles à voix haute, en respectant les tons, puis fais-toi interroger par un camarade.
\item Le jour de l'exposé, structure ta présentation selon ces quatre familles : c'est un plan tout prêt !
\end{enumerate}
\end{methode}

\courssub{Prononciation ciblée : les tons du thème}

Trois mots du thème demandent un entraînement particulier :

\begin{center}
\begin{tabularx}{\linewidth}{|X|X|X|}
\hline
\rowcolor{popblueL} Mot & Tons & Conseil de prononciation \\
\hline
市长 shìzhǎng & 4\up{e} + 3\up{e} & Attaque forte et descendante sur \emph{shì}, puis descends et remonte sur \emph{zhǎng}. \\
\hline
居民 jūmín & 1\up{er} + 2\up{e} & Tiens \emph{jū} haut et plat, puis monte sur \emph{mín} comme une question. \\
\hline
环境 huánjìng & 2\up{e} + 4\up{e} & Monte sur \emph{huán}, puis tombe net sur \emph{jìng}. \\
\hline
\end{tabularx}
\end{center}

Exerce-toi sur la phrase : \zh{市长和居民谈环境问题。} \emph{Shìzhǎng hé jūmín tán huánjìng wèntí.} « Le maire discute des problèmes d'environnement avec les habitants. » Elle enchaîne les trois mots difficiles.

\courssub{Correspondances Cameroun -- Chine}

\begin{center}
\begin{tabularx}{\linewidth}{|X|X|X|}
\hline
\rowcolor{popblueL} Chinois & Chine & Cameroun \\
\hline
市政府 shì zhèngfǔ & gouvernement municipal & la mairie, le conseil municipal \\
\hline
市长 shìzhǎng & le maire (nommé dans un système à double direction) & le maire, élu à la tête du conseil municipal \\
\hline
社区 shèqū & unité administrative de base en ville & le quartier, la communauté urbaine \\
\hline
区 qū & district d'une grande ville & l'arrondissement \\
\hline
\end{tabularx}
\end{center}

\begin{savaistu}[Qui dirige une ville chinoise ?]
En Chine, le maire (\zh{市长}) travaille avec un secrétaire du Parti communiste, qui est souvent la première autorité de la ville. Au Cameroun, le maire est élu à la tête du conseil municipal, et certaines grandes villes ont aussi un délégué du gouvernement. Citer cette différence dans ton exposé oral montre au jury que tu compares les deux cultures : un vrai point fort !
\end{savaistu}

\begin{exoresolu}[Classer et compléter]
\textbf{Énoncé.} 1) Classe les mots suivants dans les quatre familles (institution, service, action, personne) : \zh{市长}, \zh{医院}, \zh{解决}, \zh{居民}, \zh{公园}, \zh{建设}. 2) Complète : \zh{市政府计划\_\_\_一条新路。}
\tcblower
\textbf{Solution.} 1) Institution : \zh{市长} (le maire). Services : \zh{医院} (hôpital), \zh{公园} (parc). Actions : \zh{解决} (résoudre), \zh{建设} (construire). Personne : \zh{居民} (habitants). 2) On dit \zh{修路} pour une route : \zh{市政府计划修一条新路。} \emph{Shì zhèngfǔ jìhuà xiū yì tiáo xīn lù.} « La mairie prévoit de construire une nouvelle route. » On n'oublie pas le classificateur \zh{条} devant \zh{路}.
\end{exoresolu}

\begin{exercice}[À toi de jouer]
Traduis en chinois : 1) « Notre quartier a une école et un marché. » 2) « Le maire veut résoudre ce problème. » 3) « Les habitants de la ville aiment ce parc. » Puis classe chaque mot chinois utilisé dans l'une des quatre familles.
\end{exercice}

\begin{corrige}[Exercice « À toi de jouer »]
1) \zh{我们的社区有学校和市场。} \emph{Wǒmen de shèqū yǒu xuéxiào hé shìchǎng.} 2) \zh{市长想解决这个问题。} \emph{Shìzhǎng xiǎng jiějué zhège wèntí.} 3) \zh{市民们喜欢这个公园。} \emph{Shìmínmen xǐhuan zhège gōngyuán.} Familles : \zh{社区} institution/lieu ; \zh{学校}, \zh{市场}, \zh{公园} services ; \zh{解决} action ; \zh{市长}, \zh{市民} personnes.
\end{corrige}

\begin{aretenir}[L'essentiel de la section]
\begin{itemize}
\item Le lexique des collectivités locales s'organise en quatre familles : institutions (\zh{城市}, \zh{市政府}, \zh{市长}), services (\zh{医院}, \zh{学校}, \zh{市场}, \zh{公园}, \zh{路}), actions (\zh{建设}, \zh{修}, \zh{管理}, \zh{解决}, \zh{计划}), habitants (\zh{居民}, \zh{市民}, \zh{生活}, \zh{环境}).
\item Apprends les collocations : \zh{解决问题}, \zh{修路}, \zh{建设城市}.
\item Attention au piège \zh{城市} / \zh{市场} : mêmes caractères, ordre inversé.
\item Travaille les tons de \zh{市长} (4+3), \zh{居民} (1+2), \zh{环境} (2+4).
\end{itemize}
\end{aretenir}

\courssub{Dialogue modèle 1 : présenter sa commune}

\textbf{Contexte.} Un élève de Terminale (B) accueille un délégué de Suzhou (A) à la mairie de Bafoussam. L'élève présente sa commune en utilisant le lexique vu précédemment.

\begin{textebox}[Dialogue : Présentation de Bafoussam]
A: 你好！欢迎来到巴富萨。我是中文课的学生，我来介绍我们的社区。\\
\emph{Nǐ hǎo! Huānyíng lái dào Bāfùsà. Wǒ shì Zhōngwén kè de xuésheng, wǒ lái jièshào wǒmen de shèqū.}\\
« Bonjour ! Bienvenue à Bafoussam. Je suis élève de chinois, je vais présenter notre quartier. »

B: 你好！巴富萨是大城市吗？\\
\emph{Nǐ hǎo! Bāfùsà shì dà chéngshì ma?}\\
« Bonjour ! Bafoussam est-elle une grande ville ? »

A: 是的，巴富萨是西区的首府。我们市政府在市中心。我们社区有学校、医院、市场和公园。\\
\emph{Shì de, Bāfùsà shì Xīqū de shǒufǔ. Wǒmen shì zhèngfǔ zài shì zhōngxīn. Wǒmen shèqū yǒu xuéxiào, yīyuàn, shìchǎng hé gōngyuán.}\\
« Oui, Bafoussam est le chef-lieu de la Région de l'Ouest. Notre mairie est au centre-ville. Notre quartier a école, hôpital, marché et parc. »

B: 很好。市政府有什么计划？\\
\emph{Hěn hǎo. Shì zhèngfǔ yǒu shénme jìhuà?}\\
« Très bien. La mairie a quels projets ? »

A: 市政府计划修一条新路，还要建设一个新市场。居民们很高兴。\\
\emph{Shì zhèngfǔ jìhuà xiū yì tiáo xīn lù, hái yào jiànshè yí ge xīn shìchǎng. Jūmínmen hěn gāoxìng.}\\
« La mairie prévoit de construire une nouvelle route et de bâtir un nouveau marché. Les habitants sont contents. »

B: 祝你们成功！\\
\emph{Zhù nǐmen chénggōng!}\\
« Je vous souhaite du succès ! »
\end{textebox}

\begin{exemplebox}[Structures clés du dialogue]
\begin{itemize}
\item \zh{欢迎来到 + Lieu} : « Bienvenue à [Lieu] ». Ex: \zh{欢迎来到喀麦隆。}
\item \zh{是……的首府} (\emph{shì … de shǒufǔ}) : « est le chef-lieu de… ».
\item \zh{在 + Lieu} : localisation. \zh{市政府在市中心。}
\item \zh{有 + Noms (séparés par 、)} : énumération. \zh{有学校、医院和市场。}
\item \zh{计划 + Verbe / Verbe + Objet} : « prévoir de… ». \zh{计划修路} / \zh{计划建设新市场}.
\item \zh{祝 + Sujet + Verbe/Adj} : « souhaiter que… ». \zh{祝你们成功。}
\end{itemize}
\end{exemplebox}

\begin{methode}[Réussir son exposé de présentation (2-3 minutes)]
Suis ce plan en 4 étapes (les 4 familles du lexique) :
\begin{enumerate}
\item \textbf{Situer :} Nom de la ville, statut (chef-lieu ?), localisation de la mairie. \emph{Mots-clés :} \zh{城市}, \zh{首府}, \zh{市政府}, \zh{在}.
\item \textbf{Décrire les services :} Énumérer ce qui existe (école, hôpital, marché, parc, routes). \emph{Mots-clés :} \zh{社区有……}, \zh{和}, \zh{还有}.
\item \textbf{Annoncer les projets :} Ce que la mairie \emph{prévoit} de faire (construire, réparer, gérer). \emph{Mots-clés :} \zh{计划}, \zh{建设}, \zh{修}, \zh{解决}.
\item \textbf{Parler des habitants :} Leur réaction, leurs besoins, l'environnement. \emph{Mots-clés :} \zh{居民}, \zh{市民}, \zh{生活}, \zh{环境}, \zh{高兴}, \zh{希望}.
\end{enumerate}
Entraîne-toi à dire ton exposé sans lire, en te chronométrant.
\end{methode}

\courssub{Formules fonctionnelles : opiner, réagir, relancer}

Pour un débat ou un échange, tu dois exprimer ton avis, réagir à celui de l'autre et relancer la conversation. Voici les formules indispensables, classées par fonction.

\begin{center}
\begin{tabularx}{\linewidth}{|X|X|X|}
\hline
\rowcolor{popblueL} Fonction & Chinois (Pinyin) & Français \\
\hline
\multicolumn{3}{|c|}{\textbf{Donner son avis}} \\ \hline
Opinion personnelle & 我认为…… (\emph{wǒ rènwéi}) & Je pense que / Je considère que… \\
Sentiment & 我觉得…… (\emph{wǒ juéde}) & Je trouve que / J'ai l'impression que… \\
Proposition & 我建议…… (\emph{wǒ jiànyì}) & Je propose que… / Je suggère… \\
\hline
\multicolumn{3}{|c|}{\textbf{Exprimer accord / désaccord}} \\ \hline
Accord total & 我完全同意。 (\emph{wǒ wánquán tóngyì}) & Je suis tout à fait d'accord. \\
Accord partiel & 有道理，但是…… (\emph{yǒu dàoli, dànshì}) & C'est vrai, mais… \\
Désaccord poli & 我不太同意。 (\emph{wǒ bù tài tóngyì}) & Je ne suis pas très d'accord. \\
Désaccord franc & 我不同意，因为…… (\emph{wǒ bù tóngyì, yīnwèi}) & Je ne suis pas d'accord, car… \\
\hline
\multicolumn{3}{|c|}{\textbf{Relancer / Questionner}} \\ \hline
Demander l'avis & 你怎么看？ (\emph{nǐ zěnme kàn?}) & Qu'en penses-tu ? \\
Demander la raison & 为什么？ / 为什么呢？ (\emph{wèishénme? / nǐ wèishénme?}) & Pourquoi ? \\
Demander des détails & 具体来说呢？ (\emph{jùtǐ lái shuō ne?}) & Concrètement ? / Plus précisément ? \\
Changer de sujet / Ajouter & 另外，…… (\emph{lìngwài}) & Par ailleurs / En plus,… \\
\hline
\multicolumn{3}{|c|}{\textbf{Clôturer / Résumer}} \\ \hline
Résumer & 总之…… (\emph{zǒngzhī}) & En résumé / Bref… \\
Conclure & 所以，我们应该…… (\emph{suǒyǐ, wǒmen yīnggāi}) & Donc, nous devrions… \\
\end{tabularx}
\end{center}

\begin{attention}[Pièges grammaticaux des formules]
\begin{itemize}
\item \zh{认为} (\emph{rènwéi}) est plus formel que \zh{觉得} (\emph{juéde}). À l'oral, \zh{我觉得} est plus naturel.
\item \zh{建议} (\emph{jiànyì}) est suivi d'une phrase complète (sujet + verbe), pas d'un infinitif. \textbf{Correct :} \zh{我建议市政府修路。} \textbf{Incorrect :} \zh{我建议修路。} (sauf si le sujet est clair par le contexte).
\item \zh{为什么} (\emph{wèishénme}) se place \textbf{après} le sujet en chinois : \zh{你为什么不同意？} (\emph{Nǐ wèishénme bù tóngyì?}), pas au début de la phrase comme en français.
\item \zh{呢} (\emph{ne}) en fin de phrase (\zh{你怎么看呢？}) adoucit la question et invite à la réponse.
\end{itemize}
\end{attention}

\courssub{Dialogue modèle 2 : débat sur un projet municipal}

\textbf{Contexte.} Trois lycéens (A, B, C) simulent un mini-conseil municipal : faut-il dépenser le budget pour un \zh{新市场} (nouveau marché) ou pour \zh{修路} (réparer les routes) ?

\begin{textebox}[Débat : Marché ou Routes ?]
A: 大家好。市政府有钱，我们要建设新市场，还是修路？\\
\emph{Dàjiā hǎo. Shì zhèngfǔ yǒu qián, wǒmen yào jiànshè xīn shìchǎng, háishi xiū lù?}\\
« Bonjour à tous. La mairie a de l'argent, on construit un nouveau marché ou on répare les routes ? »

B: 我认为应该先修路。因为现在的路很坏，交通不方便，居民很不高兴。\\
\emph{Wǒ rènwéi yīnggāi xiān xiū lù. Yīnwèi xiànzài de lù hěn huài, jiāotōng bù fāngbiàn, jūmín hěn bù gāoxìng.}\\
« Je pense qu'il faut d'abord réparer les routes. Car les routes actuelles sont mauvaises, la circulation n'est pas pratique, les habitants ne sont pas contents. »

C: 我不同意。我觉得市场更重要。商人没有地方卖东西，经济会不好。\\
\emph{Wǒ bù tóngyì. Wǒ juéde shìchǎng gèng zhòngyào. Shāngrén méiyǒu dìfang mài dōngxi, jīngjì huì bù hǎo.}\\
« Je ne suis pas d'accord. Je trouve le marché plus important. Les commerçants n'ont pas de lieu pour vendre, l'économie ira mal. »

A: 有道理。但是，如果路不好，顾客也来不了市场。\\
\emph{Yǒu dàoli. Dànshì, rúguǒ lù bù hǎo, gùkè yě lái bu liǎo shìchǎng.}\\
« C'est vrai. Mais si les routes sont mauvaises, les clients ne pourront pas non plus venir au marché. »

B: 那我们建议：先修主路，再建设市场。\\
\emph{Nà wǒmen jiànyì: xiān xiū zhǔlù, zài jiànshè shìchǎng.}\\
« Alors nous proposons : d'abord réparer la route principale, ensuite construire le marché. »

C: 好主意！我同意这个计划。\\
\emph{Hǎo zhǔyì! Wǒ tóngyì zhège jìhuà.}\\
« Bonne idée ! J'approuve ce plan. »
\end{textebox}

\begin{exemplebox}[Analyse des arguments (Structure : Opinion + Argument + Conséquence)]
\begin{tabularx}{\linewidth}{|X|X|}
\hline
\rowcolor{popblueL} Locuteur & Structure de l'argument \\
\hline
B & \zh{我认为应该先修路} (Opinion) + \zh{因为路很坏} (Cause) + \zh{居民不高兴} (Conséquence). \\
\hline
C & \zh{我觉得市场更重要} (Opinion + Comparatif \zh{更}) + \zh{商人没有地方} (Cause) + \zh{经济会不好} (Conséquence future \zh{会}). \\
\hline
A & \zh{有道理，但是……} (Concession + Opposition) + \zh{如果……也……} (Hypothèse \zh{如果} + \zh{也}). \\
\hline
B & \zh{建议：先……再……} (Proposition séquentielle). \\
\hline
\end{tabularx}
\end{exemplebox}

\courssub{Travail phonétique : les tons du thème}

Au-delà des trois mots vus en section 1, le thème « collectivités locales » concentre des difficultés tonales précises : les tons 3\up{e} (creux), les successions de tons 4\up{e} (chutants) et les changements de ton (\cle{sandhi}).

\begin{propriete}[Rappel : Sandhi du 3\up{e} ton (三声变调)]
Quand deux syllabes de 3\up{e} ton se suivent, la \textbf{première} se prononce comme un 2\up{e} ton (montant).
\begin{itemize}
\item \zh{管理} (\emph{guǎnlǐ} $\to$ \emph{guánlǐ}) : 3+3 $\to$ 2+3.
\item \zh{建设} (\emph{jiànshè}) : 4+4 (pas de sandhi, mais double chute = fatigue vocale, bien articuler).
\item \zh{解决} (\emph{jiějué}) : 3+2 (le 3\up{e} ton reste creux car suivi d'un 2\up{e}).
\item \zh{市长} (\emph{shìzhǎng}) : 4+3 (attention à bien descendre sur \zh{shì} pour monter sur \zh{zhǎng}).
\end{itemize}
\end{propriete}

\begin{center}
\begin{tabularx}{\linewidth}{|X|X|X|X|}
\hline
\rowcolor{popblueL} Mot & Tons citation & Tons en contexte (sandhi) & Exercice de chaîne \\
\hline
管理 & 3 + 3 & \textbf{2} + 3 (\emph{guánlǐ}) & \zh{谁管理市场？} (\emph{Shéi \textbf{guánlǐ} shìchǎng?}) \\
\hline
建设 & 4 + 4 & 4 + 4 & \zh{建设城市} (\emph{Jiànshè chéngshì}) — bien marquer les deux chutes. \\
\hline
解决 & 3 + 2 & 3 + 2 (inchangé) & \zh{解决问题} (\emph{Jiějué wèntí}) — creux puis monté. \\
\hline
市长 & 4 + 3 & 4 + 3 & \zh{市长很忙} (\emph{Shìzhǎng hěn máng}) — chute puis creux. \\
\hline
环境 & 2 + 4 & 2 + 4 & \zh{环境很好} (\emph{Huánjìng hěn hǎo}) — monté puis chute. \\
\hline
一条路 & 1 + 2 + 4 & \textbf{4} + 2 + 4 (\emph{yì tiáo lù}) & \zh{修一条路} (\emph{Xiū \textbf{yì} tiáo lù}) — \zh{一} devient 4\up{e} ton devant \zh{条} (4\up{e} ton). \\
\hline
\end{tabularx}
\end{center}

\begin{experience}[Entraînement phonétique en binôme (5 min)]
\begin{enumerate}
\item \textbf{Lecture en miroir :} L'élève A lit une ligne du tableau ci-dessus, l'élève B répète en exagérant le mouvement de la voix (main qui monte/descend).
\item \textbf{Dictée tonale :} L'enseignant dit une phrase du dialogue (ex: \zh{市政府计划修一条新路}), les élèves notent seulement les numéros de tons sur papier (4-3-4-4-2-4-1-2-4). Correction collective.
\item \textbf{Conteste de fluidité :} Qui peut lire le dialogue du « Débat » (Section 4) le plus vite sans erreur de ton ? (Le speed ne doit pas sacrifier la justesse).
\end{enumerate}
\end{experience}

\courssub{Jeux de rôle et débat guidé}

\textbf{Consignes générales :} Travail en groupes de 3 ou 4. Préparez 5 minutes, jouez 3 minutes. L'enseignant circule et note : justesse des tons, utilisation du lexique ciblé (min. 5 mots), utilisation des formules fonctionnelles (min. 3), cohérence de l'argumentation.

\begin{experience}[Jeu de rôle 1 : « Le Maire et les Habitants » (Simulation de réunion publique)]
\textbf{Rôles :}
\begin{itemize}
\item \textbf{Maire (市长)} : Présente un projet (ex: \zh{建设公园} / \zh{修路} / \zh{管理市场}). Utilise \zh{计划}, \zh{建设}, \zh{为了……} (pour…).
\item \textbf{Habitant 1 (居民 A)} : Soutient le projet. Utilise \zh{我同意}, \zh{我觉得……很好}, \zh{这对……有好处}.
\item \textbf{Habitant 2 (居民 B)} : S'oppose ou demande des précisions. Utilise \zh{我不同意}, \zh{为什么……}, \zh{有没有……}, \zh{环境}, \zh{问题}.
\item \textbf{Observateur (optionnel)} : Coche sur une grille les mots-clés entendus et les formules utilisées.
\end{itemize}
\textbf{Déroulement :} Maire (1 min) $\to$ Questions/Réponses Habitants (2 min) $\to$ Maire conclut (30 s) : \zh{总之，我们会解决问题。}
\end{experience}

\begin{experience}[Jeu de rôle 2 : « Jumelage Yaoundé -- Suzhou » (Exposé + Questions)]
\textbf{Scénario :} Tu es l'élève camerounais sélectionné pour le jumelage. Tu fais un exposé de 2 minutes devant la délégation de Suzhou.
\textbf{Contraintes obligatoires :}
\begin{enumerate}
\item Saluer et te présenter (\zh{我叫……，来自喀麦隆}).
\item Situer Yaoundé (\zh{首府}, \zh{大城市}).
\item Citer 3 services existants (\zh{有医院、学校、市场}).
\item Annoncer 2 projets concrets (\zh{计划修路}, \zh{建设公园}).
\item Parler des habitants (\zh{市民}, \zh{生活}, \zh{环境}).
\item Conclure par un vœu de coopération (\zh{希望合作}, \zh{祝友谊长久}).
\end{enumerate}
\textbf{Phase 2 (Questions) :} Les « délégués » (camarades) posent 2 questions en utilisant \zh{请问……}, \zh{你们怎么……?}, \zh{有没有……?}. Tu réponds avec \zh{因为……所以……}.
\end{experience}

\begin{experience}[Débat guidé : « Priorités du budget communal »]
\textbf{Sujet :} \zh{市政府只有100万，应该修路，还是建医院？} (La mairie n'a que 1 million, faut-il réparer les routes ou construire un hôpital ?)
\textbf{Groupes :} Groupe A (Pro-Routes), Groupe B (Pro-Hôpital), Groupe C (Médiateurs/Journalistes).
\textbf{Étapes :}
\begin{enumerate}
\item \textbf{Préparation (5 min) :} Chaque groupe prépare 3 arguments structurés : \zh{第一，…… 第二，…… 第三，……}. Utiliser \zh{解决问题}, \zh{方便}, \zh{生命安全} (sécurité de la vie), \zh{经济发展}.
\item \textbf{Débat (8 min) :} Tour de table. Un argument par groupe à tour de rôle. Les médiateurs notent les arguments et forcent l'usage de \zh{你怎么看？}, \zh{我建议……}.
\item \textbf{Synthèse (2 min) :} Les médiateurs proposent un compromis en chinois : \zh{我们建议：一半修路，一半建医院。} ou \zh{先修主路，再建诊所。} (clinique).
\end{enumerate}
\end{experience}

\begin{methode}[Grille d'auto-évaluation pour l'oral (Bac blanc)]
Après chaque jeu de rôle, coche honnêtement :
\begin{center}
\begin{tabularx}{\linewidth}{|X|c|c|c|}
\hline
\rowcolor{popblueL} Critère & \textbf{Acquis} (3) & \textbf{En cours} (2) & \textbf{À travailler} (1) \\
\hline
Prononciation des tons (mots-clés) & & & \\
\hline
Utilisation du lexique thème (min 5 mots) & & & \\
\hline
Structures grammaticales correctes (ordre mots, 了/过 si besoin) & & & \\
\hline
Formules fonctionnelles (avis, accord, relance) & & & \\
\hline
Cohérence / Logique du discours & & & \\
\hline
Fluidité / Vitesse de débit & & & \\
\hline
\end{tabularx}
\end{center}
\textbf{Objectif Terminale :} Moyenne $\ge$ 2,5 / 3 sur chaque critère pour viser la mention à l'oral.
\end{methode}

\begin{savaistu}[Les « Villes nouvelles » chinoises : le cas de Xiong'an]
Depuis 2017, la Chine construit \zh{雄安新区} (\emph{Xióng'ān Xīnqū}), une « ville du futur » au sud de Pékin, pour désengorger la capitale. C'est un projet de \zh{建设} colossal : routes, hôpitaux, écoles, marchés, parcs, tout est planifié (\zh{计划}) avant construction. Au Cameroun, les villes grandissent souvent de façon spontanée. Comparer \zh{雄安} (planification centrale) et l'urbanisation de \zh{雅温得} ou \zh{杜阿拉} (croissance organique) est un excellent sujet d'exposé pour impressionner le jury : tu montres que tu connais l'actualité chinoise et que tu sais comparer (\zh{比较}) les deux modèles.
\end{savaistu}

\begin{exoresolu}[Traduction thématique (Version) : Chinois $\to$ Français]
\textbf{Énoncé.} Traduis le paragraphe suivant en français naturel :
\zh{我们社区的环境越来越好。市政府计划建设一个大公园，修两条新路。居民们很高兴，因为生活会更方便。市长说：``我们要为市民服务，解决所有的问题。''}
\tcblower
\textbf{Solution.}
« L'environnement de notre quartier devient de mieux en mieux. La mairie prévoit de construire un grand parc et de réparer deux nouvelles routes. Les habitants sont très contents, car la vie sera plus pratique. Le dit : "Nous devons servir les citoyens, résoudre tous les problèmes." »
\textbf{Points de vigilance :}
\begin{itemize}
\item \zh{越来越好} : structure comparatif de progression « de plus en plus bien ».
\item \zh{两条新路} : classificateur \zh{条} pour routes, numérateur \zh{两} (deux) et non \zh{二} devant classificateur.
\item \zh{为……服务} (\emph{wèi … fúwù}) : structure « servir… / être au service de… ».
\item \zh{所有的问题} : « tous les problèmes » (\zh{所有} + \zh{的} + nom).
\end{itemize}
\end{exoresolu}

\begin{exercice}[Expression écrite guidée : Rédige ton exposé de jumelage]
Rédige en caractères (avec pinyin si besoin) ton exposé pour la délégation de Suzhou (cf. Jeu de rôle 2). Respecte les 6 contraintes obligatoires. Nombre de caractères visé : 80--120.
\textbf{Amorce :} \zh{尊敬的各位领导，大家好。我叫……，来自喀麦隆的首都雅温得。……}
\end{exercice}

\begin{corrige}[Exercice « Expression écrite guidée » — Proposition de modèle]
\shortstack[l]{\zh{尊敬的各位领导，大家好！}\\ \zh{我叫李明，来自喀麦隆的首都雅温得。雅温得是一个大城市，市政府在市中心。}\\ \zh{我们社区有学校、医院、市场和公园，环境很好。}\\ \zh{市政府计划修两条新路，建设一个新体育馆。}\\ \zh{市民的生活越来越方便，居民们很高兴。}\\ \zh{希望中喀友谊长久，祝苏州和雅温得合作成功！谢谢大家。}}

\textbf{Pinyin :} \emph{Zūnjìng de gèwèi lǐngdǎo, dàjiā hǎo! Wǒ jiào Lǐ Míng, lái zì Kāmàilóng de shǒudū Yǎwēndé. Yǎwēndé shì yí ge dà chéngshì, shì zhèngfǔ zài shì zhōngxīn. Wǒmen shèqū yǒu xuéxiào, yīyuàn, shìchǎng hé gōngyuán, huánjìng hěn hǎo. Shì zhèngfǔ jìhuà xiū liǎng tiáo xīn lù, jiànshè yí ge xīn tǐyùguǎn. Shìmín de shēnghuó yuè lái yuè fāngbiàn, jūmínmen hěn gāoxìng. Xīwàng Zhōng-Kā yǒuyì chángjiǔ, zhù Sūzhōu hé Yǎwēndé hézuò chénggōng! Xièxie dàjiā.}

\textbf{Vérification des contraintes :} 1) Salutation/Présentation \checkmark 2) Statut Yaoundé \checkmark 3) 4 services \checkmark 4) 2 projets (\zh{修路}, \zh{建设体育馆}) \checkmark 5) Habitants (\zh{市民}, \zh{居民}, \zh{生活}) \checkmark 6) Vœu coopération \checkmark. Lexique thème utilisé : 12+ mots. Structures : \zh{计划}, \zh{越来越}, \zh{为……服务} (implicite dans le vœu).
\end{corrige}

\begin{aretenir}[Synthèse finale de la Leçon 20 — Expression orale : Collectivités locales]
\begin{itemize}
\item \textbf{Lexique :} 4 familles (Institutions, Services, Actions, Habitants) — Maîtrise des collocations (\zh{修路}, \zh{解决问题}, \zh{计划建设}).
\item \textbf{Grammaire / Structures :} \zh{计划+V}, \zh{因为……所以……}, \zh{如果……就/也……}, \zh{先……再……}, \zh{为+谁+服务}, \zh{越来越+Adj}.
\item \textbf{Fonctionnel :} \zh{我认为/觉得/建议}, \zh{同意/不同意}, \zh{有道理}, \zh{你怎么看?}, \zh{总之}.
\item \textbf{Phonétique :} Sandhi 3\up{e} ton (\zh{管理} $\to$ \emph{guánlǐ}), \zh{一} (\emph{yī} $\to$ \emph{yì} devant 4\up{e} ton), enchaînements tons 4+4 (\zh{建设}), 4+3 (\zh{市长}).
\item \textbf{Socioculturel :} Différences Maire (Chine vs Cameroun), Unité \zh{社区} vs Quartier/Arrondissement, Projet \zh{雄安新区} vs Urbanisation camerounaise.
\item \textbf{Compétence Bac :} Exposé structuré 2-3 min (Plan 4 familles) + Débat interactif 3-4 min (Formules fonctionnelles).
\end{itemize}
\end{aretenir}

\coursec{Dialogue modèle 1 : présenter sa commune}

Dans la section précédente, nous avons construit le \cle{lexique des collectivités locales} : \zh{城市} (chéngshì, ville), \zh{社区} (shèqū, quartier), \zh{市政府} (shì zhèngfǔ, mairie), \zh{市场} (shìchǎng, marché), \zh{居民} (jūmín, habitants). Ces mots ne prennent tout leur sens que lorsqu'on les met en bouche. C'est exactement l'objectif de cette section : apprendre à \cle{présenter oralement sa commune} en chinois, à travers un dialogue modèle entre un journaliste chinois et un élève camerounais. Lisez d'abord le dialogue en silence, puis à voix haute avec le pinyin, et enfin sans le pinyin.

\courssub{Le dialogue modèle : un journaliste chinois interviewe un élève de Yaoundé}

\begin{textebox}[Dialogue 1 — 请你介绍一下你的城市 (Présente ta ville, s'il te plaît)]
A（记者，jìzhě，journaliste）：你好！我是中国记者。请你介绍一下你的城市。\\
Nǐ hǎo! Wǒ shì Zhōngguó jìzhě. Qǐng nǐ jièshào yíxià nǐ de chéngshì.\\
« Bonjour ! Je suis journaliste chinois. Présente-moi ta ville, s'il te plaît. »\\[4pt]
B（学生，xuésheng，élève）：好的。我住在雅温得，它是喀麦隆的首都，有大约三百万人。\\
Hǎo de. Wǒ zhù zài Yǎwēndé, tā shì Kāmàilóng de shǒudū, yǒu dàyuē sānbǎi wàn rén.\\
« D'accord. J'habite à Yaoundé, c'est la capitale du Cameroun, il y a environ trois millions d'habitants. »\\[4pt]
A：你们的城市有什么？\\
Nǐmen de chéngshì yǒu shénme?\\
« Qu'y a-t-il dans votre ville ? »\\[4pt]
B：我们的城市有大学、医院和两个大市场。市政府现在计划修新的路。\\
Wǒmen de chéngshì yǒu dàxué, yīyuàn hé liǎng ge dà shìchǎng. Shì zhèngfǔ xiànzài jìhuà xiū xīn de lù.\\
« Notre ville a une université, des hôpitaux et deux grands marchés. La mairie prévoit maintenant de construire de nouvelles routes. »\\[4pt]
A：雅温得的天气怎么样？\\
Yǎwēndé de tiānqì zěnmeyàng?\\
« Comment est le temps à Yaoundé ? »\\[4pt]
B：不太热，也不太冷。欢迎大家来参观！\\
Bú tài rè, yě bú tài lěng. Huānyíng dàjiā lái cānguān!\\
« Il ne fait ni trop chaud ni trop froid. Bienvenue à tous pour la visite ! »\\[4pt]
A：谢谢你！\\
Xièxie nǐ!\\
« Merci ! »\\[4pt]
B：不客气！\\
Bú kèqi!\\
« De rien ! »
\end{textebox}

\textbf{Glossaire du dialogue}

\begin{center}
\begin{tabularx}{\linewidth}{|X|X|X|X|}
\hline
\rowcolor{popblueL} Caractères & Pinyin & Nature & Traduction \\
\hline
记者 & jìzhě & nom & journaliste \\
\hline
介绍 & jièshào & verbe & présenter, introduire \\
\hline
首都 & shǒudū & nom & capitale \\
\hline
大约 & dàyuē & adverbe & environ, approximativement \\
\hline
医院 & yīyuàn & nom & hôpital \\
\hline
计划 & jìhuà & verbe / nom & prévoir, planifier ; un plan \\
\hline
修 & xiū & verbe & construire, réparer \\
\hline
参观 & cānguān & verbe & visiter \\
\hline
不客气 & bú kèqi & expression & de rien, je vous en prie \\
\hline
雅温得 & Yǎwēndé & nom propre & Yaoundé \\
\hline
杜阿拉 & Dù'ālā & nom propre & Douala \\
\hline
\end{tabularx}
\end{center}

\textbf{Questions de compréhension (répondez en chinois)}

\begin{enumerate}
\item 学生住在哪个城市？(Xuésheng zhù zài nǎ ge chéngshì? — Dans quelle ville habite l'élève ?)
\item 雅温得有多少人口？(Yǎwēndé yǒu duōshao rénkǒu? — Combien d'habitants Yaoundé compte-t-elle ?)
\item 市政府计划做什么？(Shì zhèngfǔ jìhuà zuò shénme? — Que prévoit de faire la mairie ?)
\end{enumerate}

\courssub{Observer : les formules de présentation}

Relisons les phrases clés du dialogue et observons leur construction avant d'énoncer la règle.

\begin{exemplebox}[Les deux formules pivots de la présentation]
\zh{我来介绍一下我们的社区。}\\
\emph{Wǒ lái jièshào yíxià wǒmen de shèqū.}\\
« Laissez-moi vous présenter notre quartier. »\\[6pt]
\zh{我们的城市在喀麦隆的中部。}\\
\emph{Wǒmen de chéngshì zài Kāmàilóng de zhōngbù.}\\
« Notre ville se trouve dans le centre du Cameroun. »\\[6pt]
\zh{欢迎大家来参观！}\\
\emph{Huānyíng dàjiā lái cānguān!}\\
« Bienvenue à tous pour la visite ! »
\end{exemplebox}

\begin{definition}[La formule 我来介绍一下…]
\zh{我来介绍一下…} (wǒ lái jièshào yíxià…) signifie littéralement « moi, je viens présenter un peu… », c'est-à-dire « laissez-moi (vous) présenter brièvement… ». C'est la formule d'ouverture la plus naturelle pour commencer un exposé oral en chinois. Le verbe \zh{来} (lái, venir) placé devant un autre verbe exprime ici l'intention : « je vais faire… ».
\end{definition}

\begin{attention}[Le rôle de 一下 (yíxià)]
\zh{一下} (yíxià) placé \textbf{après le verbe} adoucit le ton et marque une action brève : \zh{介绍一下} = « présenter un peu / brièvement ». Comparez : \zh{介绍我们的城市} (présentez notre ville — sec, presque un ordre) et \zh{介绍一下我们的城市} (présentez-nous un peu votre ville — poli, chaleureux). À l'oral, les Chinois emploient très souvent \zh{一下} : \zh{看一下} (kàn yíxià, jeter un coup d'œil), \zh{等一下} (děng yíxià, attendre un instant), \zh{问一下} (wèn yíxià, demander un renseignement). Erreur fréquente des francophones : placer \zh{一下} avant le verbe, comme en français (« un peu présenter ») — c'est impossible en chinois, il suit toujours le verbe.
\end{attention}

\courssub{La description chiffrée simple}

Pour décrire sa commune, il faut savoir donner des chiffres. Observons les phrases du dialogue : \zh{有大约三百万人} et \zh{有两个大市场}. La structure est toujours la même : \textbf{(lieu) + 有 + (environ) + nombre + (classificateur) + nom}.

\begin{propriete}[La place de 大约 et le classificateur 个]
\zh{大约} (dàyuē, « environ ») se place \textbf{avant le chiffre}, jamais après : \zh{大约三百万人} (dàyuē sānbǎi wàn rén) = « environ trois millions de personnes ». On ne dit jamais \zh{三百万人大约}.\\[4pt]
Devant un nom, un nombre exige un \cle{classificateur}. Le classificateur général est \zh{个} (ge) : \zh{两个大市场} (liǎng ge dà shìchǎng) = « deux grands marchés ». Attention : avec le nombre 2, on dit \zh{两} (liǎng) et non \zh{二} (èr) devant un classificateur : \zh{两个市场} et non \zh{二个市场}.\\[4pt]
Pour les grands nombres : \zh{万} (wàn) = dix mille. \zh{三百万} (sānbǎi wàn) = 300 × 10 000 = 3 000 000. Le chinois compte par dizaines de mille, pas par milliers : c'est un réflexe à acquérir.
\end{propriete}

\begin{center}
\begin{tabularx}{\linewidth}{|X|X|X|}
\hline
\rowcolor{popblueL} Structure & Exemple (caractères + pinyin) & Traduction \\
\hline
有 + 大约 + nombre + 人 & 有大约三百万人 (yǒu dàyuē sānbǎi wàn rén) & il y a environ 3 millions d'habitants \\
\hline
有 + nombre + 个 + nom & 有两个大市场 (yǒu liǎng ge dà shìchǎng) & il y a deux grands marchés \\
\hline
ville + 在 + lieu & 我们的城市在喀麦隆的中部 (…zài Kāmàilóng de zhōngbù) & notre ville se trouve au centre du Cameroun \\
\hline
市政府 + 计划 + verbe & 市政府计划修新的路 (…jìhuà xiū xīn de lù) & la mairie prévoit de construire de nouvelles routes \\
\hline
\end{tabularx}
\end{center}

\courssub{Travail phonétique : les tons du thème}

Deux groupes de mots du dialogue méritent un entraînement spécial, car les francophones y perdent souvent leurs tons.

\begin{itemize}
\item \zh{两个} (liǎng ge) : 3\textsuperscript{e} ton + ton neutre. Le 3\textsuperscript{e} ton n'est prononcé à moitié (demi-3\textsuperscript{e} ton, descente sans remontée) devant le ton neutre. Répétez : liǎng ge, liǎng ge rén, liǎng ge shìchǎng.
\item \zh{介绍} (jièshào) : 4\textsuperscript{e} ton + 4\textsuperscript{e} ton. Deux tons descendants secs, comme deux petites exclamations : jiè! shào! Ne laissez pas le second ton « fondre » : chaque syllabe garde sa chute complète.
\item \zh{大约} (dàyuē) : 4\textsuperscript{e} ton + 1\textsuperscript{er} ton. Chute puis note haute et plate : dà — yuē.
\end{itemize}

\begin{methode}[Les trois étapes d'un bon exposé oral]
Un bon exposé oral en chinois suit toujours trois étapes :
\begin{enumerate}
\item \textbf{开头 (kāitóu, introduction)} : saluer et annoncer le sujet — \zh{大家好！我来介绍一下我的城市。} (Dàjiā hǎo! Wǒ lái jièshào yíxià wǒ de chéngshì. — Bonjour à tous ! Laissez-moi vous présenter ma ville.)
\item \textbf{中间 (zhōngjiān, développement)} : donner 3 ou 4 informations — situation (\zh{在…}), population (\zh{有大约…人}), lieux (\zh{有大学、医院和市场}), projets (\zh{市政府计划…}).
\item \textbf{结尾 (jiéwěi, conclusion)} : conclure et remercier — \zh{欢迎大家来！谢谢大家！} (Huānyíng dàjiā lái! Xièxie dàjiā! — Bienvenue à tous ! Merci à tous !)
\end{enumerate}
\end{methode}

\begin{popfigure}
\begin{center}
\begin{tikzpicture}[
bulle/.style={ellipse, draw=popblue, thick, fill=popblueL, align=center, minimum width=3.2cm, minimum height=1.6cm, inner sep=2pt},
fleche/.style={-{Stealth[length=3mm]}, thick, poporange}]
\node[bulle, align=center] (a) at (0,0){\zh{开头}\\introduction\\\zh{我来介绍…}};
\node[bulle, fill=poporangeL, draw=poporange, align=center] (b) at (5.2,0){\zh{中间}\\développement\\lieux, chiffres, projets};
\node[bulle, fill=popgreenL, draw=popgreen, align=center] (c) at (10.4,0){\zh{结尾}\\conclusion\\\zh{欢迎大家来！}};
\draw[fleche] (a) -- (b);
\draw[fleche] (b) -- (c);
\end{tikzpicture}
\end{center}
\legende{Les trois temps de la présentation orale d'une commune : introduction, développement, conclusion.}
\end{popfigure}

\courssub{Exercice résolu et reprise du dialogue en binôme}

\begin{exoresolu}[Mettre les phrases dans l'ordre pour reconstruire une présentation]
Énoncé — Remettez ces quatre phrases dans l'ordre logique d'un exposé oral :\\
1) 我们的城市有两个大市场和一所大学。(Wǒmen de chéngshì yǒu liǎng ge dà shìchǎng hé yì suǒ dàxué.)\\
2) 谢谢大家！(Xièxie dàjiā!)\\
3) 大家好！我来介绍一下我的城市杜阿拉。(Dàjiā hǎo! Wǒ lái jièshào yíxià wǒ de chéngshì Dù'ālā.)\\
4) 杜阿拉有大约两百万人。(Dù'ālā yǒu dàyuē liǎngbǎi wàn rén.)
\tcblower
Solution détaillée — On applique la méthode en trois étapes. L'\textbf{introduction} vient d'abord : phrase 3 (salutation + \zh{我来介绍一下…}). Le \textbf{développement} suit, du général au particulier : d'abord la population, phrase 4 (\zh{有大约两百万人}), puis les lieux, phrase 1 (\zh{有两个大市场和一所大学}). La \textbf{conclusion} termine : phrase 2 (\zh{谢谢大家！}). Ordre correct : 3 → 4 → 1 → 2. Traduction complète : « Bonjour à tous ! Laissez-moi vous présenter ma ville, Douala. Douala compte environ deux millions d'habitants. Notre ville a deux grands marchés et une université. Merci à tous ! »
\end{exoresolu}

\begin{attention}[Pièges fréquents des francophones dans ce dialogue]
\begin{itemize}
\item Ne traduisez pas « il y a » mot à mot : le français exige « il », le chinois dit simplement \zh{有}. \zh{城市有三百万人} suffit, sans sujet « il ».
\item « Deux marchés » = \zh{两个市场} (liǎng), jamais \zh{二个市场} (èr). \zh{二} sert à compter (un, deux, trois…), \zh{两} sert devant un classificateur.
\item Dans \zh{我住在雅温得}, la préposition \zh{在} (zài, « à ») se place \textbf{avant le verbe} avec le lieu : \zh{住在雅温得} = « habiter à Yaoundé ». Ne dites pas \zh{我住雅温得在}.
\item La capitale : \zh{它是喀麦隆的首都} — l'ordre est « Cameroun + de + capitale », le possesseur précède toujours le possédé avec \zh{的}.
\end{itemize}
\end{attention}

\begin{experience}[Reprise du dialogue en binôme avec substitution]
Par groupes de deux, rejouez le dialogue en remplaçant les informations soulignées : choisissez une ville (\zh{杜阿拉} Dù'ālā, Douala ; \zh{布埃亚} Bù'āiyà, Buea ; \zh{加鲁阿} Jiālǔ'ā, Garoua), changez la population (\zh{大约五十万人} dàyuē wǔshí wàn rén, environ 500 000 habitants ; \zh{大约一百万人}, environ un million) et les lieux (\zh{有三个市场} yǒu sān ge shìchǎng ; \zh{有一个大医院} yǒu yí ge dà yīyuàn). L'élève A joue le journaliste, l'élève B présente sa commune, puis échangez les rôles. Exigence : prononcez correctement \zh{两个} (liǎng ge) et \zh{介绍} (jièshào), et terminez toujours par \zh{欢迎大家来！谢谢大家！}
\end{experience}

Vous savez désormais présenter une commune en trois temps : annoncer, décrire avec des chiffres, conclure. La section suivante enrichira votre palette d'expression avec les \cle{formules fonctionnelles} pour donner votre avis, réagir et relancer la parole — des outils indispensables pour le débat qui suivra.

\coursec{Formules fonctionnelles : opiner, réagir, relancer}

Dans la section précédente, nous avons appris à présenter notre commune avec le dialogue modèle 1 : nom de la ville, quartiers, équipements, problèmes. Mais une présentation ne suffit pas : un vrai citoyen sait aussi \cle{donner son avis}, \cle{réagir} à celui des autres et \cle{relancer} la conversation. C'est exactement ce que nous allons apprendre ici : les formules fonctionnelles du débat. Ces formules sont les « outils de parole » qui transformeront vos exposés en véritables échanges.

\courssub{Observer d'abord : un mini-échange entre deux élèves}

Avant toute règle, observons ce court échange entre trois élèves de Terminale à Yaoundé, qui discutent d'un projet de leur mairie.

\begin{textebox}[Un petit débat entre camarades]
甲：我觉得市政府应该先修路。\\
Jiǎ : Wǒ juéde shì zhèngfǔ yīnggāi xiān xiū lù.\\
A : Je pense que la mairie devrait d'abord réparer les routes.\\[4pt]
乙：我不太同意，因为学校也很重要。\\
Yǐ : Wǒ bú tài tóngyì, yīnwèi xuéxiào yě hěn zhòngyào.\\
B : Je ne suis pas tout à fait d'accord, car l'école est aussi très importante.\\[4pt]
甲：你说得对。可是没有好路，学生怎么去学校呢？\\
Jiǎ : Nǐ shuō de duì. Kěshì méiyǒu hǎo lù, xuésheng zěnme qù xuéxiào ne?\\
A : Tu as raison. Mais sans bonnes routes, comment les élèves iraient-ils à l'école ?\\[4pt]
乙：嗯……你有什么看法，丙？\\
Yǐ : Ǹg…… Nǐ yǒu shénme kànfǎ, Bǐng?\\
B : Hmm…… Et toi, C, quel est ton avis ?\\[4pt]
丙：我觉得我们应该先修路，再建新学校。\\
Bǐng : Wǒ juéde wǒmen yīnggāi xiān xiū lù, zài jiàn xīn xuéxiào.\\
C : Je pense que nous devrions d'abord réparer les routes, puis construire de nouvelles écoles.
\end{textebox}

Observons : chaque prise de parole suit le même mouvement — on \textbf{donne un avis} (\zh{我觉得…}), on le \textbf{justifie} (\zh{因为…}), puis on \textbf{relance} un camarade (\zh{你有什么看法？}). C'est ce cycle que vous devez maîtriser.

\courssub{Donner son avis : 我觉得，我认为，在我看来}

\begin{definition}[看法 kànfǎ]
\zh{看法} (kànfǎ) est un nom qui signifie « point de vue, opinion ». La question \zh{你有什么看法？} (Nǐ yǒu shénme kànfǎ ?) signifie « Quel est ton point de vue ? ». C'est la question-reine du débat en classe de chinois.
\end{definition}

En chinois, il existe trois formules principales pour donner son avis, de la plus courante à la plus formelle :

\begin{center}
\begin{tabularx}{\linewidth}{|X|X|X|}
\hline
\rowcolor{popblueL} Formule & Registre & Traduction \\
\hline
\zh{我觉得…} Wǒ juéde… & courant, oral & je trouve que…, il me semble que… \\
\hline
\zh{我认为…} Wǒ rènwéi… & soutenu, exposé & je pense que…, j'estime que… \\
\hline
\zh{在我看来…} Zài wǒ kànlái… & formel, écrit ou débat & à mon avis…, selon moi… \\
\hline
\end{tabularx}
\end{center}

\begin{propriete}[Structure de l'avis]
\textbf{Structure :} Formule d'avis + phrase complète (Sujet + Verbe + Complément).\\
La formule d'avis se place \textbf{au début} de la phrase, comme en français. Après \zh{觉得}, \zh{认为} ou \zh{在我看来}, on met une phrase entière, jamais un mot isolé.\\
\zh{我觉得杜阿拉很大。} Wǒ juéde Dù'ālā hěn dà. « Je trouve que Douala est très grande. »\\
\zh{我认为中文很有用。} Wǒ rènwéi Zhōngwén hěn yǒuyòng. « Je pense que le chinois est très utile. »
\end{propriete}

\begin{exemplebox}[Avis sur les collectivités locales]
\zh{我觉得我们的城市很干净。} Wǒ juéde wǒmen de chéngshì hěn gānjìng. « Je trouve que notre ville est très propre. »\\
\zh{我认为市政府工作很努力。} Wǒ rènwéi shì zhèngfǔ gōngzuò hěn nǔlì. « Je pense que la mairie travaille avec beaucoup d'efforts. »\\
\zh{在我看来，市场太远了。} Zài wǒ kànlái, shìchǎng tài yuǎn le. « À mon avis, le marché est trop loin. »\\
\zh{我觉得这个区需要一家医院。} Wǒ juéde zhège qū xūyào yì jiā yīyuàn. « Je trouve que ce quartier a besoin d'un hôpital. »
\end{exemplebox}

Remarquez le classificateur \zh{家} (jiā) devant \zh{医院} : c'est le classificateur des établissements (hôpitaux, écoles, banques, entreprises). On dit \zh{一家医院}, \zh{一家学校}, rarement \zh{一个医院}.

Comparaison avec le français : le français dit « je trouve ça bien » avec un pronom « ça » ; le chinois exige une phrase complète après \zh{觉得} : on ne dit jamais \zh{我觉得好} tout seul dans un débat, on précise \emph{ce que} l'on trouve bien.

\courssub{Exprimer l'accord}

Pour montrer que l'on est d'accord, trois formules, de la plus simple à la plus chaleureuse :

\begin{itemize}
\item \zh{我同意。} Wǒ tóngyì. « Je suis d'accord. » — formule directe et neutre.
\item \zh{你说得对。} Nǐ shuō de duì. « Tu as raison » (littéralement : « ce que tu dis est juste »). Notez le \zh{得} qui introduit le complément de manière : \zh{说得对} = « parler juste ».
\item \zh{我也是这么想的。} Wǒ yě shì zhème xiǎng de. « Je pense pareil » (littéralement : « moi aussi je pense ainsi »). Formule très naturelle à l'oral.
\end{itemize}

\begin{exemplebox}[Réactions d'accord]
— \zh{我们应该多种树。} Wǒmen yīnggāi duō zhòng shù. « Nous devrions planter plus d'arbres. »\\
— \zh{我同意！你说得对。} Wǒ tóngyì! Nǐ shuō de duì. « Je suis d'accord ! Tu as raison. »\\
— \zh{雅温得的交通太乱了。} Yǎwēndé de jiāotōng tài luàn le. « La circulation à Yaoundé est trop désordonnée. »\\
— \zh{我也是这么想的。} Wǒ yě shì zhème xiǎng de. « Je pense exactement pareil. »
\end{exemplebox}

\courssub{Exprimer le désaccord poli}

En Chine comme au Cameroun, dire brutalement « non » à quelqu'un peut blesser. Le chinois dispose d'un outil merveilleux pour adoucir le désaccord : \zh{不太} (bú tài), « pas tout à fait ».

\begin{propriete}[Le désaccord atténué avec 不太]
\textbf{Structure :} 我 + 不太 + 同意，因为 + justification.\\
\zh{我不太同意} (Wǒ bú tài tóngyì) signifie « je ne suis pas tout à fait d'accord ». C'est la forme polie du désaccord. La forme franche \zh{我不同意} (Wǒ bù tóngyì, « je ne suis pas d'accord ») existe, mais elle est plus dure : on la réserve aux débats très engagés.\\
La conjonction \zh{可是} (kěshì, « mais ») permet aussi d'introduire une objection douce, souvent après avoir reconnu un point positif : \zh{你说得对，可是…} « Tu as raison, mais… ».
\end{propriete}

\begin{exemplebox}[Désaccords polis]
\zh{我不太同意，因为公园也很重要。} Wǒ bú tài tóngyì, yīnwèi gōngyuán yě hěn zhòngyào. « Je ne suis pas tout à fait d'accord, car les parcs sont aussi très importants. »\\
\zh{你说得对，可是我们没有那么多钱。} Nǐ shuō de duì, kěshì wǒmen méiyǒu nàme duō qián. « Tu as raison, mais nous n'avons pas autant d'argent. »\\
\zh{这个想法很好，可是太难了。} Zhège xiǎngfǎ hěn hǎo, kěshì tài nán le. « Cette idée est bonne, mais elle est trop difficile à réaliser. »
\end{exemplebox}

\begin{attention}[N'oubliez jamais 因为 !]
Les francophones oublient souvent \zh{因为} (yīnwèi) pour justifier leur avis. En français, on peut dire « je ne suis pas d'accord » et s'arrêter là. En chinois oral, un avis sans justification paraît \textbf{sec, voire impoli}. Règle d'or : après \zh{我觉得…} ou \zh{我不太同意…}, ajoutez toujours \zh{因为…} avec une raison. Mauvais réflexe : \zh{我不同意。} (sec). Bon réflexe : \zh{我不太同意，因为…} (poli et argumenté).
\end{attention}

\courssub{Relancer la parole}

Un bon débatteur ne monopolise pas la parole : il la fait circuler. Trois relances à connaître par cœur :

\begin{itemize}
\item \zh{你觉得呢？} Nǐ juéde ne? « Et toi, qu'en penses-tu ? » — la particule \zh{呢} renvoie la question à l'autre, sans répéter toute la phrase.
\item \zh{你有什么看法？} Nǐ yǒu shénme kànfǎ? « Quel est ton avis ? » — plus formelle, idéale pour passer la parole à un camarade précis.
\item \zh{为什么？} Wèishénme? « Pourquoi ? » — pour demander une justification et faire durer l'échange.
\end{itemize}

\begin{exemplebox}[Faire circuler la parole]
\zh{我觉得应该先建市场。你觉得呢？} Wǒ juéde yīnggāi xiān jiàn shìchǎng. Nǐ juéde ne? « Je pense qu'on devrait d'abord construire un marché. Et toi ? »\\
\zh{小李，你有什么看法？} Xiǎo Lǐ, nǐ yǒu shénme kànfǎ? « Petit Li, quel est ton avis ? »\\
— \zh{我不同意。} — \zh{为什么？} « — Je ne suis pas d'accord. — Pourquoi ? »
\end{exemplebox}

\courssub{Proposer avec 应该 et 可以 : la nuance qui sauve}

Pour faire des propositions, le chinois utilise deux verbes modaux qui ne disent pas la même chose :

\begin{center}
\begin{tabularx}{\linewidth}{|X|X|X|}
\hline
\rowcolor{popblueL} Modal & Force & Exemple + traduction \\
\hline
\zh{应该} yīnggāi & devrait (conseil ferme) & \zh{我们应该先修路。} « Nous devrions d'abord réparer les routes. » \\
\hline
\zh{可以} kěyǐ & pourrait (suggestion souple) & \zh{市政府可以建一个公园。} « La mairie pourrait construire un parc. » \\
\hline
\end{tabularx}
\end{center}

\begin{propriete}[Nuancer pour ne pas paraître autoritaire]
\textbf{Structure :} Sujet + 应该 / 可以 + Verbe (+ Complément).\\
\zh{应该} exprime ce qui est souhaitable, presque un devoir ; \zh{可以} exprime une simple possibilité. Dans un débat, enchaîner les \zh{应该} donne un ton de donneur de leçons. Alternez : \zh{我们应该…} pour votre priorité, \zh{市政府可以…} pour les options secondaires. C'est exactement la nuance française entre « il faut » et « on pourrait ».
\end{propriete}

\begin{propriete}[先…再… : ordonner ses propositions]
\textbf{Structure :} 先 + action 1，再 + action 2 = « d'abord…, ensuite… ».\\
Attention à la place : \zh{先} et \zh{再} se placent \textbf{avant le verbe}, jamais après. Cette structure est parfaite pour hiérarchiser un projet municipal :\\
\zh{我们应该先修路，再建公园。} Wǒmen yīnggāi xiān xiū lù, zài jiàn gōngyuán. « Nous devrions d'abord réparer les routes, puis construire un parc. »\\
\zh{先打扫市场，再建新学校。} Xiān dǎsǎo shìchǎng, zài jiàn xīn xuéxiào. « D'abord nettoyer le marché, ensuite construire une nouvelle école. »
\end{propriete}

\courssub{L'échelle d'opinion : de l'accord fort au désaccord franc}

La figure suivante résume les quatre degrés de réaction, du plus chaleureux au plus franc. En débat, choisissez votre position sur cette échelle selon votre interlocuteur et le contexte.

\begin{popfigure}
\begin{tikzpicture}[
  node distance=0.9cm,
  box/.style={rounded corners, draw=popline, thick, align=center, minimum width=6.2cm, minimum height=1cm, font=\small},
  lab/.style={font=\small\itshape, text=popmuted}
]
\node[box, fill=popgreenL, align=center] (a){\zh{我同意} Wǒ tóngyì\\ « Je suis d'accord »};
\node[box, fill=popblueL, below=of a, align=center] (b){\zh{我觉得可以} Wǒ juéde kěyǐ\\ « Je trouve que c'est possible »};
\node[box, fill=poporangeL, below=of b, align=center] (c){\zh{我不太同意} Wǒ bú tài tóngyì\\ « Je ne suis pas tout à fait d'accord »};
\node[box, fill=poppinkL, below=of c, align=center] (d){\zh{我不同意} Wǒ bù tóngyì\\ « Je ne suis pas d'accord »};
\node[lab, left=0.3cm of a] {accord fort};
\node[lab, left=0.3cm of b] {accord modéré};
\node[lab, left=0.3cm of c] {désaccord poli};
\node[lab, left=0.3cm of d] {désaccord franc};
\draw[-{Stealth}, very thick, popdark] (a.east) ++(0.8,0.3) -- ++(0,-5.2);
\end{tikzpicture}
\legende{Échelle d'opinion : du vert (accord chaleureux) au rose (désaccord franc). Plus on descend, plus il faut justifier avec 因为.}
\end{popfigure}

\courssub{Tableau récapitulatif des formules fonctionnelles}

\begin{center}
\begin{tabularx}{\linewidth}{|X|X|X|}
\hline
\rowcolor{popblueL} Fonction & Formule chinoise + pinyin & Traduction \\
\hline
Donner son avis & \zh{我觉得…} Wǒ juéde… & je trouve que… \\
\hline
Donner son avis (soutenu) & \zh{我认为…} Wǒ rènwéi… & je pense que… \\
\hline
Donner son avis (formel) & \zh{在我看来…} Zài wǒ kànlái… & à mon avis… \\
\hline
Accord & \zh{我同意。} Wǒ tóngyì. & je suis d'accord \\
\hline
Accord chaleureux & \zh{你说得对。} Nǐ shuō de duì. & tu as raison \\
\hline
Accord identique & \zh{我也是这么想的。} Wǒ yě shì zhème xiǎng de. & je pense pareil \\
\hline
Désaccord poli & \zh{我不太同意，因为…} Wǒ bú tài tóngyì, yīnwèi… & je ne suis pas tout à fait d'accord, parce que… \\
\hline
Objection douce & \zh{可是…} Kěshì… & mais… \\
\hline
Relancer & \zh{你觉得呢？} Nǐ juéde ne? & et toi, qu'en penses-tu ? \\
\hline
Relancer (formel) & \zh{你有什么看法？} Nǐ yǒu shénme kànfǎ? & quel est ton avis ? \\
\hline
Demander une raison & \zh{为什么？} Wèishénme? & pourquoi ? \\
\hline
Proposer (ferme) & \zh{我们应该…} Wǒmen yīnggāi… & nous devrions… \\
\hline
Proposer (souple) & \zh{市政府可以…} Shì zhèngfǔ kěyǐ… & la mairie pourrait… \\
\hline
Ordonner & \zh{先…，再…} xiān…, zài… & d'abord…, ensuite… \\
\hline
\end{tabularx}
\end{center}

\courssub{Méthode : le cycle « avis → justification → relance »}

\begin{methode}[Réussir une prise de parole en trois temps]
En débat, chaque intervention suit un cycle en trois étapes :
\begin{enumerate}
\item \textbf{Avis} : annoncez votre position avec \zh{我觉得…} ou \zh{我认为…}.
\item \textbf{Justification} : donnez une raison avec \zh{因为…}. Sans raison, votre avis paraît sec.
\item \textbf{Relance} : passez la parole avec \zh{你觉得呢？} ou \zh{你有什么看法？}.
\end{enumerate}
Modèle complet : \zh{我觉得我们应该先修路，因为路不好，车不能走。你觉得呢？} (Wǒ juéde wǒmen yīnggāi xiān xiū lù, yīnwèi lù bù hǎo, chē bù néng zǒu. Nǐ juéde ne?) « Je pense que nous devrions d'abord réparer les routes, car quand les routes sont mauvaises, les voitures ne peuvent pas passer. Et toi, qu'en penses-tu ? »
\end{methode}

\begin{exoresolu}[Compléter un échange]
Énoncé : complétez ce dialogue avec les formules fonctionnelles étudiées.\\
甲：我们的城市有一个问题：垃圾太多。\_\_\_\_\_\_\_\_\_\_（donner un avis : la mairie devrait d'abord nettoyer les rues）\\
乙：\_\_\_\_\_\_\_\_\_\_（désaccord poli + justification : l'eau potable est plus importante）\\
甲：\_\_\_\_\_\_\_\_\_\_（demander l'avis de 丙）\\
丙：我同意乙。我们应该先解决水的问题，再打扫街道。
\tcblower
Solution détaillée :\\
甲：\zh{我觉得市政府应该先打扫街道。} (Wǒ juéde shì zhèngfǔ yīnggāi xiān dǎsǎo jiēdào.) — formule d'avis \zh{我觉得} + modal \zh{应该} + \zh{先} placé avant le verbe pour marquer la priorité.\\
乙：\zh{我不太同意，因为干净的水更重要。} (Wǒ bú tài tóngyì, yīnwèi gānjìng de shuǐ gèng zhòngyào.) — désaccord adouci par \zh{不太}, obligatoirement suivi de \zh{因为} + justification ; notez \zh{更} (gèng, « encore plus ») pour la comparaison d'importance.\\
甲 relance 丙 par : \zh{丙，你有什么看法？} (Bǐng, nǐ yǒu shénme kànfǎ?) — relance formelle avec le nom du camarade placé en tête.\\
Remarquez que 丙 conclut avec la structure \zh{先…再…} : « d'abord résoudre le problème de l'eau, ensuite nettoyer les rues » — le cycle avis → justification → synthèse est complet.
\end{exoresolu}

\begin{attention}[Pièges fréquents des francophones]
\begin{itemize}
\item \textbf{Oublier la justification} : dire seulement \zh{我不同意} sonne brutal. Dites \zh{我不太同意，因为…}.
\item \textbf{Traduire « je suis d'accord avec toi » mot à mot} : le français « avec toi » ne se traduit pas ici ; \zh{我同意} suffit, ou \zh{我同意你的看法} (je suis d'accord avec ton point de vue).
\item \textbf{Confondre 觉得 et 认为} : \zh{觉得} est oral et spontané, \zh{认为} est réfléchi et soutenu. Dans un exposé noté, préférez \zh{我认为}.
\item \textbf{Placer 先 après le verbe} : \zh{先} se met \emph{avant} le verbe (\zh{先修路}), jamais après, contrairement au français « réparer les routes d'abord ».
\item \textbf{Tons de 看法} : kànfǎ, quatrième ton puis troisième ton — ne dites pas kánfǎ. Attention aussi à \zh{不太} : bú tài, avec le changement de ton de \zh{不} devant un quatrième ton.
\end{itemize}
\end{attention}

\begin{aretenir}[L'essentiel de la section]
Trois verbes pour opiner : \zh{我觉得} (oral), \zh{我认为} (soutenu), \zh{在我看来} (formel). Accord : \zh{我同意}, \zh{你说得对}, \zh{我也是这么想的}. Désaccord poli : \zh{我不太同意，因为…}. Relances : \zh{你觉得呢？}, \zh{你有什么看法？}, \zh{为什么？}. Propositions nuancées : \zh{应该} (ferme) contre \zh{可以} (souple), ordonnées par \zh{先…再…}. Et la règle d'or du débat : avis → justification avec \zh{因为} → relance.
\end{aretenir}

Ces formules sont maintenant vos outils. Dans la section suivante, nous les verrons à l'œuvre dans un dialogue modèle complet : un débat sur un projet municipal.

\coursec{Dialogue modèle 2 : débat sur un projet municipal}

Dans la section précédente, nous avons appris les formules fonctionnelles pour donner son avis (\zh{我觉得…}, \zh{我认为…}), réagir (\zh{我同意}, \zh{我不太同意}) et relancer la parole (\zh{你呢？}, \zh{你说呢？}). Il est temps de mettre ces outils en action dans une situation authentique et très camerounaise : un débat entre deux élèves à propos d'un projet de leur mairie. C'est exactement le type d'échange que l'on vous demandera de produire à l'oral en Terminale.

\courssub{Situation de communication}

La mairie de votre commune veut construire un \cle{nouveau marché}. Deux élèves, A et B, débattent : A est \cle{pour} le projet, B pense qu'il faut d'abord \cle{réparer les routes}. Observez d'abord le dialogue, repérez les mots qui relient les idées, puis nous analyserons chaque structure.

\begin{textebox}[Dialogue : un nouveau marché pour la commune ?]
A ：我觉得我们应该建一个新市场，因为旧市场太小了。\\
\emph{Wǒ juéde wǒmen yīnggāi jiàn yí ge xīn shìchǎng, yīnwèi jiù shìchǎng tài xiǎo le.}\\
« Je pense que nous devrions construire un nouveau marché, car l'ancien marché est trop petit. »\\[4pt]
B ：我不太同意。我认为应该先修路，如果路不好，市场就没有用。\\
\emph{Wǒ bù tài tóngyì. Wǒ rènwéi yīnggāi xiān xiū lù, rúguǒ lù bù hǎo, shìchǎng jiù méiyǒu yòng.}\\
« Je ne suis pas tout à fait d'accord. Je pense qu'il faut d'abord réparer les routes ; si la route est mauvaise, le marché ne sert à rien. »\\[4pt]
A ：可是市场也很重要，而且很多人需要工作。\\
\emph{Kěshì shìchǎng yě hěn zhòngyào, érqiě hěn duō rén xūyào gōngzuò.}\\
« Mais le marché est aussi très important, et de plus beaucoup de gens ont besoin de travail. »\\[4pt]
B ：你说得对，所以我们应该先修路，再建市场。\\
\emph{Nǐ shuō de duì, suǒyǐ wǒmen yīnggāi xiān xiū lù, zài jiàn shìchǎng.}\\
« Tu as raison, donc nous devrions d'abord réparer les routes, puis construire le marché. »
\end{textebox}

\textbf{Glossaire du dialogue}

\begin{center}
\begin{tabularx}{\linewidth}{|X|X|X|X|}
\hline
\rowcolor{popblueL} Caractères & Pinyin & Nature & Traduction \\
\hline
建 & jiàn & verbe & construire \\
\hline
市场 & shìchǎng & nom & marché \\
\hline
因为 & yīnwèi & conjonction & parce que, car \\
\hline
旧 & jiù & adjectif & vieux, ancien (objets) \\
\hline
同意 & tóngyì & verbe & être d'accord \\
\hline
先 & xiān & adverbe & d'abord \\
\hline
修路 & xiū lù & VO (verbe-objet) & réparer les routes \\
\hline
如果…就… & rúguǒ… jiù… & conjonction & si… alors… \\
\hline
没有用 & méiyǒu yòng & expression & ne servir à rien, être inutile \\
\hline
可是 & kěshì & conjonction & mais \\
\hline
而且 & érqiě & conjonction & de plus, et aussi \\
\hline
需要 & xūyào & verbe & avoir besoin de \\
\hline
所以 & suǒyǐ & conjonction & donc, c'est pourquoi \\
\hline
再 & zài & adverbe & ensuite, puis (action future) \\
\hline
\end{tabularx}
\end{center}

\courssub{Observation : les connecteurs logiques du débat}

Relevez dans le dialogue les mots qui organisent les arguments. Ce sont les \cle{connecteurs logiques} : sans eux, un débat n'est qu'une suite de phrases isolées. Voici le tableau récapitulatif des cinq connecteurs de cette leçon, avec la comparaison chinois / français.

\begin{center}
\begin{tabularx}{\linewidth}{|X|X|X|}
\hline
\rowcolor{popblueL} Structure & Exemple (caractères + pinyin) & Traduction \\
\hline
因为 + cause & 因为旧市场太小了。\newline \emph{Yīnwèi jiù shìchǎng tài xiǎo le.} & Parce que l'ancien marché est trop petit. \\
\hline
可是 (tête de phrase 2) & 可是市场也很重要。\newline \emph{Kěshì shìchǎng yě hěn zhòngyào.} & Mais le marché est aussi très important. \\
\hline
而且 + ajout & 而且很多人需要工作。\newline \emph{Érqiě hěn duō rén xūyào gōngzuò.} & De plus, beaucoup de gens ont besoin de travail. \\
\hline
所以 + conséquence & 所以我们应该先修路。\newline \emph{Suǒyǐ wǒmen yīnggāi xiān xiū lù.} & Donc nous devrions d'abord réparer les routes. \\
\hline
如果…就… & 如果路不好，市场就没有用。\newline \emph{Rúguǒ lù bù hǎo, shìchǎng jiù méiyǒu yòng.} & Si la route est mauvaise, le marché ne sert à rien. \\
\hline
\end{tabularx}
\end{center}

\begin{propriete}[Les connecteurs du débat : forme et emploi]
\begin{itemize}
\item \zh{因为} (yīnwèi, « parce que ») introduit la \cle{cause}. Deux positions possibles :
      \begin{itemize}
      \item \textbf{Position standard (antéposée)} : \zh{因为旧市场太小了，所以我们应该建新市场。} (Comme l'ancien marché est trop petit, nous devrions en construire un nouveau.)
      \item \textbf{Position postposée (orale, très fréquente)} : \zh{我们应该建新市场，因为旧市场太小了。} (Nous devrions construire un nouveau marché, car l'ancien est trop petit.) $\leftarrow$ C'est la structure du dialogue.
      \end{itemize}
\item \zh{可是} (kěshì, « mais ») introduit une \cle{opposition}. Il se place en tête de la deuxième phrase. Synonyme soutenu : \zh{但是} (dànshì).
\item \zh{而且} (érqiě, « de plus ») ajoute un \cle{argument supplémentaire}. Il se place devant le verbe ou en tête de phrase : \zh{而且很多人需要工作。}
\item \zh{所以} (suǒyǐ, « donc ») introduit la \cle{conséquence}. \textbf{Particularité majeure} : en chinois, \zh{因为…所以…} peuvent coexister dans la même phrase (« parce que… donc… »), ce qui est fautif en français (*« Parce que la route est mauvaise, donc le marché ne sert à rien »).
      \begin{center}
      \zh{因为路不好，所以市场没有用。} \emph{(Correct en chinois, incorrect en français)}
      \end{center}
\item \zh{如果…就…} (rúguǒ… jiù…, « si… alors… ») exprime la \cle{condition}. Structure : \textbf{如果 + condition，sujet + 就 + résultat}. Le mot \zh{就} (jiù) se place \textbf{juste avant le verbe} du résultat et ne se traduit pas seul ; il marque la corrélation conditionnelle.
\end{itemize}
\end{propriete}

\begin{exemplebox}[Exemples traduits]
\zh{如果路不好，市场就没有用。}\\
\emph{Rúguǒ lù bù hǎo, shìchǎng jiù méiyǒu yòng.}\\
« Si la route est mauvaise, le marché ne sert à rien. »\\[4pt]
\zh{而且很多人需要工作。}\\
\emph{Érqiě hěn duō rén xūyào gōngzuò.}\\
« De plus, beaucoup de gens ont besoin de travail. »\\[4pt]
\zh{如果明天下雨，我们就不去市场。}\\
\emph{Rúguǒ míngtiān xià yǔ, wǒmen jiù bù qù shìchǎng.}\\
« S'il pleut demain, nous n'irons pas au marché. »\\[4pt]
\zh{因为很多人都同意，所以市长决定建一个新市场。}\\
\emph{Yīnwèi hěn duō rén dōu tóngyì, suǒyǐ shìzhǎng juédìng jiàn yí ge xīn shìchǎng.}\\
« Comme beaucoup de gens sont d'accord, le maire a décidé de construire un nouveau marché. »
\end{exemplebox}

\courssub{Travail phonétique : la règle du sandhi du 3e ton}

Prononcez à voix haute : \zh{所以} (suǒyǐ), \zh{市长} (shìzhǎng), \zh{很好} (hěn hǎo), \zh{可以} (kěyǐ). Vous remarquez quelque chose ? Personne, pas même les Chinois natifs, ne prononce deux 3e tons complets à la suite : ce serait trop lourd et trop lent. Le premier 3e ton change automatiquement.

\begin{propriete}[Règle du sandhi du 3e ton (变调 biàndiào)]
Quand \textbf{deux 3e tons se suivent}, le \textbf{premier se prononce comme un 2e ton} (montant) ; le deuxième garde son 3e ton (bas-descendant-remontant).
\begin{center}
3e ton + 3e ton \quad $\rightarrow$ \quad 2e ton + 3e ton \quad (à l'oral uniquement)
\end{center}
Ainsi \zh{所以} s'écrit \emph{suǒyǐ} mais se prononce \emph{suóyǐ}. \textbf{Important :} le pinyin écrit garde toujours le ton d'origine (citation form) ; on n'écrit jamais « suóyǐ », le changement n'existe qu'à l'oral.
\end{propriete}

\begin{popfigure}
\begin{tikzpicture}[
  bloc/.style={rectangle, rounded corners=4pt, draw=popblue, fill=popblueL, minimum width=2.2cm, minimum height=1.1cm, align=center, font=\small},
  res/.style={rectangle, rounded corners=4pt, draw=popgreen, fill=popgreenL, minimum width=2.2cm, minimum height=1.1cm, align=center, font=\small},
  fleche/.style={-{Stealth[length=3mm]}, very thick, poporange}
]
\node[bloc, align=center] (a) at (0,0){suǒ\\ 3e ton};
\node[bloc, align=center] (b) at (2.8,0){yǐ\\ 3e ton};
\node[res, align=center] (c) at (7.2,0){suó\\ 2e ton};
\node[res, align=center] (d) at (10,0){yǐ\\ 3e ton};
\draw[fleche] (3.9,0) -- (5.9,0) node[midway, above, font=\small\itshape, text=popdark] {à l'oral};
\node[font=\small, text=popmuted] at (1.4,-1.1) {écrit : suǒyǐ};
\node[font=\small, text=popmuted] at (8.6,-1.1) {prononcé : suóyǐ};
\end{tikzpicture}
\legende{Le sandhi du 3e ton : 所以 s'écrit suǒyǐ mais se prononce « suóyǐ ».}
\end{popfigure}

La même règle s'applique à tous les mots de notre thème : \zh{市长} (shìzhǎng $\rightarrow$ \emph{shízhǎng}, « le maire »), \zh{很好} (hěn hǎo $\rightarrow$ \emph{hén hǎo}, « très bien »), \zh{可以} (kěyǐ $\rightarrow$ \emph{kéyǐ}, « pouvoir »). En revanche, \zh{可是} (kěshì) combine un 3e ton et un 4e ton : il n'y a \textbf{pas} de sandhi « 2e ton », mais le 3e ton devant un ton non-3e se prononce souvent « à moitié » (ton bas, sans remontée finale, dit \emph{half third tone}). Dites simplement un \emph{kě} bas et court, puis un \emph{shì} franc et descendant.

\begin{attention}[Erreurs fréquentes des francophones]
\begin{itemize}
\item \textbf{Sandhi ignoré} : Beaucoup d'élèves prononcent \zh{所以} avec deux 3e tons complets (descente-remontée, puis descente-remontée). Résultat : la parole est lente, lourde, et cela coûte des points à l'évaluation orale. Entraînez-vous à dire \emph{suóyǐ} d'un seul élan, comme un seul mot.
\item \textbf{Confusion auditive} : Ne confondez pas à l'oreille \zh{所以} (suǒyǐ, « donc ») et \zh{可以} (kěyǐ, « pouvoir ») : les deux subissent le sandhi (premier ton $\rightarrow$ 2e ton), mais la consonne initiale diffère (\emph{s-} / \emph{k-}).
\item \textbf{Position de 就} : Dans \zh{如果…就…}, \zh{just} précède le verbe (\zh{就没有用}, \zh{就不去}). Ne l'oubliez pas : *\zh{如果路不好，市场没有用就} est incorrect.
\end{itemize}
\end{attention}

\courssub{Méthode : apprendre et reproduire le débat}

\begin{methode}[Mémoriser le débat par blocs, pas mot à mot]
Pour gagner en fluidité à l'oral, ne récitez jamais le dialogue phrase par phrase comme un poème. Procédez par \cle{blocs d'idées} (chunks) :
\begin{enumerate}
\item \textbf{Avis} : chaque élève annonce sa position (\zh{我觉得我们应该…} / \zh{我不太同意。我认为…}).
\item \textbf{Justification} : chacun donne sa raison avec \zh{因为} (cause) ou \zh{如果…就…} (condition).
\item \textbf{Réaction} : on répond à l'autre avec \zh{可是} (opposition) ou \zh{你说得对} (concession partielle).
\item \textbf{Conclusion} : on propose une solution commune avec \zh{所以…先…再…} (conséquence + séquencement).
\end{enumerate}
Retenez seulement le « squelette » : \textbf{Avis $\rightarrow$ Justification $\rightarrow$ Réaction $\rightarrow$ Conclusion}. Les mots exacts peuvent varier : c'est ce qui rendra votre débat naturel et spontané.
\end{methode}

\begin{exoresolu}[Compléter le débat avec les connecteurs]
Énoncé — Complétez les phrases avec le connecteur qui convient : \zh{因为}，\zh{可是}，\zh{而且}，\zh{所以}，\zh{如果…就…}\\
1. 我们应该先修路，\_\_\_\_\_\_ 路太破了。\\
2. 新市场很漂亮，\_\_\_\_\_\_ 路不好，\_\_\_\_\_\_ 没有人来。\\
3. 市场很重要，\_\_\_\_\_\_ 它能给很多人工作。\\
4. 你说得对，\_\_\_\_\_\_ 我们先修路，再建市场。
\tcblower
Solution détaillée :
\begin{enumerate}
\item \zh{因为} : on donne la cause (« parce que la route est trop abîmée »). \zh{我们应该先修路，因为路太破了。} (Position postposée orale).
\item \zh{如果…就…} : c'est une condition et son résultat. \zh{新市场很漂亮，如果路不好，就没有人来。} Remarquez que \zh{就} se place avant le verbe \zh{没有} (ou \zh{来}).
\item \zh{而且} : on ajoute un argument (« de plus, il peut donner du travail à beaucoup de gens »).
\item \zh{所以} : on conclut (« donc, réparons d'abord la route, puis construisons le marché »).
\end{enumerate}
\textbf{Vérification à l'oral} : lisez la phrase 4 à voix haute en appliquant le sandhi — \zh{所以} se prononce \emph{suóyǐ}.
\end{exoresolu}

\courssub{De l'écoute au jeu de rôle}

\begin{experience}[Écoute, chœur, puis débat en binôme]
\begin{enumerate}
\item \textbf{Écoute active} : le professeur lit le dialogue une première fois, livre fermé. Les élèves lèvent la main à chaque connecteur entendu (\zh{因为}, \zh{可是}, \zh{而且}, \zh{所以}, \zh{如果}).
\item \textbf{Répétition en chœur} : la classe répète phrase par phrase, en insistant sur le sandhi (\emph{suóyǐ}, \emph{shízhǎng}, \emph{kéyǐ}) et sur le 3e ton bas (\emph{half third tone}) de \zh{可是} (\emph{kěshì}).
\item \textbf{Jeu de rôle en binôme (guidé)} : chaque paire rejoue le dialogue original, puis inverse les rôles (A devient B).
\item \textbf{Variante libre (transfert)} : remplacez le projet « nouveau marché » par un autre projet local (une école, un terrain de football, un hôpital, un forage) en gardant le squelette \textbf{Avis $\rightarrow$ Justification $\rightarrow$ Réaction $\rightarrow$ Conclusion}. Utilisez le vocabulaire de la leçon 19 (collectivités locales).
\item \textbf{Défi oral} : un binôme volontaire présente son débat devant la classe sans regarder sa feuille ; la classe note sur une fiche d'évaluation les connecteurs utilisés et la justesse des tons (sandhi).
\end{enumerate}
\end{experience}

\begin{aretenir}[L'essentiel de la section]
\begin{itemize}
\item Un débat chinois s'organise en quatre blocs : \cle{avis} $\rightarrow$ \cle{justification} (\zh{因为} / \zh{如果…就…}) $\rightarrow$ \cle{réaction} (\zh{可是}, \zh{你说得对}) $\rightarrow$ \cle{conclusion} (\zh{所以…先…再…}).
\item \zh{因为…所以…} peuvent coexister dans la même phrase chinoise (structure parataxique), contrairement au français (hypotaxe).
\item \textbf{Sandhi du 3e ton} : 3e + 3e $\rightarrow$ 2e + 3e à l'oral (\zh{所以} = \emph{suóyǐ}, \zh{市长} = \emph{shízhǎng}), mais l'écriture pinyin ne change jamais.
\item \zh{就} dans \zh{如果…就…} se place \textbf{immédiatement avant le verbe} de la conséquence.
\item Mémorisez par blocs d'idées (chunks), jamais mot à mot : c'est la clé de la fluidité et de la naturel à l'oral.
\end{itemize}
\end{aretenir}

\coursec{Travail phonétique : les tons du thème}

Après le débat sur le projet municipal, nous avons entendu beaucoup de mots nouveaux : \zh{市长} « le maire », \zh{修路} « réparer une route », \zh{公园} « parc », \zh{居民} « habitants ». Mais un mot mal prononcé peut changer de sens : dire \zh{路} \emph{lù} (4\textsuperscript{e} ton) au lieu de \zh{绿} \emph{lǜ} (3\textsuperscript{e} ton), c'est confondre « route » et « vert ». Cette section est donc consacrée à l'entraînement phonétique : nous allons réviser les tons, travailler les paires minimales du thème, maîtriser les suites de 3\textsuperscript{e} tons, puis lire le dialogue 2 avec une prononciation modèle.

\courssub{Révision rapide : les quatre tons et le ton neutre}

\begin{definition}[Le ton en chinois]
Le chinois est une langue \cle{tonale} : chaque syllabe est prononcée avec une mélodie (un \cle{ton}) qui fait partie du mot, au même titre que ses consonnes et ses voyelles. Le mandarin standard possède \cle{quatre tons} et un \cle{ton neutre} (léger et court, sans marque). Contrairement au français, où la mélodie exprime seulement l'émotion ou la question, en chinois elle \emph{change le sens du mot}.
\end{definition}

Observons d'abord les quatre tons sur des mots de notre thème « collectivités locales » :

\begin{center}
\begin{tabularx}{\linewidth}{|X|X|X|X|}\hline
\rowcolor{popblueL} Ton & Mouvement mélodique & Mot du thème & Traduction \\ \hline
1\textsuperscript{er} ton (—) & haut et plat & \zh{修} \emph{xiū} & réparer \\ \hline
2\textsuperscript{e} ton (/) & montant & \zh{园} \emph{yuán} & jardin, parc \\ \hline
3\textsuperscript{e} ton (V) & descendant puis montant & \zh{管} \emph{guǎn} & gérer, administrer \\ \hline
4\textsuperscript{e} ton (\textbackslash) & descendant brusque & \zh{市} \emph{shì} & ville, municipal \\ \hline
Ton neutre & court et léger & \zh{们} \emph{men} (dans \zh{我们}) & (marque du pluriel) \\ \hline
\end{tabularx}
\end{center}

\begin{popfigure}
\begin{tikzpicture}[scale=1.0, every node/.style={font=\small}]
% Ton 1
\begin{scope}[shift={(0,0)}]
\draw[popline, dashed] (0,0) -- (2.6,0);
\draw[popline, dashed] (0,1.6) -- (2.6,1.6);
\draw[line width=1.6pt, popblue] (0.2,1.5) -- (2.4,1.5);
\node[align=center] at (1.3,-0.7) {1\textsuperscript{er} ton\\ \zh{修} \emph{xiū} « réparer »};
\end{scope}
% Ton 2
\begin{scope}[shift={(3.6,0)}]
\draw[popline, dashed] (0,0) -- (2.6,0);
\draw[popline, dashed] (0,1.6) -- (2.6,1.6);
\draw[line width=1.6pt, poporange] (0.2,0.3) -- (2.4,1.5);
\node[align=center] at (1.3,-0.7) {2\textsuperscript{e} ton\\ \zh{园} \emph{yuán} « parc »};
\end{scope}
% Ton 3
\begin{scope}[shift={(7.2,0)}]
\draw[popline, dashed] (0,0) -- (2.6,0);
\draw[popline, dashed] (0,1.6) -- (2.6,1.6);
\draw[line width=1.6pt, popgreen] (0.2,1.1) .. controls (0.9,0.1) and (1.5,0.1) .. (2.4,1.4);
\node[align=center] at (1.3,-0.7) {3\textsuperscript{e} ton\\ \zh{管} \emph{guǎn} « gérer »};
\end{scope}
% Ton 4
\begin{scope}[shift={(10.8,0)}]
\draw[popline, dashed] (0,0) -- (2.6,0);
\draw[popline, dashed] (0,1.6) -- (2.6,1.6);
\draw[line width=1.6pt, poppink] (0.2,1.5) -- (2.4,0.2);
\node[align=center] at (1.3,-0.7) {4\textsuperscript{e} ton\\ \zh{市} \emph{shì} « ville »};
\end{scope}
\end{tikzpicture}
\legende{Frise des quatre tons du mandarin avec un mot du thème sous chaque courbe.}
\end{popfigure}

\begin{methode}[La gestuelle des tons : la mémoire du corps]
Pour ancrer les tons dans la mémoire, utilisez votre \cle{main} pendant que vous prononcez :
\begin{enumerate}
\item \textbf{1\textsuperscript{er} ton} : main à plat, haute, qui avance horizontalement — la voix reste haute et stable, comme une note tenue.
\item \textbf{2\textsuperscript{e} ton} : main qui monte en diagonale — comme quand on demande « hein ? » en français.
\item \textbf{3\textsuperscript{e} ton} : main qui descend puis remonte (un creux) — c'est le ton le plus long ; prenez le temps de descendre \emph{avant} de remonter.
\item \textbf{4\textsuperscript{e} ton} : main qui tombe brusquement, comme un karatéka qui coupe — court, sec, décidé, comme un « non ! » ferme.
\item \textbf{Ton neutre} : petit coup sec et léger, sans mélodie.
\end{enumerate}
Dessiner le ton dans l'air en parlant aide la \cle{mémoire corporelle} : c'est particulièrement efficace pour le 3\textsuperscript{e} ton, que les francophones ont tendance à raccourcir.
\end{methode}

\courssub{Paires minimales à risque du thème}

Une \cle{paire minimale} est un couple de mots qui ne diffèrent que par un seul trait — ici, le ton ou une voyelle proche. Les confondre peut changer complètement le sens de votre phrase à l'oral.

\begin{exemplebox}[Paires minimales du thème « collectivités »]
\begin{itemize}
\item \zh{市} \emph{shì} « ville » et \zh{是} \emph{shì} « être » : \textbf{homophones parfaits} (même ton !). Seul le contexte permet de distinguer : \zh{这是市长。} \emph{Zhè shì shìzhǎng.} « Voici le maire. » — les deux \emph{shì} se prononcent pareil, mais ne s'écrivent pas pareil.
\item \zh{修} \emph{xiū} (1\textsuperscript{er} ton) « réparer » vs \zh{休} \emph{xiū} (1\textsuperscript{er} ton) « se reposer » : homophones également ; attention à l'écriture : \zh{休息} \emph{xiūxi} « se reposer », \zh{修路} \emph{xiū lù} « réparer la route ».
\item \zh{路} \emph{lù} (4\textsuperscript{e} ton) « route » vs \zh{绿} \emph{lǜ} (3\textsuperscript{e} ton) « vert » : ici, c'est bien le \textbf{ton} qui change le sens. \zh{绿色的路} \emph{lǜsè de lù} « la route verte » combine les deux !
\end{itemize}
\end{exemplebox}

\begin{attention}[Piège : la voyelle ü après l, n, j, q, x]
\zh{绿} \emph{lǜ} contient la voyelle \cle{ü} (lèvres arrondies comme pour « u » français, langue comme pour « i »). Après \emph{j, q, x}, on écrit \emph{u} mais on prononce toujours \emph{ü} : \zh{去} \emph{qù} se prononce « tchü », et \zh{居民} \emph{jūmín} « habitants » se prononce « jümin ». Après \emph{l} et \emph{n}, on écrit les deux points : \zh{绿} \emph{lǜ}, \zh{女} \emph{nǚ}. Ne prononcez jamais \emph{lǜ} comme « lou » : votre auditeur chinois entendrait un autre mot.
\end{attention}

\courssub{Les suites de 3\textsuperscript{e} tons : la règle de changement de ton}

Lisez à voix haute : \zh{市长} \emph{shìzhǎng}, \zh{很好} \emph{hěn hǎo}, \zh{可以} \emph{kěyǐ}, \zh{所以} \emph{suǒyǐ}. Vous remarquez qu'il est difficile d'enchaîner deux 3\textsuperscript{e} tons complets (descente-remontée, puis descente-remontée) : la langue fatigue. Le chinois a donc une règle automatique.

\begin{propriete}[Changement de ton du 3\textsuperscript{e} ton]
\textbf{Règle 1 — Deux 3\textsuperscript{e} tons qui se suivent} : le premier 3\textsuperscript{e} ton se prononce comme un \cle{2\textsuperscript{e} ton} (montant). On écrit toujours le ton d'origine dans le pinyin, mais on prononce le changement : \zh{很好} s'écrit \emph{hěn hǎo} mais se prononce \emph{hén hǎo}.

\textbf{Règle 2 — 3\textsuperscript{e} ton devant un autre ton} (1\textsuperscript{er}, 2\textsuperscript{e}, 4\textsuperscript{e} ou neutre) : le 3\textsuperscript{e} ton devient un \cle{demi-3\textsuperscript{e} ton} : on ne fait que la descente, sans la remontée. C'est le cas le plus fréquent dans la langue parlée.
\end{propriete}

\begin{exemplebox}[La règle en action dans notre thème]
\begin{itemize}
\item \zh{很好} \emph{hěn hǎo} → prononcé \emph{hén hǎo} « très bien » (règle 1).
\item \zh{可以} \emph{kěyǐ} → prononcé \emph{kéyǐ} « pouvoir » (règle 1).
\item \zh{所以} \emph{suǒyǐ} → prononcé \emph{suóyǐ} « donc, c'est pourquoi » (règle 1).
\item \zh{市长} \emph{shìzhǎng} « maire » : le 3\textsuperscript{e} ton de \zh{长} est en fin de mot, il reste complet ; mais dans \zh{北京} \emph{Běijīng}, le \zh{北} \emph{běi} devant le 1\textsuperscript{er} ton \emph{jīng} devient un demi-3\textsuperscript{e} ton : descente seule (règle 2).
\end{itemize}
\end{exemplebox}

\begin{attention}[Le demi-3\textsuperscript{e} ton : l'erreur n° 1 des francophones]
Devant un 1\textsuperscript{er}, 2\textsuperscript{e}, 4\textsuperscript{e} ton ou un ton neutre, le 3\textsuperscript{e} ton se prononce presque toujours « à moitié » : \textbf{descente seule, sans remontée}. Exemples : \zh{北京} \emph{Běijīng} (le \emph{běi} descend seulement), \zh{老师} \emph{lǎoshī} (le \emph{lǎo} descend seulement), \zh{我们} \emph{wǒmen} (le \emph{wǒ} descend seulement devant le ton neutre). Si vous remontez à chaque fois, votre phrase devient lente et artificielle. En revanche, n'oubliez pas la remontée quand le 3\textsuperscript{e} ton est \textbf{seul ou en fin de phrase} : \zh{好！} \emph{Hǎo !} « D'accord ! ».
\end{attention}

\begin{center}
\begin{tabularx}{\linewidth}{|X|X|X|}\hline
\rowcolor{popblueL} Structure tonale & Exemple (écriture → prononciation) & Traduction \\ \hline
3\textsuperscript{e} + 3\textsuperscript{e} & \zh{很好} \emph{hěn hǎo} → \emph{hén hǎo} & très bien \\ \hline
3\textsuperscript{e} + 3\textsuperscript{e} & \zh{可以} \emph{kěyǐ} → \emph{kéyǐ} & pouvoir \\ \hline
3\textsuperscript{e} + 3\textsuperscript{e} & \zh{所以} \emph{suǒyǐ} → \emph{suóyǐ} & donc \\ \hline
3\textsuperscript{e} + 1\textsuperscript{er} & \zh{北京} \emph{Běijīng} → demi-3\textsuperscript{e} + 1\textsuperscript{er} & Pékin \\ \hline
4\textsuperscript{e} + 3\textsuperscript{e} & \zh{市长} \emph{shìzhǎng} (pas de changement) & maire \\ \hline
1\textsuperscript{er} + 4\textsuperscript{e} & \zh{修路} \emph{xiū lù} (pas de changement) & réparer la route \\ \hline
\end{tabularx}
\end{center}

\courssub{Les tons en contexte : lire le dialogue 2}

Reprenons un extrait du dialogue modèle 2 (débat sur le projet municipal). La méthode de lecture se fait en \textbf{deux étapes} : d'abord en \cle{exagérant} les tons (lentement, avec la gestuelle), puis à \cle{vitesse normale}, en appliquant les changements de ton.

\begin{textebox}[Extrait du dialogue 2 — le débat municipal]
甲：我觉得我们市应该修一条新路。\\
\emph{Jiǎ : Wǒ juéde wǒmen shì yīnggāi xiū yì tiáo xīn lù.}\\
A : Je pense que notre ville devrait construire une nouvelle route.\\[4pt]
乙：我不同意。所以我认为应该先修学校。\\
\emph{Yǐ : Wǒ bù tóngyì. Suǒyǐ wǒ rènwéi yīnggāi xiān xiū xuéxiào.}\\
B : Je ne suis pas d'accord. C'est pourquoi je pense qu'on devrait d'abord construire des écoles.\\[4pt]
甲：你说得很好，可是居民也需要公园。\\
\emph{Jiǎ : Nǐ shuō de hěn hǎo, kěshì jūmín yě xūyào gōngyuán.}\\
A : Tu parles très bien, mais les habitants ont aussi besoin de parcs.
\end{textebox}

Points de vigilance tonale dans cet extrait :
\begin{itemize}
\item \zh{觉得} \emph{juéde} : 2\textsuperscript{e} ton + ton neutre — le \emph{de} est très léger.
\item \zh{所以} \emph{suǒyǐ} → prononcé \emph{suóyǐ} (règle des deux 3\textsuperscript{e} tons).
\item \zh{很好} \emph{hěn hǎo} → prononcé \emph{hén hǎo}.
\item \zh{一条} \emph{yì tiáo} : rappel de la règle de changement de \zh{一} — seul ou dans un nombre, \zh{一} est au 1\textsuperscript{er} ton (\emph{yī}) ; devant un 4\textsuperscript{e} ton, il devient 2\textsuperscript{e} ton (\emph{yí}) ; devant un 1\textsuperscript{er}, 2\textsuperscript{e} ou 3\textsuperscript{e} ton, il devient 4\textsuperscript{e} ton (\emph{yì}). Ici \zh{条} \emph{tiáo} est au 2\textsuperscript{e} ton, donc on prononce \emph{yì tiáo}.
\item \zh{不同意} \emph{bù tóngyì} : \zh{不} reste au 4\textsuperscript{e} ton devant un 2\textsuperscript{e} ton ; il ne devient 2\textsuperscript{e} ton (\emph{bú}) que devant un 4\textsuperscript{e} ton : \zh{不是} \emph{bú shì}.
\end{itemize}

\begin{exoresolu}[Repérer et corriger les tons]
Énoncé : un élève de Tle A lit la phrase \zh{我们市长想修一条绿色的路。} en prononçant : \emph{wǒmen shìzhāng xiǎng xiū yì tiáo lǜsè de lù} avec \emph{zhāng} au 1\textsuperscript{er} ton et \emph{lǜ} prononcé « lou ». Relevez les deux erreurs et donnez la lecture correcte.
\tcblower
Solution détaillée :
\begin{enumerate}
\item \textbf{Erreur de ton} : \zh{长} dans \zh{市长} se prononce \emph{zhǎng} (3\textsuperscript{e} ton), pas \emph{zhāng}. En fin de mot, le 3\textsuperscript{e} ton est complet : descente \emph{puis} remontée. Lecture correcte : \emph{shìzhǎng}.
\item \textbf{Erreur de voyelle} : \zh{绿} \emph{lǜ} contient la voyelle \emph{ü} (lèvres arrondies), pas « ou ». Lecture correcte : \emph{lǜsè}.
\end{enumerate}
Phrase modèle complète : \zh{我们市长想修一条绿色的路。} \emph{Wǒmen shìzhǎng xiǎng xiū yì tiáo lǜsè de lù.} « Notre maire veut construire une route verte. » Remarquez aussi : \zh{条} \emph{tiáo} est au 2\textsuperscript{e} ton, donc \zh{一} se prononce au 4\textsuperscript{e} ton : \emph{yì tiáo}. Et \zh{想} \emph{xiǎng}, devant le 1\textsuperscript{er} ton \emph{xiū}, se prononce en demi-3\textsuperscript{e} ton (descente seule).
\end{exoresolu}

\courssub{Auto-évaluation par enregistrement}

\begin{experience}[S'enregistrer, comparer, progresser]
En binôme ou individuellement, avec un téléphone ou un enregistreur :
\begin{enumerate}
\item \textbf{Écoutez} le modèle (le professeur ou l'enregistrement du manuel) deux fois, sans parler.
\item \textbf{Enregistrez-vous} en lisant l'extrait du dialogue 2, d'abord lentement en exagérant les tons, puis à vitesse normale.
\item \textbf{Réécoutez-vous} et comparez avec le modèle : cochez dans la grille ci-dessous les points réussis.
\item \textbf{Recommencez} une seconde fois en corrigeant un seul point à la fois (d'abord les tons, puis la fluidité).
\item \textbf{Échangez} avec votre voisin : chacun écoute l'enregistrement de l'autre et lui donne un compliment et un conseil, en utilisant les formules du débat : \zh{你说得很好，可是……} \emph{Nǐ shuō de hěn hǎo, kěshì……} « Tu parles très bien, mais… ».
\end{enumerate}
\end{experience}

\begin{center}
\begin{tabularx}{\linewidth}{|X|X|X|}\hline
\rowcolor{popblueL} Critère d'auto-évaluation & Réussi & À retravailler \\ \hline
Les 4 tons sont distincts (haut, montant, creux, descendant) & $\square$ & $\square$ \\ \hline
Deux 3\textsuperscript{e} tons : le premier devient 2\textsuperscript{e} (\zh{很好}, \zh{所以}) & $\square$ & $\square$ \\ \hline
Demi-3\textsuperscript{e} ton devant les autres tons (\zh{北京}, \zh{老师}) & $\square$ & $\square$ \\ \hline
La voyelle ü est correcte (\zh{绿} \emph{lǜ}, \zh{居民} \emph{jūmín}) & $\square$ & $\square$ \\ \hline
Les tons neutres sont courts et légers (\zh{我们}, \zh{觉得}) & $\square$ & $\square$ \\ \hline
\end{tabularx}
\end{center}

\begin{savaistu}[Les tons au baccalauréat]
À l'épreuve orale de chinois du baccalauréat, la \cle{correction des tons} fait officiellement partie des critères d'évaluation, au même titre que le contenu de votre exposé et la fluidité de votre expression. Un exposé riche en idées mais aux tons confus perd des points ; à l'inverse, une prononciation claire valorise immédiatement votre prestation devant le jury. C'est pourquoi les professeurs de chinois du Cameroun insistent sur la gestuelle des tons et l'auto-enregistrement dès la classe de Première : l'oreille et la bouche s'entraînent comme un muscle, régulièrement, un peu chaque jour.
\end{savaistu}

\begin{aretenir}[L'essentiel phonétique de la leçon]
\begin{itemize}
\item Quatre tons + ton neutre : 1\textsuperscript{er} haut et plat, 2\textsuperscript{e} montant, 3\textsuperscript{e} en creux, 4\textsuperscript{e} descendant brusque.
\item Deux 3\textsuperscript{e} tons qui se suivent : le premier devient un 2\textsuperscript{e} ton (\zh{很好} \emph{hěn hǎo} → \emph{hén hǎo} ; \zh{可以} → \emph{kéyǐ} ; \zh{所以} → \emph{suóyǐ}).
\item 3\textsuperscript{e} ton devant un autre ton : demi-3\textsuperscript{e} ton, descente seule (\zh{北京} \emph{Běijīng}).
\item Changements de \zh{一} et \zh{不} : \zh{一} devient \emph{yí} devant un 4\textsuperscript{e} ton et \emph{yì} devant les autres tons ; \zh{不} devient \emph{bú} seulement devant un 4\textsuperscript{e} ton.
\item Homophones du thème : \zh{市}/\zh{是} \emph{shì}, \zh{修}/\zh{休} \emph{xiū} ; paire de tons : \zh{路} \emph{lù} / \zh{绿} \emph{lǜ}.
\item La gestuelle de la main et l'auto-enregistrement sont vos deux meilleurs outils de progrès.
\end{itemize}
\end{aretenir}

\coursec{Jeux de rôle et débat guidé}

Après avoir travaillé les tons du thème dans la section précédente, vous êtes maintenant prêts à les mettre en pratique dans des situations réelles de communication. C'est l'objectif de cette dernière étape : passer de la phrase isolée à l'échange vivant. Nous allons d'abord observer un modèle de débat, puis dégager les règles utiles, avant de vous lancer dans trois jeux de rôle et un grand débat de classe.

\courssub{Observer avant d'agir : un extrait de débat}

\begin{textebox}[Au conseil municipal]
主持人：我们今天讨论一个问题：我们的城市应该先修路还是先建学校？\\
(Zhǔchírén: Wǒmen jīntiān tǎolùn yí ge wèntí: Wǒmen de chéngshì yīnggāi xiān xiū lù háishi xiān jiàn xuéxiào?)\\
L'animateur : Nous discutons aujourd'hui d'une question : notre ville devrait-elle d'abord réparer les routes ou construire des écoles ?\\[4pt]
赞成方：我认为应该先修路，因为路好了，经济才能发展。\\
(Zànchéngfāng: Wǒ rènwéi yīnggāi xiān xiū lù, yīnwèi lù hǎo le, jīngjì cái néng fāzhǎn.)\\
Le camp « pour » : Je pense qu'il faut d'abord réparer les routes, parce que quand les routes sont bonnes, l'économie peut se développer.\\[4pt]
反对方：我不同意。我觉得应该先建学校，因为教育最重要。没有学校，孩子们怎么办？\\
(Fǎnduìfāng: Wǒ bù tóngyì. Wǒ juéde yīnggāi xiān jiàn xuéxiào, yīnwèi jiàoyù zuì zhòngyào. Méiyǒu xuéxiào, háizimen zěnme bàn?)\\
Le camp « contre » : Je ne suis pas d'accord. Je trouve qu'il faut d'abord construire des écoles, parce que l'éducation est le plus important. Sans école, que font les enfants ?\\[4pt]
主持人：还有别的意见吗？居民代表，您怎么看？\\
(Zhǔchírén: Hái yǒu biéde yìjian ma? Jūmín dàibiǎo, nín zěnme kàn?)\\
L'animateur : Y a-t-il d'autres avis ? Représentant des habitants, qu'en pensez-vous ?\\[4pt]
居民：我赞成先建学校。我们学校太旧了，学生太多了。\\
(Jūmín: Wǒ zànchéng xiān jiàn xuéxiào. Wǒmen xuéxiào tài jiù le, xuésheng tài duō le.)\\
L'habitant : Je suis favorable à la construction d'écoles d'abord. Notre école est trop vieille, il y a trop d'élèves.\\[4pt]
主持人：好，谢谢大家。我们来总结一下……\\
(Zhǔchírén: Hǎo, xièxie dàjiā. Wǒmen lái zǒngjié yíxià…)\\
L'animateur : Bien, merci à tous. Faisons la synthèse…
\end{textebox}

Observez la mécanique : chaque prise de parole suit presque toujours le même schéma en trois temps — \cle{avis} (\zh{我认为} / \zh{我赞成} / \zh{我不同意}), \cle{justification} introduite par \zh{因为}, puis éventuellement une \cle{relance} (\zh{您怎么看？}). C'est exactement le triptyque que vous devrez reproduire.

\begin{methode}[Le triptyque de la prise de parole]
Avant de parler, préparez mentalement trois blocs :
\begin{enumerate}
\item \textbf{Avis} : \zh{我认为…} / \zh{我觉得…} / \zh{我赞成…} / \zh{我不同意…}
\item \textbf{Justification} : \zh{因为} + raison (\zh{因为教育最重要。})
\item \textbf{Relance} : \zh{您怎么看？} / \zh{还有别的意见吗？} / \zh{你同意吗？}
\end{enumerate}
Conseil d'or : \textbf{parler lentement vaut mieux que parler vite avec des tons faux}. Un débit lent et des tons corrects donnent une impression de maîtrise ; un débit rapide et des tons aplatis donnent l'impression inverse. Comptez une courte pause après chaque virgule.
\end{methode}

\courssub{La question alternative avec 还是}

La question centrale de notre débat utilise une structure grammaticale essentielle :

\begin{exemplebox}[Question alternative]
\zh{我们的城市应该先修路还是先建学校？}\\
\emph{Wǒmen de chéngshì yīnggāi xiān xiū lù háishi xiān jiàn xuéxiào?}\\
« Notre ville devrait-elle d'abord réparer les routes ou construire des écoles ? »
\end{exemplebox}

\begin{propriete}[还是 (háishi) et 或者 (huòzhě) : deux « ou » différents]
\begin{itemize}
\item \zh{还是} (háishi) s'emploie dans les \textbf{questions alternatives} : il propose un choix entre deux options. Structure : \textbf{A + 还是 + B + ？}
\item \zh{或者} (huòzhě) s'emploie dans les \textbf{phrases affirmatives} : « ou » au sens de l'alternative non interrogative. Structure : \textbf{A + 或者 + B + 。}
\end{itemize}
Exemples :\\
\zh{你想喝茶还是喝咖啡？} \emph{Nǐ xiǎng hē chá háishi hē kāfēi?} « Tu veux boire du thé ou du café ? » (question)\\
\zh{周末我想看电影或者去市场。} \emph{Zhōumò wǒ xiǎng kàn diànyǐng huòzhě qù shìchǎng.} « Le week-end, je veux voir un film ou aller au marché. » (affirmation)\\
\textbf{Exception} : dans une question rapportée (style indirect), on garde \zh{还是} : \zh{他问我应该先修路还是建学校。} \emph{Tā wèn wǒ yīnggāi xiān xiū lù háishi jiàn xuéxiào.} « Il m'a demandé s'il fallait d'abord réparer les routes ou construire des écoles. »
\end{propriete}

\begin{center}
\begin{tabularx}{\linewidth}{|X|X|X|}
\hline
\rowcolor{popblueL} Structure & Exemple (caractères + pinyin) & Traduction \\
\hline
Question : A + 还是 + B ? & \zh{先修路还是先建学校？} Xiān xiū lù háishi xiān jiàn xuéxiào? & D'abord les routes ou les écoles ? \\
\hline
Affirmation : A + 或者 + B & \zh{我们可以修路或者建市场。} Wǒmen kěyǐ xiū lù huòzhě jiàn shìchǎng. & Nous pouvons réparer les routes ou construire un marché. \\
\hline
Donner son avis & \zh{我赞成先建学校。} Wǒ zànchéng xiān jiàn xuéxiào. & Je suis favorable aux écoles d'abord. \\
\hline
Relancer & \zh{您怎么看？} Nín zěnme kàn? & Qu'en pensez-vous ? \\
\hline
\end{tabularx}
\end{center}

\begin{attention}[« Je suis pour » ne se traduit pas mot à mot]
Ne dites jamais \zh{我是为了} pour « je suis pour » : c'est un calque du français qui ne veut rien dire en chinois. Dites simplement :
\begin{itemize}
\item \zh{我同意。} \emph{Wǒ tóngyì.} « Je suis d'accord. »
\item \zh{我赞成。} \emph{Wǒ zànchéng.} « Je suis favorable. »
\item \zh{我支持。} \emph{Wǒ zhīchí.} « Je soutiens. »
\item \zh{我反对。} \emph{Wǒ fǎnduì.} « Je suis contre. »
\end{itemize}
Autre piège : \zh{因为} (parce que) et \zh{所以} (donc) s'utilisent ensemble en chinois (\zh{因为路不好，所以经济不发展。}) alors qu'en français « parce que… donc » est une faute de syntaxe (redondance). Ne traduisez jamais « donc » seul par \zh{因为}.
\end{attention}

\courssub{Jeu de rôle 1 : le maire face aux journalistes}

\begin{experience}[La conférence de presse du maire]
\textbf{Répartition} : un élève joue le maire (\zh{市长} shìzhǎng), trois ou quatre élèves jouent les journalistes (\zh{记者} jìzhě).\\
\textbf{Situation} : le maire de votre commune présente le projet de l'année (par exemple : réparer la route de Mvog-Ada à Yaoundé) et répond aux questions.\\
\textbf{Phrases obligatoires pour le maire} :
\begin{enumerate}
\item Avis : \zh{我认为这个项目很重要。} \emph{Wǒ rènwéi zhège xiàngmù hěn zhòngyào.} « Je pense que ce projet est très important. »
\item Justification : \zh{因为这条路太旧了，车走得很慢。} \emph{Yīnwèi zhè tiáo lù tài jiù le, chē zǒu de hěn màn.} « Parce que cette route est trop vieille, les voitures roulent lentement. »
\item Relance : \zh{大家有什么问题吗？} \emph{Dàjiā yǒu shénme wèntí ma?} « Avez-vous des questions ? »
\end{enumerate}
\textbf{Phrases obligatoires pour les journalistes} : \zh{这个项目要多少钱？} \emph{Zhège xiàngmù yào duōshao qián?} « Combien coûtera ce projet ? » ; \zh{什么时候开始？} \emph{Shénme shíhou kāishǐ?} « Quand est-ce que ça commence ? » ; \zh{居民同意吗？} \emph{Jūmín tóngyì ma?} « Les habitants sont-ils d'accord ? »
\end{experience}

\courssub{Jeu de rôle 2 : la réunion du conseil}

\begin{experience}[Trois habitants, trois priorités]
\textbf{Répartition} : un animateur (\zh{主持人}) et trois habitants (\zh{居民} jūmín). Chaque habitant défend une priorité : la route (\zh{路}), l'école (\zh{学校}) ou le marché (\zh{市场}).\\
\textbf{Déroulement} : l'animateur pose la question \zh{我们的城市应该先做什么？} \emph{Wǒmen de chéngshì yīnggāi xiān zuò shénme?} « Que doit faire notre ville en premier ? », puis chaque habitant parle une minute.\\
\textbf{Fiches d'appui (3 phrases obligatoires par rôle)} :
\begin{itemize}
\item Défenseur de la route : \zh{我觉得应该先修路。} / \zh{因为路不好，农民的东西运不出去。} \emph{Yīnwèi lù bù hǎo, nóngmín de dōngxi yùn bù chūqù.} « Parce que les routes sont mauvaises, les produits des paysans ne peuvent pas être transportés. » / \zh{你们同意吗？} \emph{Nǐmen tóngyì ma?}
\item Défenseur de l'école : \zh{我认为应该先建学校。} / \zh{因为没有教育，就没有未来。} \emph{Yīnwèi méiyǒu jiàoyù, jiù méiyǒu wèilái.} « Parce que sans éducation, il n'y a pas d'avenir. » / \zh{您怎么看？} \emph{Nín zěnme kàn?}
\item Défenseur du marché : \zh{我赞成先建市场。} / \zh{因为市场好了，大家的生活就方便了。} \emph{Yīnwèi shìchǎng hǎo le, dàjiā de shēnghuó jiù fāngbiàn le.} « Parce qu'avec un bon marché, la vie de tous devient pratique. » / \zh{还有别的意见吗？} \emph{Hái yǒu biéde yìjian ma?}
\end{itemize}
\end{experience}

\courssub{Jeu de rôle 3 : l'interview par un journaliste chinois}

\begin{experience}[Un journaliste de Beijing interviewe un élève camerounais]
\textbf{Répartition} : par binômes. L'un est le journaliste chinois en reportage à Douala, l'autre un élève camerounais.\\
\textbf{Questions du journaliste} : \zh{你的城市叫什么名字？} \emph{Nǐ de chéngshì jiào shénme míngzi?} « Comment s'appelle ta ville ? » ; \zh{你们的市长是谁选的？} \emph{Nǐmen de shìzhǎng shì shéi xuǎn de?} « Qui élit votre maire ? » ; \zh{你的城市有什么问题？} \emph{Nǐ de chéngshì yǒu shénme wèntí?} « Quels problèmes y a-t-il dans ta ville ? »\\
\textbf{Phrases obligatoires de l'élève} : \zh{我住在雅温得。} \emph{Wǒ zhù zài Yǎwēndé.} « J'habite à Yaoundé. » ; \zh{我们的市长是居民选的。} \emph{Wǒmen de shìzhǎng shì jūmín xuǎn de.} « Notre maire est élu par les habitants. » ; \zh{我觉得我们应该先修路，因为…} (avis + justification).\\
\textbf{Variante} : inversez les rôles, puis changez de ville (Douala, Buea, Garoua).
\end{experience}

\courssub{Débat guidé en classe entière}

\begin{experience}[Le grand débat : routes ou écoles ?]
\textbf{Question} : \zh{我们的城市应该先修路还是先建学校？}\\
\textbf{Organisation} : la classe est divisée en deux camps — \zh{赞成方} (zànchéngfāng, le camp « pour » les routes) et \zh{反对方} (fǎnduìfāng, le camp « contre », pour les écoles). Un élève est l'animateur (\zh{主持人}), un autre fait la synthèse finale (\zh{总结} zǒngjié).\\
\textbf{Déroulement} : (1) chaque camp prépare trois arguments en cinq minutes ; (2) l'animateur donne la parole en alternance ; (3) chaque orateur applique le triptyque avis + \zh{因为} + relance ; (4) le rapporteur conclut : \zh{今天我们讨论了… 大部分人认为…} \emph{Jīntiān wǒmen tǎolùn le… Dàbùfēn rén rènwéi…} « Aujourd'hui nous avons discuté de… La plupart pensent que… »
\end{experience}

\begin{popfigure}
\begin{tikzpicture}[node distance=0.9cm and 1.6cm,
box/.style={draw=popblue, fill=popblueL, rounded corners, align=center, minimum height=0.9cm, font=\small},
box2/.style={draw=poporange, fill=poporangeL, rounded corners, align=center, minimum height=0.9cm, font=\small},
box3/.style={draw=popgreen, fill=popgreenL, rounded corners, align=center, minimum height=0.9cm, font=\small},
arr/.style={-{Stealth[length=2.5mm]}, thick, popdark}]
\node[box, align=center] (anim){主持人\\ \emph{animateur}};
\node[box2, below left=of anim, align=center] (pour){赞成方\\ \emph{pour : 修路}};
\node[box2, below right=of anim, align=center] (contre){反对方\\ \emph{contre : 建学校}};
\node[box3, below=2.6cm of anim, align=center] (synth){总结\\ \emph{synthèse}};
\draw[arr] (anim) -- (pour);
\draw[arr] (anim) -- (contre);
\draw[arr, poporange] (pour.east) to[bend left=15] node[above, font=\scriptsize]{relance} (contre.west);
\draw[arr, poporange] (contre.west) to[bend left=15] node[below, font=\scriptsize]{réponse} (pour.east);
\draw[arr] (pour.south) -- (synth.west);
\draw[arr] (contre.south) -- (synth.east);
\end{tikzpicture}
\legende{Organigramme du débat : l'animateur distribue la parole, les deux camps échangent, le rapporteur conclut.}
\end{popfigure}

\begin{exoresolu}[Construire une prise de parole complète]
Énoncé : vous êtes le défenseur du marché. Construisez une prise de parole de trois phrases (avis + justification avec \zh{因为} + relance) en utilisant le vocabulaire du thème.
\tcblower
Solution détaillée :
\begin{enumerate}
\item Avis : \zh{我觉得应该先建市场。} \emph{Wǒ juéde yīnggāi xiān jiàn shìchǎng.} « Je trouve qu'il faut d'abord construire un marché. »
\item Justification : \zh{因为我们城市没有大市场，买东西很不方便。} \emph{Yīnwèi wǒmen chéngshì méiyǒu dà shìchǎng, mǎi dōngxi hěn bù fāngbiàn.} « Parce que notre ville n'a pas de grand marché, faire ses courses est très peu pratique. »
\item Relance : \zh{居民们，你们同意吗？} \emph{Jūmínmen, nǐmen tóngyì ma?} « Habitants, êtes-vous d'accord ? »
\end{enumerate}
Vérification des tons : \zh{市场} shìchǎng (4\textsuperscript{e} + 3\textsuperscript{e} ton), \zh{方便} fāngbiàn (1\textsuperscript{er} + 4\textsuperscript{e} ton). Prononcez lentement, en marquant bien le 4\textsuperscript{e} ton descendant de \zh{建} (jiàn). Attention au 3\textsuperscript{e} ton de \zh{反对} (fǎnduì $\to$ fánduì en sandhi) et de \zh{主持人} (zhǔchírén $\to$ zhúchírén en sandhi).
\end{exoresolu}

\courssub{Regards croisés : Cameroun et Chine}

\begin{savaistu}[Communes camerounaises et quartiers chinois]
Au Cameroun, les collectivités locales décentralisées sont les \textbf{communes}, dirigées par des \textbf{maires élus} au suffrage universel ; les grandes villes comme Douala et Yaoundé sont organisées en \textbf{communautés urbaines} regroupant plusieurs communes d'arrondissement. En Chine, il n'y a pas de maire élu au même sens : chaque quartier, appelé \zh{社区} (shèqū), possède un \textbf{comité de résidents}, \zh{居民委员会} (jūmín wěiyuánhuì), qui gère la vie quotidienne du quartier (propreté, médiation, petits services). Savoir citer \zh{居民委员会} dans votre exposé — et expliquer la différence avec nos maires élus — est un argument culturel qui impressionne toujours le jury au baccalauréat.
\end{savaistu}

\begin{center}
\begin{tabularx}{\linewidth}{|X|X|X|X|}
\hline
\rowcolor{popblueL} Caractères & Pinyin & Nature & Traduction \\
\hline
市长 & shìzhǎng & nom & maire \\
\hline
居民 & jūmín & nom & habitant, résident \\
\hline
社区 & shèqū & nom & quartier, communauté \\
\hline
居民委员会 & jūmín wěiyuánhuì & nom & comité de résidents \\
\hline
赞成 & zànchéng & verbe & approuver, être favorable \\
\hline
反对 & fǎnduì & verbe & s'opposer, être contre \\
\hline
主持人 & zhǔchírén & nom & animateur, présentateur \\
\hline
总结 & zǒngjié & verbe / nom & résumer ; synthèse \\
\hline
项目 & xiàngmù & nom & projet \\
\hline
记者 & jìzhě & nom & journaliste \\
\hline
修路 & xiū lù & verbe + obj & réparer les routes \\
\hline
建学校 & jiàn xuéxiào & verbe + obj & construire des écoles \\
\hline
市场 & shìchǎng & nom & marché \\
\hline
农民 & nóngmín & nom & paysan \\
\hline
未来 & wèilái & nom & avenir \\
\hline
方便 & fāngbiàn & adj & pratique, commode \\
\hline
雅温得 & Yǎwēndé & n.p. & Yaoundé \\
\hline
选 & xuǎn & verbe & élire \\
\hline
问题 & wèntí & nom & problème, question \\
\hline
多少钱 & duōshao qián & locution & combien d'argent \\
\hline
时候 & shíhou & nom & moment, quand \\
\hline
开始 & kāishǐ & verbe & commencer \\
\hline
\end{tabularx}
\end{center}

\courssub{Grille d'évaluation orale}

Chaque prestation (jeu de rôle ou débat) est notée sur 20 points selon cinq critères :

\begin{center}
\begin{tabularx}{\linewidth}{|X|X|}
\hline
\rowcolor{popblueL} Critère & Points \\
\hline
Prononciation des tons (tons justes, 3\textsuperscript{e} ton bien creusé, sandhi respecté) & 4 pts \\
\hline
Vocabulaire du thème (市长, 居民, 社区, 项目, 修路, 建学校…) & 4 pts \\
\hline
Formules fonctionnelles (avis, 因为, relance, 还是/或者) & 4 pts \\
\hline
Fluidité et interaction (écoute, réponse, débit régulier) & 4 pts \\
\hline
Contenu (arguments clairs et pertinents) & 4 pts \\
\hline
\end{tabularx}
\end{center}

\begin{exercice}[À vous de jouer]
En binômes, préparez en cinq minutes un mini-débat sur la question : \zh{我们应该先建市场还是先修路？} \emph{Wǒmen yīnggāi xiān jiàn shìchǎng háishi xiān xiū lù?} « Faut-il d'abord construire un marché ou réparer les routes ? » Chacun défend une position avec le triptyque complet (avis + \zh{因为} + relance), puis présentez devant la classe. Vos camarades vous évaluent avec la grille ci-dessus.
\end{exercice}

\begin{aretenir}[L'essentiel de la section]
\begin{itemize}
\item Toute prise de parole suit le triptyque : \textbf{avis} (\zh{我认为/我赞成}) + \textbf{justification} (\zh{因为}) + \textbf{relance} (\zh{您怎么看？}).
\item « Ou » se dit \zh{还是} dans les questions, \zh{或者} dans les affirmations (sauf question rapportée).
\item « Je suis pour » se dit \zh{我赞成}, \zh{我同意} ou \zh{我支持}, jamais un calque du français (\zh{我是为了}).
\item \zh{因为} et \zh{所以} s'associent en chinois ; « parce que… donc » est fautif en français.
\item Mieux vaut parler lentement avec des tons corrects que vite avec des tons faux.
\item Le débat s'organise autour de quatre rôles : \zh{主持人}, \zh{赞成方}, \zh{反对方}, \zh{总结}.
\item Connaître la différence institutionnelle : maire élu (Cameroun) vs comité de résidents \zh{居民委员会} (Chine).
\end{itemize}
\end{aretenir}

\newpage

\coursec{Activité d'intégration}

\coursec{Leçon 20 — Expression orale : collectivités locales}

\courssub{Objectifs de la leçon}

À l'issue de cette leçon, tu seras capable de :
\begin{itemize}
\item mobiliser le \cle{lexique des collectivités locales} (mairie, conseil municipal, habitants, projets d'infrastructure, vote) dans un échange oral authentique ;
\item présenter oralement un exposé simple et structuré en chinois (présentation de la situation, énumération des choix, conclusion) ;
\item exprimer ton opinion, la justifier et relancer l'interlocuteur à l'aide de \cle{formules fonctionnelles} figées (\zh{我觉得…}, \zh{因为…所以…}, \zh{你有什么看法？}, \zh{你同意吗？}) ;
\item participer à un \cle{débat simulé} en respectant les tours de parole : prendre la parole, argumenter, réagir, conclure ;
\item prononcer correctement les \cle{tons} des mots-clés du thème (paires tonales 4-3, 2-4, 1-2, ton neutre) ;
\item évaluer ta production et celle de tes pairs à l'aide d'une grille critériée (tons, vocabulaire, structures, fluidité, contenu).
\end{itemize}

\courssub{Situation de communication : Le conseil municipal de Bafoussam}

La ville de \textbf{Bafoussam} (Cameroun) est jumelée (\zh{友好城市}, \emph{yǒuhǎo chéngshì}) avec la ville chinoise de \textbf{Kunming} (昆明, \emph{Kūnmíng}). Pour célébrer ce jumelage, la mairie organise une séance publique du conseil municipal : il faut désigner \emph{le projet prioritaire de l'année}. Trois projets sont en lice : réparer la route principale, construire un nouveau marché, rénover l'école primaire. Le budget ne permet d'en financer qu'un seul. La classe simule cette séance : le maire présente les projets, les habitants défendent leurs priorités, les journalistes interrogent, l'animateur assure la fluidité des échanges, et le public vote.

\begin{textebox}[Affiche officielle de la mairie — 市政府公告]
市政府公告 \emph{Shì zhèngfǔ gōnggào} « Annonce de la mairie » \\
各位居民：\emph{Gèwèi jūmín :} « Chers habitants, » \\
我们市今年有一个重要项目。有三个选择：\emph{Wǒmen shì jīnnián yǒu yí ge zhòngyào xiàngmù. Yǒu sān ge xuǎnzé :} « Notre ville a un projet important cette année. Il y a trois options : » \\
一、修路；二、建市场；三、修学校。\emph{Yī, xiū lù ; èr, jiàn shìchǎng ; sān, xiū xuéxiào.} « 1. Réparer la route ; 2. Construire un marché ; 3. Rénover l'école. » \\
欢迎大家来开会，说说你们的看法！\emph{Huānyíng dàjiā lái kāihuì, shuōshuo nǐmen de kànfǎ !} « Tout le monde est invité à la réunion pour donner son avis ! » \\
时间：星期五下午三点。地点：市政府大厅。\emph{Shíjiān : Xīngqīwǔ xiàwǔ sān diǎn. Dìdiǎn : Shì zhèngfǔ dàtīng.} « Date : vendredi 15 h. Lieu : grand hall de la mairie. »
\end{textebox}

\begin{savaistu}[Le jumelage Cameroun–Chine : 友好城市 / 姐妹城市]
Plusieurs villes camerounaises (Yaoundé, Douala, Bafoussam, Buea…) entretiennent des relations de \cle{jumelage} (\zh{友好城市} \emph{yǒuhǎo chéngshì} « villes amies » ou \zh{姐妹城市} \emph{jiěmèi chéngshì} « villes sœurs ») avec des villes chinoises (Pékin, Shanghai, Kunming, Shenzhen…). Ces accords portent sur la culture, l'éducation, la santé et les infrastructures. En Chine comme au Cameroun, la \cle{commune} (\zh{市} \emph{shì} pour une ville-préfecture, \zh{县} \emph{xiàn} pour un district) est l'échelon administratif le plus proche des citoyens : c'est là que se décident les routes, les marchés, les écoles et la gestion des déchets. Le \cle{maire} (\zh{市长} \emph{shìzhǎng}) y préside le \cle{conseil municipal} (\zh{市政会} \emph{shìzhènghuì} ou \zh{人民代表大会} \emph{rénmín dàibiǎo dàhuì}).
\end{savaistu}

\courssub{Lexique thématique : Les collectivités locales}

\begin{center}\begin{tabularx}{\linewidth}{|X|X|X|X|}\hline
\rowcolor{popblueL} Caractères & Pinyin & Nature & Traduction \\ \hline
市长 & shìzhǎng & nom & maire \\ \hline
副市长 & fù shìzhǎng & nom & maire adjoint \\ \hline
居民 & jūmín & nom & habitant, résident \\ \hline
记者 & jìzhě & nom & journaliste \\ \hline
主持人 & zhǔchírén & nom & animateur, modérateur \\ \hline
市政府 & shì zhèngfǔ & nom & mairie, gouvernement municipal \\ \hline
市政会 & shìzhènghuì & nom & conseil municipal \\ \hline
项目 & xiàngmù & nom & projet \\ \hline
预算 & yùsuàn & nom & budget \\ \hline
优先事项 & yōuxiān shìxiàng & nom & priorité \\ \hline
修路 & xiū lù & v. + obj. & réparer / construire une route \\ \hline
建市场 & jiàn shìchǎng & v. + obj. & construire un marché \\ \hline
修学校 & xiū xuéxiào & v. + obj. & rénover / réparer une école \\ \hline
学校 & xuéxiào & nom & école \\ \hline
市场 & shìchǎng & nom & marché \\ \hline
路 & lù & nom & route, chemin \\ \hline
看法 & kànfǎ & nom & opinion, point de vue \\ \hline
重要 & zhòngyào & adj. & important \\ \hline
选择 & xuǎnzé & nom / v. & choix / choisir \\ \hline
投票 & tóupiào & v. / nom & voter / vote \\ \hline
通过 & tōngguò & v. & adopter (une motion), passer \\ \hline
友好城市 & yǒuhǎo chéngshì & nom & ville jumelée (ville amie) \\ \hline
\end{tabularx}\end{center}

\begin{propriete}[Point de grammaire : \zh{修} vs \zh{建}]
\begin{itemize}
\item \zh{修} \emph{xiū} : réparer, rénover, entretenir (ex. \zh{修路} réparer la route, \zh{修学校} rénover l'école, \zh{修桥} réparer le pont).
\item \zh{建} \emph{jiàn} : construire, édifier du neuf (ex. \zh{建市场} construire un marché, \zh{建医院} construire un hôpital, \zh{建设} construction/développement).
\item \textbf{Attention} : On peut dire \zh{建学校} (construire une nouvelle école) mais \zh{修学校} implique une rénovation de l'existant. Dans l'affiche, \zh{修学校} = rénover.
\end{itemize}
\end{propriete}

\courssub{Dialogue modèle 1 : Le maire présente les projets (Exposé structuré)}

Ce dialogue modèle montre comment structurer une prise de parole officielle : ouverture \textrightarrow énumération \textrightarrow invitation à débattre.

\begin{textebox}[Discours d'ouverture du maire]
市长：各位居民，大家好！\emph{Shìzhǎng : Gèwèi jūmín, dàjiā hǎo !} « Maire : Chers habitants, bonjour à tous ! » \\
欢迎大家参加今天的市政会。\emph{Huānyíng dàjiā cānjiā jīntiān de shìzhènghuì.} « Bienvenue à la séance du conseil municipal d'aujourd'hui. » \\
我们市今年有一个重要项目，预算有限，只能选一个。\emph{Wǒmen shì jīnnián yǒu yí ge zhòngyào xiàngmù, yùsuàn yǒuxiàn, zhǐ néng xuǎn yí ge.} « Notre ville a un projet important cette année, le budget est limité, nous ne pouvons en choisir qu'un. » \\
我们有三个选择：\emph{Wǒmen yǒu sān ge xuǎnzé :} « Nous avons trois options : » \\
第一，修路。\emph{Dì-yī, xiū lù.} « Premièrement, réparer la route. » \\
第二，建市场。\emph{Dì-èr, jiàn shìchǎng.} « Deuxièmement, construire un marché. » \\
第三，修学校。\emph{Dì-sān, xiū xuéxiào.} « Troisièmement, rénover l'école. » \\
请大家发表看法。\emph{Qǐng dàjiā fābiǎo kànfǎ.} « Je vous invite à donner vos avis. »
\end{textebox}

\begin{propriete}[Structure de l'énumération formelle]
\begin{center}\begin{tabularx}{\linewidth}{|X|X|X|}\hline
\rowcolor{popblueL} Structure & Exemple (Caractères + Pinyin) & Traduction \\ \hline
第一，…；第二，…；第三，… & \zh{第一，修路；第二，建市场；第三，修学校。} \emph{Dì-yī, xiū lù; dì-èr, jiàn shìchǎng; dì-sān, xiū xuéxiào.} & Premièrement…, deuxièmement…, troisièmement… \\ \hline
有三个选择：A，B，C。 & \zh{我们有三个选择：修路，建市场，修学校。} \emph{Wǒmen yǒu sān ge xuǎnzé : xiū lù, jiàn shìchǎng, xiū xuéxiào.} & Nous avons trois options : A, B, C. \\ \hline
只能选一个。 & \zh{预算有限，只能选一个。} \emph{Yùsuàn yǒuxiàn, zhǐ néng xuǎn yí ge.} & Budget limité, on ne peut en choisir qu'un. \\ \hline
\end{tabularx}\end{center}
\end{propriete}

\courssub{Formules fonctionnelles : Opiner, justifier, relancer, (dés)accorder}

Ces formules sont des \cle{blocs figés} à mémoriser et à automatiser pour la fluidité du débat.

\begin{propriete}[Boîte à outils orale — 表达观点与互动]
\begin{center}\begin{tabularx}{\linewidth}{|X|X|X|}\hline
\rowcolor{popblueL} Fonction & Formule chinoise (Caractères + Pinyin) & Équivalent français \\ \hline
\bfseries Donner son avis & \zh{我觉得…} \emph{Wǒ juéde…} & Je pense que… / À mon avis… \\ \hline
& \zh{我认为…} \emph{Wǒ rènwéi…} (plus formel) & J'estime que… / Je considère que… \\ \hline
\bfseries Justifier (Cause \textrightarrow Conséquence) & \zh{因为…，所以…} \emph{Yīnwèi…, suǒyǐ…} & Parce que…, donc… \\ \hline
& \zh{因为…} \emph{Yīnwèi…} (seul, oral) & Parce que… (réponse à \zh{为什么？}) \\ \hline
\bfseries Superlatif (Priorité) & \zh{最} + adj. \emph{zuì + adj.} & Le plus… (ex. \zh{最重要} le plus important) \\ \hline
\bfseries Nuancer / Ajouter & \zh{也} \emph{yě} & Aussi / Également (ex. \zh{我也觉得…} Moi aussi je pense que…) \\ \hline
\bfseries Demander l'avis (Relance) & \zh{你有什么看法？} \emph{Nǐ yǒu shénme kànfǎ ?} & Quel est ton avis ? \\ \hline
& \zh{你怎么看？} \emph{Nǐ zěnme kàn ?} & Comment tu vois les choses ? \\ \hline
& \zh{为什么？} \emph{Wèishénme ?} & Pourquoi ? \\ \hline
\bfseries (Dés)accorder & \zh{我同意。} \emph{Wǒ tóngyì.} & Je suis d'accord. \\ \hline
& \zh{我不同意。} \emph{Wǒ bù tóngyì.} & Je ne suis pas d'accord. \\ \hline
& \zh{有道理。} \emph{Yǒu dàoli.} & C'est logique / Il y a du sens. \\ \hline
\bfseries Relancer un timide & \zh{你呢？} \emph{Nǐ ne ?} & Et toi ? \\ \hline
& \zh{你同意吗？} \emph{Nǐ tóngyì ma ?} & Tu es d'accord ? \\ \hline
\bfseries Conclure / Voter & \zh{现在我们投票。} \emph{Xiànzài wǒmen tóupiào.} & Maintenant nous votons. \\ \hline
& \zh{我们选第X个项目。} \emph{Wǒmen xuǎn dì-X ge xiàngmù.} & Nous choisissons le projet X. \\ \hline
\end{tabularx}\end{center}
\end{propriete}

\begin{attention}[Pièges fréquents pour francophones — \cle{Ordre des mots} et \cle{Tons}]
\begin{itemize}
\item \textbf{Ordre Cause/Effet} : En français « Je choisis l'école \emph{parce que} les enfants en ont besoin ». En chinois : \zh{因为孩子们需要学校，所以我选学校。} (\emph{Yīnwèi háizi xūyào xuéxiào, suǒyǐ wǒ xuǎn xuéxiào.}) \textrightarrow La cause \zh{因为} précède \textbf{toujours} la conséquence \zh{所以}. Ne commence pas par la conclusion.
\item \textbf{\zh{觉得} vs \zh{认为}} : \zh{觉得} \emph{juéde} (2 + ton neutre) = ressenti, avis personnel (oral). \zh{认为} \emph{rènwéi} (4-2) = opinion réfléchie, formel (écrit/oral formel). Ne prononce pas \emph{juédé} (2-4).
\item \textbf{Minimal pairs tons/initiales} :
\begin{itemize}
\item \zh{市长} \emph{shìzhǎng} (4-3) « maire » vs \zh{市场} \emph{shìchǎng} (4-3) « marché » \textrightarrow Distingue \emph{zh} (rétroflexe) et \emph{ch} (aspiré).
\item \zh{学校} \emph{xuéxiào} (2-4) « école » vs \zh{选择} \emph{xuǎnzé} (3-2) « choix ».
\item \zh{重要} \emph{zhòngyào} (4-4) vs \zh{项目} \emph{xiàngmù} (4-4).
\end{itemize}
\item \textbf{Classificateur \zh{个}} : \zh{一个项目} \emph{yí ge xiàngmù} (ton 2 \emph{yí} devant ton 4), \zh{三个选择} \emph{sān ge xuǎnzé}.
\item \textbf{\zh{了} final} : Marque le changement d'état / action accomplie. \zh{会议结束了} \emph{Huìyì jiéshù le} « La réunion est finie ». \zh{我们选了第二个} \emph{Wǒmen xuǎn le dì-èr ge} « Nous avons choisi le deuxième ». Ne l'oublie pas au passé/accompli.
\end{itemize}
\end{attention}

\courssub{Travail phonétique : Les tons du thème (Atelier prononciation)}

\begin{methode}[Entraînement ciblé sur les paires tonales du débat]
\begin{enumerate}
\item \textbf{Écoute-Modèle} : Le professeur (ou l'audio) lit la liste ci-dessous. Les élèves répètent en chœur, puis individuellement.
\item \textbf{Focus paires critiques} :
\begin{itemize}
\item \textbf{4-3} : \zh{市长} (shìzhǎng), \zh{市场} (shìchǎng), \zh{看法} (kànfǎ), \zh{项目} (xiàngmù — 4-4 mais proche), \zh{重要} (zhòngyào — 4-4).
\item \textbf{2-4} : \zh{学校} (xuéxiào), \zh{投票} (tóupiào), \zh{居民} (jūmín — 1-2), \zh{记者} (jìzhě — 4-3).
\item \textbf{Ton neutre} : \zh{觉得} (juéde), \zh{什么} (shénme), \zh{吗} (ma), \zh{呢} (ne), \zh{的} (de).
\end{itemize}
\item \textbf{Exercice de discrimination} : Le prof dit un mot, l'élève montre la carte \emph{4-3} ou \emph{2-4} ou \emph{Neutre}.
\item \textbf{Phonétique intégrée} : Relire le dialogue modèle 1 en insistant sur la \cle{proodie de l'énumération} (montée sur \zh{第一}、\zh{第二}、\zh{第三}, chute finale sur \zh{看法}).
\end{enumerate}
\end{methode}

\begin{popfigure}
\begin{tikzpicture}[
  node distance=1.2cm and 1.5cm,
  box/.style={draw=popblue, thick, rounded corners, fill=popblueL, text width=2.8cm, align=center, font=\small},
  arrow/.style={-Stealth, thick, popdark}
]
\node[box, align=center] (maire){Maire \\ \zh{市长}\\\emph{shìzhǎng}\\(4-3)};
\node[box, right=of maire, align=center] (hab1){Habitant 1 \\ \zh{居民}\\\emph{jūmín}\\(1-2)};
\node[box, right=of hab1, align=center] (hab2){Habitant 2 \\ \zh{居民}\\\emph{jūmín}};
\node[box, below=of maire, align=center] (jour){Journaliste \\ \zh{记者}\\\emph{jìzhě}\\(4-3)};
\node[box, below=of hab1, align=center] (anim){Animateur \\ \zh{主持人}\\\emph{zhǔchírén}\\(3-2-2)};
\node[box, below=of hab2, align=center] (public){Public / Jury \\ \zh{观众}\\\emph{guānzhòng}\\(1-4)};

\draw[arrow] (maire) -- node[above, sloped, font=\tiny] {Présente projets} (hab1);
\draw[arrow] (maire) -- (hab2);
\draw[arrow] (hab1) -- node[right, font=\tiny] {Donne avis \zh{我觉得…}} (anim);
\draw[arrow] (hab2) -- (anim);
\draw[arrow] (jour) -- node[left, font=\tiny] {Interroge \zh{为什么？}} (maire);
\draw[arrow] (jour) -- (hab1);
\draw[arrow] (anim) -- node[below, font=\tiny, align=center]{Relance \zh{你呢？} \\ Conclut \zh{投票}} (public);
\draw[arrow, dashed] (public) -- node[below, font=\tiny] {Vote \zh{投票}} (anim);
\end{tikzpicture}
\legende{Schéma des interactions et mots-clés tonaux du conseil municipal simulé}
\end{popfigure}

\courssub{Dialogue modèle 2 : Débat simulé (Version complète annotée)}

Ce dialogue intègre toutes les formules fonctionnelles. Il sert de \cle{modèle cible} pour le jeu de rôle final.

\begin{textebox}[Conseil municipal simulé — Débat complet]
\zh{主持人：各位居民，大家好！欢迎来到今天的市政会。今天我们要选一个重要项目。} \\
\emph{Zhǔchírén : Gèwèi jūmín, dàjiā hǎo ! Huānyíng lái dào jīntiān de shìzhènghuì. Jīntiān wǒmen yào xuǎn yí ge zhòngyào xiàngmù.} \\
« Animateur : Chers habitants, bonjour ! Bienvenue à la séance du conseil municipal d'aujourd'hui. Aujourd'hui nous devons choisir un projet important. »

\zh{市长：我们有三个选择：第一，修路；第二，建市场；第三，修学校。预算有限，只能选一个。请大家发表看法。} \\
\emph{Shìzhǎng : Wǒmen yǒu sān ge xuǎnzé : dì-yī, xiū lù ; dì-èr, jiàn shìchǎng ; dì-sān, xiū xuéxiào. Yùsuàn yǒuxiàn, zhǐ néng xuǎn yí ge. Qǐng dàjiā fābiǎo kànfǎ.} \\
« Maire : Trois options : 1. réparer la route ; 2. construire un marché ; 3. rénover l'école. Budget limité, un seul choix. Donnez vos avis. »

\zh{居民一：我觉得修学校最重要，因为孩子们需要好学校。因为学校好了，所以城市的未来就好了。} \\
\emph{Jūmín yī : Wǒ juéde xiū xuéxiào zuì zhòngyào, yīnwèi háizi men xūyào hǎo xuéxiào. Yīnwèi xuéxiào hǎo le, suǒyǐ chéngshì de wèilái jiù hǎo le.} \\
« Habitant 1 : Je pense que rénover l'école est le plus important, car les enfants ont besoin d'une bonne école. Parce que l'école va bien, l'avenir de la ville ira bien. »

\zh{记者：市长，你有什么看法？你选哪一个？为什么？} \\
\emph{Jìzhě : Shìzhǎng, nǐ yǒu shénme kànfǎ ? Nǐ xuǎn nǎ yí ge ? Wèishénme ?} \\
« Journaliste : Monsieur le maire, quel est votre avis ? Lequel choisissez-vous ? Pourquoi ? »

\zh{市长：我觉得建市场也很好，因为很多居民要买东西、卖东西，市场能带来收入。} \\
\emph{Shìzhǎng : Wǒ juéde jiàn shìchǎng yě hěn hǎo, yīnwèi hěn duō jūmín yào mǎi dōngxi, mài dōngxi, shìchǎng néng dàilái shōurù.} \\
« Maire : Je pense que construire un marché est bien aussi, car beaucoup d'habitants achètent et vendent, le marché apporte des revenus. »

\zh{居民二：我不同意市长。我觉得修路最重要。因为路坏了，车走不了，市场也没用。} \\
\emph{Jūmín èr : Wǒ bù tóngyì shìzhǎng. Wǒ juéde xiū lù zuì zhòngyào. Yīnwèi lù huài le, chē zǒu bu liǎo, shìchǎng yě méi yòng.} \\
« Habitant 2 : Je ne suis pas d'accord avec le maire. Je pense que réparer la route est le plus important. Parce que la route est abîmée, les voitures ne peuvent pas passer, le marché ne sert à rien. »

\zh{主持人：好，两位居民的看法都很有道理。还有别的看法吗？没有的话，现在我们投票！} \\
\emph{Zhǔchírén : Hǎo, liǎng wèi jūmín de kànfǎ dōu hěn yǒu dàoli. Hái yǒu bié de kànfǎ ma ? Méi yǒu de huà, xiànzài wǒmen tóupiào !} \\
« Animateur : Bien, les avis des deux habitants sont fondés. D'autres avis ? Sinon, votons maintenant ! »

\zh{主持人：举手表决。同意修路的举手……同意建市场的举手……同意修学校的举手……} \\
\emph{Zhǔchírén : Jǔshǒu biǎojué. Tóngyì xiū lù de jǔshǒu… Tóngyì jiàn shìchǎng de jǔshǒu… Tóngyì xiū xuéxiào de jǔshǒu…} \\
« Animateur : Vote à main levée. Ceux pour la route lèvent la main… Pour le marché… Pour l'école… »

\zh{主持人：修学校得票最多。我们选第三个项目：修学校！会议结束。谢谢大家！} \\
\emph{Zhǔchírén : Xiū xuéxiào dépiào zuì duō. Wǒmen xuǎn dì-sān ge xiàngmù : xiū xuéxiào ! Huìyì jiéshù. Xièxie dàjiā !} \\
« Animateur : L'école a le plus de voix. Nous choisissons le 3e projet : rénover l'école ! Séance levée. Merci à tous ! »
\end{textebox}

\begin{exemplebox}[Analyse linguistique du dialogue modèle]
\begin{itemize}
\item \textbf{Marqueurs de structure} : \zh{第一/第二/第三} (énumération), \zh{因为…所以…} (argumentation), \zh{好/没有的话} (gestion du tour de parole).
\item \textbf{Contraste d'opinions} : \zh{也} (aussi, ajout) chez le maire vs \zh{不同意} (désaccord) chez l'habitant 2.
\item \textbf{Aspect accompli} : \zh{路坏了} (la route \emph{est devenue} mauvaise), \zh{学校好了} (l'école \emph{est devenue} bonne), \zh{会议结束了} (la réunion \emph{s'est terminée}). \zh{了} marque le changement d'état.
\item \textbf{Vocabulaire de vote} : \zh{举手表决} (vote à main levée), \zh{得票最多} (obtenir le plus de voix), \zh{选第三个项目} (choisir le 3e projet).
\end{itemize}
\end{exemplebox}

\courssub{Jeu de rôle guidé : Simulation du conseil municipal (Tâche finale)}

\begin{experience}[Mise en situation : « Le grand débat de Bafoussam » — 40 min]
\textbf{Matériel :} Cartes de rôle (voir ci-dessous), affiche projetée, grille d'évaluation (une par élève public), minuteur.

\textbf{Déroulement en 3 phases :}

\noindent\textbf{Phase 1 : Préparation en îlots (15 min)}
\begin{enumerate}
\item \textbf{Distribution des rôles} (1 carte par élève, les autres = public/jury) :
\begin{itemize}
\item 1 \textbf{Maire} (\zh{市长}) : Prépare l'ouverture + présentation des 3 projets (\zh{第一…第二…第三…}) + 1 argument personnel.
\item 4 \textbf{Habitants} (\zh{居民}) : Chacun choisit UN projet différent. Rédit 2 phrases : \zh{我觉得…最重要，因为…所以…}.
\item 2 \textbf{Journalistes} (\zh{记者}) : Préparent 3 questions chacun : \zh{你有什么看法？为什么？你同意…吗？}.
\item 1 \textbf{Animateur} (\zh{主持人}) : Prépare : Ouverture (\zh{欢迎…今天我们选…}), 2 relances (\zh{你呢？还有别的看法吗？}), Conclusion + annonce vote (\zh{现在我们投票…我们选第X个项目}).
\end{itemize}
\item \textbf{Écriture du brouillon} : Chaque élève écrit ses répliques en caractères + pinyin. Le professeur circule, corrige les tons et l'ordre \zh{因为…所以…}.
\item \textbf{Répétition phonétique} : Lecture à voix basse dans le groupe, focus sur les mots-clés du rôle.
\end{enumerate}

\noindent\textbf{Phase 2 : Simulation devant la classe (20 min)}
\begin{enumerate}
\item Installation : Maire et animateur face à la classe. Habitants et journalistes assis au premier rang.
\item L'animateur lance la séance. Le maire présente. Les habitants s'expriment (tour de table). Journalistes interrogent (maire puis habitants). Animateur relance les silencieux (\zh{王同学，你有什么看法？}). Animateur appelle au vote, dépouille (main levée), annonce résultat.
\item \textbf{Contraintes obligatoires} : Chaque participant (sauf public) doit placer \textbf{au minimum} : 1 formule d'avis (\zh{我觉得/我认为}), 1 justification (\zh{因为…}), 1 relance (pour animateur/journalistes : \zh{你呢？/为什么？}). Respect des tons ciblés.
\end{enumerate}

\noindent\textbf{Phase 3 : Évaluation par les pairs \& Méta-réflexion (5 min)}
\begin{itemize}
\item Le public remplit la \textbf{Grille d'évaluation} (0=absent, 1=fragile, 2=maîtrisé) pour 3 intervenants de son choix :
\begin{center}\begin{tabularx}{\linewidth}{|X|c|c|c|}\hline
\rowcolor{popblueL} Critère & 0 & 1 & 2 \\ \hline
Tons des mots-clés (shìzhǎng, xuéxiào, juéde, yīnwèi…) & & & \\ \hline
Vocabulaire thème (min. 3 mots : 项目, 选择, 市场, 学校, 路…) & & & \\ \hline
Structures cibles (我觉得, 因为…所以…, 你有什么看法？, 最重要) & & & \\ \hline
Fluidité / Débit / Cohérence & & & \\ \hline
Contenu / Pertinence argumentaire & & & \\ \hline
\end{tabularx}\end{center}
\item Dépouillement rapide : l'animateur annonce le projet gagnant en chinois.
\item \textbf{Bilan oral collectif} : « Qu'est-ce qui était facile ? Difficile ? Quel mot de vocabulaire a manqué ? »
\end{itemize}
\end{experience}

\courssub{Production écrite d'entraînement : Compte rendu de la réunion}

\begin{exercice}[Rédaction : Compte rendu du conseil municipal]
Rédige en chinois (caractères + pinyin) un court compte rendu de la simulation (6--8 phrases). Utilise le passé accompli (\zh{了}), les connecteurs logiques et le vocabulaire de la leçon.
\begin{itemize}
\item \textbf{Amorce} : \zh{昨天我们班开了一个会议…} (\emph{Zuótiān wǒmen bān kāi le yí ge huìyì…})
\item \textbf{Contraintes} : Présenter le contexte (jumelage, 3 projets), citer au moins 2 avis différents avec \zh{因为…所以…}, décrire le vote (\zh{投票}, \zh{举手}, \zh{得票最多}), annoncer le résultat (\zh{我们选了…项目}).
\item \textbf{Objectif} : Réinvestir \zh{了} (accompli), \zh{因为…所以…}, \zh{第一/第二/第三}, \zh{觉得}, \zh{最重要}, vocabulaire thème.
\end{itemize}
\end{exercice}

\begin{corrige}[Exercice — Compte rendu (Proposition de corrigé)]
\textbf{Version élève attendue (HSK 3-4) :}

\zh{昨天我们班开了一个市政会，因为彝良市和昆明市是友好城市。} \\
\emph{Zuótiān wǒmen bān kāi le yí ge shìzhènghuì, yīnwèi Yílíng shì hé Kūnmíng shì yǒuhǎo chéngshì.} \\
« Hier, notre classe a tenu un conseil municipal, car la ville de Bafoussam et Kunming sont villes jumelées. »

\zh{市长介绍了三个项目：第一，修路；第二，建市场；第三，修学校。} \\
\emph{Shìzhǎng jièshào le sān ge xiàngmù : dì-yī, xiū lù ; dì-èr, jiàn shìchǎng ; dì-sān, xiū xuéxiào.} \\
« Le maire a présenté trois projets : 1. réparer la route ; 2. construire un marché ; 3. rénover l'école. »

\zh{居民一觉得修学校最重要，因为孩子们需要好学校，所以城市未来会好。} \\
\emph{Jūmín yī juéde xiū xuéxiào zuì zhòngyào, yīnwèi háizi men xūyào hǎo xuéxiào, suǒyǐ chéngshì wèilái huì hǎo.} \\
« L'habitant 1 pense que rénover l'école est le plus important, car les enfants ont besoin d'une bonne école, donc l'avenir de la ville sera bon. »

\zh{居民二不同意，他觉得修路最重要，因为路坏了，车走不了。} \\
\emph{Jūmín èr bù tóngyì, tā juéde xiū lù zuì zhòngyào, yīnwèi lù huài le, chē zǒu bu liǎo.} \\
« L'habitant 2 n'est pas d'accord, il pense que réparer la route est le plus important, car la route est abîmée, les voitures ne peuvent pas passer. »

\zh{最后我们投票了。修学校得票最多。我们选了第三个项目。} \\
\emph{Zuìhòu wǒmen tóupiào le. Xiū xuéxiào dépiào zuì duō. Wǒmen xuǎn le dì-sān ge xiàngmù.} \\
« Finalement nous avons voté. L'école a obtenu le plus de voix. Nous avons choisi le 3e projet. »

\zh{会议结束了，大家都很高兴。} \\
\emph{Huìyì jiéshù le, dàjiā dōu hěn gāoxìng.} \\
« La réunion s'est terminée, tout le monde était content. »
\end{corrige}

\courssub{Bilan et auto-évaluation}

\begin{aretenir}[Ce que tu dois retenir de la Leçon 20]
\begin{itemize}
\item \textbf{Lexique} : Maîtriser le vocabulaire de la gouvernance locale (\zh{市长, 居民, 市政府, 市政会, 项目, 预算, 优先事项, 投票, 友好城市}) et des infrastructures (\zh{修路, 建市场, 修学校, 修 vs 建}).
\item \textbf{Structures syntaxiques} :
\begin{itemize}
\item Énumération formelle : \zh{第一… 第二… 第三…}
\item Opinion + Justification : \zh{我觉得… 因为… 所以…} (Ordre immuable Cause \textrightarrow Conséquence).
\item Superlatif : \zh{最} + adjectif (\zh{最重要}).
\item Interrogation / Relance : \zh{你有什么看法？}, \zh{为什么？}, \zh{你同意吗？}, \zh{你呢？}.
\item Accord/Désaccord : \zh{我同意 / 我不同意 / 有道理}.
\item Passé / Changement d'état : \zh{了} final (\zh{路坏了, 会议结束了, 我们选了…}).
\end{itemize}
\item \textbf{Phonétique} : Automatiser les paires tonales \textbf{4-3} (\zh{shìzhǎng, shìchǎng, kànfǎ, jìzhě}), \textbf{2-4} (\zh{xuéxiào, tóupiào}), \textbf{1-2} (\zh{jūmín}), \textbf{Ton neutre} (\zh{juéde, shénme, ma, ne, de}). Distinguer \emph{zh} / \emph{ch} / \emph{sh} / \emph{r}.
\item \textbf{Compétence socioculturelle} : Comprendre le fonctionnement d'un conseil municipal, le concept de jumelage (\zh{友好城市}) Cameroun-Chine, et le rôle de la commune dans le développement local.
\item \textbf{Compétence citoyenne} : Savoir argumenter, écouter, voter, accepter le résultat démocratique.
\end{itemize}
\end{aretenir}

\begin{methode}[Pour progresser : Entraînement autonome]
\begin{enumerate}
\item \textbf{Flashcards Anki/Quizlet} : Crée des cartes « Caractère + Pinyin + Ton » pour les 20 mots du tableau vocabulaire. Révise 5 min/jour.
\item \textbf{Shadowing} : Réécoute le dialogue modèle 2 (audio prof ou synthèse vocale) et répète en décalé de 1 seconde (technique \emph{shadowing}) pour calquer la prosodie.
\item \textbf{Écriture guidée} : Refais l'exercice de compte rendu en changeant le projet gagnant (ex. \zh{建市场} gagne). Varie les justifications (\zh{因为商人多, 因为经济重要}).
\item \textbf{Débat inversé} : Avec un camarade, tirez un sujet au sort (ex. \zh{建公园 vs 修医院}) et débattez 3 min en chinois, minuteur en main. Notez les formules utilisées.
\end{enumerate}
\end{methode}

\newpage

\coursec{Méthodes et exercices résolus}

\coursec{Leçon 20 — Expression orale : collectivités locales}

\courssub{Lexique des collectivités locales}

\begin{definition}[Vocabulaire essentiel — Collectivités locales]
Le tableau ci-dessous regroupe le lexique de base pour décrire une commune, ses équipements et ses acteurs. Mémorisez les tons (marqués sur les voyelles) et les classificateurs.
\begin{center}
\begin{tabularx}{\linewidth}{|X|X|X|X|}
\hline
\rowcolor{popblueL} \textbf{Caractères} & \textbf{Pinyin} & \textbf{Nature} & \textbf{Traduction} \\
\hline
城市 & chéngshì & n. & ville, cité \\
\hline
市区 & shìqū & n. & centre-ville, zone urbaine \\
\hline
郊区 & jiāoqū & n. & banlieue, zone périurbaine \\
\hline
市中心 & shìzhōngxīn & n. & centre-ville (litt. cœur de ville) \\
\hline
市场 & shìchǎng & n. & marché \\
\hline
超市 & chāoshì & n. & supermarché \\
\hline
商店 & shāngdiàn & n. & magasin, boutique \\
\hline
政府 & zhèngfǔ & n. & gouvernement, administration locale \\
\hline
市长 & shìzhǎng & n. & maire \\
\hline
居民 & jūmín & n. & habitant, résident \\
\hline
商人 & shāngrén & n. & commerçant, marchand \\
\hline
社区 & shèqū & n. & quartier, communauté de quartier \\
\hline
街道 & jiēdào & n. & rue, sous-district administratif \\
\hline
公园 & gōngyuán & n. & parc public \\
\hline
医院 & yīyuàn & n. & hôpital \\
\hline
学校 & xuéxiào & n. & école \\
\hline
干净 & gānjìng & adj. & propre \\
\hline
热闹 & rènao & adj. & animé, vivant \\
\hline
拥挤 & yōngjǐ & adj. & bondé, encombré \\
\hline
方便 & fāngbiàn & adj. & pratique, commode \\
\hline
建 & jiàn & v. & construire, bâtir \\
\hline
搬 & bān & v. & déménager, déplacer \\
\hline
同意 & tóngyì & v. & être d'accord \\
\hline
反对 & fǎnduì & v. & s'opposer, être contre \\
\hline
讨论 & tǎolùn & v. & discuter, débattre \\
\hline
决定 & juédìng & v. & décider \\
\hline
因为 & yīnwèi & conj. & parce que \\
\hline
所以 & suǒyǐ & conj. & donc, c'est pourquoi \\
\hline
可是 & kěshì & conj. & mais, cependant \\
\hline
觉得 & juéde & v. & trouver (avis), penser \\
\hline
认为 & rènwéi & v. & estimer, considérer (plus formel) \\
\hline
\end{tabularx}
\end{center}
\textbf{Classificateurs clés :} \zh{个} (gè, générique : \zh{一个市场}), \zh{所} (suǒ, écoles/hôpitaux : \zh{一所学校}, \zh{一所医院}), \zh{条} (tiáo, routes/rues : \zh{一条路}), \zh{座} (zuò, montagnes/bâtiments imposants : \zh{一座桥}), \zh{位} (wèi, honorifique pour personnes : \zh{一位市长}).
\end{definition}

\courssub{Dialogue modèle 1 : présenter sa commune}

\begin{textebox}[Dialogue : Présentation de Yaoundé]
\textbf{Contexte :} Un élève présente sa ville lors d'un échange scolaire.
\zh{大家好！我叫保罗，今天我要介绍雅温得。} \\
\emph{Dàjiā hǎo! Wǒ jiào Bǎoluó, jīntiān wǒ yào jièshào Yǎwēndé.} \\
« Bonjour à tous ! Je m'appelle Paul, aujourd'hui je vais présenter Yaoundé. »

\zh{雅温得在喀麦隆的中部，是首都。} \\
\emph{Yǎwēndé zài Kāmàilóng de zhōngbù, shì shǒudū.} \\
« Yaoundé se trouve au centre du Cameroun, c'est la capitale. »

\zh{我们市有一所很大的大学，还有很多公园。} \\
\emph{Wǒmen shì yǒu yì suǒ hěn dà de dàxué, hái yǒu hěn duō gōngyuán.} \\
« Notre ville a une très grande université et aussi beaucoup de parcs. »

\zh{市中心很热闹，商店很多，也很方便。} \\
\emph{Shìzhōngxīn hěn rènao, shāngdiàn hěn duō, yě hěn fāngbiàn.} \\
« Le centre-ville est très animé, il y a beaucoup de magasins, c'est aussi très pratique. »

\zh{我觉得雅温得很漂亮，气候也不错。} \\
\emph{Wǒ juéde Yǎwēndé hěn piàoliang, qìhòu yě bùcuò.} \\
« Je trouve Yaoundé très jolie, le climat est aussi pas mal. »

\zh{我的介绍完了，谢谢大家！} \\
\emph{Wǒ de jièshào wán le, xièxie dàjiā!} \\
« Ma présentation est terminée, merci à tous ! »
\end{textebox}

\courssub{Savoir-faire 1 : Présenter oralement un exposé simple en chinois}

\begin{methode}[Structurer un exposé sur sa commune]
Pour présenter sa commune ou son quartier en chinois, suivez ce plan en 5 étapes. Annoncez le plan avec \zh{第一…… 第二…… 第三……} (\emph{dì-yī... dì-èr... dì-sān...}).
\begin{enumerate}
\item \textbf{Saluer et annoncer le sujet} : \zh{大家好！我叫……，今天我要介绍……} (\emph{Dàjiā hǎo! Wǒ jiào..., jīntiān wǒ yào jièshào...})
\item \textbf{Situer géographiquement} : Structure \textbf{Sujet + 在 + Lieu de référence + 方位词}. Exemple : \zh{杜阿拉在喀麦隆的西部。} (\emph{Dù'ālā zài Kāmàilóng de xībù.}) « Douala est à l'ouest du Cameroun. » \textbf{Attention :} le complément de lieu (\zh{在...}) se place \emph{avant} le verbe d'état (\zh{是}) ou d'action.
\item \textbf{Décrire (existence et qualités)} : Utiliser \zh{有} (\emph{yǒu}) pour l'existence + classificateur. \zh{我们市有一个大市场。} (\emph{Wǒmen shì yǒu yí ge dà shìchǎng.}) Utiliser \zh{很/真} + adjectif pour la qualité. \zh{环境很干净。} (\emph{Huánjìng hěn gānjìng.})
\item \textbf{Exprimer un avis personnel} : \zh{我觉得……} (\emph{Wǒ juéde...}) ou \zh{我认为……} (\emph{Wǒ rènwéi...}). Exemple : \zh{我觉得这里很适合居住。} (\emph{Wǒ juéde zhèlǐ hěn shìhé jūzhù.})
\item \textbf{Conclure et remercier} : \zh{我的介绍完了，谢谢大家！} (\emph{Wǒ de jièshào wán le, xièxie dàjiā!})
\end{enumerate}
\textbf{Réflexes oraux :} Parler lentement, articuler les tons, regarder l'auditoire, ne pas lire ses notes.
\end{methode}

\begin{exoresolu}[Exposé : présenter sa ville (Douala)]
\textbf{Énoncé.} Présente ta ville (Douala) en 5-6 phrases en suivant le plan : salutation, situation, description, avis, conclusion. Rédige le texte complet avec pinyin et traduction.
\tcblower
\textbf{Raisonnement.} Application stricte de la méthode en 5 étapes. Respect de l'ordre des mots : Sujet + (在 + Lieu) + Verbe. Utilisation des classificateurs (\zh{个} pour \zh{港口}, \zh{所} pour \zh{大学} si mentionné).

\textbf{Production orale rédigée.}

\zh{大家好！我叫玛丽，今天我要介绍杜阿拉。} \\
\emph{Dàjiā hǎo! Wǒ jiào Mǎlì, jīntiān wǒ yào jièshào Dù'ālā.} \\
« Bonjour à tous ! Je m'appelle Marie, aujourd'hui je vais présenter Douala. »

\zh{杜阿拉在喀麦隆的西部，是一个很大的城市。} \\
\emph{Dù'ālā zài Kāmàilóng de xībù, shì yí ge hěn dà de chéngshì.} \\
« Douala se trouve à l'ouest du Cameroun, c'est une très grande ville. »

\zh{我们市有一个很大的港口，还有一所著名的大学。} \\
\emph{Wǒmen shì yǒu yí ge hěn dà de gǎngkǒu, hái yǒu yì suǒ zhùmíng de dàxué.} \\
« Notre ville a un très grand port et aussi une université célèbre. »

\zh{市中心很热闹，市场也很多，很方便。} \\
\emph{Shìzhōngxīn hěn rènao, shìchǎng yě hěn duō, hěn fāngbiàn.} \\
« Le centre-ville est très animé, il y a aussi beaucoup de marchés, c'est très pratique. »

\zh{我觉得杜阿拉很有活力，但交通有点拥挤。} \\
\emph{Wǒ juéde Dù'ālā hěn yǒu huólì, dàn jiāotōng yǒudiǎn yōngjǐ.} \\
« Je trouve Douala très dynamique, mais la circulation est un peu encombrée. »

\zh{我的介绍完了，谢谢大家！} \\
\emph{Wǒ de jièshào wán le, xièxie dàjiā!} \\
« Ma présentation est terminée, merci à tous ! »

\textbf{Conclusion.} L'exposé est structuré, les tons sont notés, le vocabulaire du thème est utilisé (\zh{港口}, \zh{热闹}, \zh{拥挤}), la structure \zh{在...西部} est correcte.
\end{exoresolu}

\courssub{Formules fonctionnelles : opiner, réagir, relancer}

\begin{propriete}[Boîte à outils du débat oral]
Ces formules sont à connaître par cœur pour l'interaction orale (Bac). Elles structurent l'échange.
\begin{center}
\begin{tabularx}{\linewidth}{|X|X|X|}
\hline
\rowcolor{popblueL} \textbf{Fonction} & \textbf{Formule chinoise (Pinyin)} & \textbf{Équivalent français} \\
\hline
\multicolumn{3}{|c|}{\textbf{Donner son avis}} \\
\hline
Avis personnel & \zh{我觉得……} (\emph{Wǒ juéde...}) & Je trouve que... / J'ai l'impression que... \\
\hline
Opinion réfléchie & \zh{我认为……} (\emph{Wǒ rènwéi...}) & Je pense que... / J'estime que... \\
\hline
Point de vue subjectif & \zh{对我来说，……} (\emph{Duì wǒ lái shuō,...}) & Pour moi... / À mon avis... \\
\hline
\multicolumn{3}{|c|}{\textbf{Être d'accord}} \\
\hline
Accord total & \zh{我同意你的看法。} (\emph{Wǒ tóngyì nǐ de kànfǎ.}) & Je suis d'accord avec ton point de vue. \\
\hline
Accord partagé & \zh{对，我也这么想。} (\emph{Duì, wǒ yě zhème xiǎng.}) & Oui, je pense pareil. \\
\hline
Appui & \zh{你说得很对。} (\emph{Nǐ shuō de hěn duì.}) & Tu as tout à fait raison. \\
\hline
\multicolumn{3}{|c|}{\textbf{Ne pas être d'accord (poli)}} \\
\hline
Désaccord nuancé & \zh{我不太同意，因为……} (\emph{Wǒ bú tài tóngyì, yīnwèi...}) & Je ne suis pas tout à fait d'accord, car... \\
\hline
Objection & \zh{可是…… / 但是……} (\emph{Kěshì... / Dànshì...}) & Mais... / Cependant... \\
\hline
Autre avis & \zh{我觉得不一定，……} (\emph{Wǒ juéde bù yídìng,...}) & Je ne pense pas nécessairement que... \\
\hline
\multicolumn{3}{|c|}{\textbf{Relancer / Interroger}} \\
\hline
Demander avis & \zh{你怎么看？} (\emph{Nǐ zěnme kàn?}) & Qu'en penses-tu ? / Comment tu vois ça ? \\
\hline
Rebondir & \zh{你觉得呢？} (\emph{Nǐ juéde ne?}) & Et toi, qu'en penses-tu ? \\
\hline
Demander raison & \zh{为什么？} (\emph{Wèishénme?}) & Pourquoi ? \\
\hline
Préciser & \zh{具体说说看。} (\emph{Jùtǐ shuō shuo kàn.}) & Dis-le plus concrètement. \\
\hline
\end{tabularx}
\end{center}
\textbf{Règle d'or :} Toujours justifier avec \zh{因为……所以……} (\emph{yīnwèi... suǒyǐ...}). Adoucir le refus avec \zh{不太} (\emph{bú tài}) plutôt que \zh{不} (\emph{bù}) seul.
\end{propriete}

\courssub{Dialogue modèle 2 : débat sur un projet municipal}

\begin{textebox}[Dialogue : Réunion municipale — Nouveau marché à Bafoussam]
\textbf{Contexte :} Le maire propose un nouveau marché. Deux habitants réagissent.
\zh{市长：大家好！我们要在巴福萨建一个新的大市场。大家怎么看？} \\
\emph{Shìzhǎng: Dàjiā hǎo! Wǒmen yào zài Bāfúsà jiàn yí ge xīn de dà shìchǎng. Dàjiā zěnme kàn?} \\
« Maire : Bonjour ! Nous allons construire un grand nouveau marché à Bafoussam. Qu'en pensez-vous ? »

\zh{居民甲：我觉得这是好主意，因为现在的市场太小、太脏了。} \\
\emph{Jūmín jiǎ: Wǒ juéde zhè shì hǎo zhǔyi, yīnwèi xiànzài de shìchǎng tài xiǎo, tài zāng le.} \\
« Habitant 1 : Je trouve que c'est une bonne idée, car le marché actuel est trop petit et trop sale. »

\zh{居民乙：我不太同意。新市场离我们村太远了，不方便。} \\
\emph{Jūmín yǐ: Wǒ bú tài tóngyì. Xīn shìchǎng lí wǒmen cūn tài yuǎn le, bù fāngbiàn.} \\
« Habitant 2 : Je ne suis pas tout à fait d'accord. Le nouveau marché est trop loin de notre village, ce n'est pas pratique. »

\zh{市长：好，谢谢大家的意见。我们会再考虑一下位置。} \\
\emph{Shìzhǎng: Hǎo, xièxie dàjiā de yìjiàn. Wǒmen huì zài kǎolǜ yí xià wèizhi.} \\
« Maire : Bien, merci pour vos avis. Nous allons reconsidérer l'emplacement. »
\end{textebox}

\courssub{Savoir-faire 2 : Échanger dans un petit débat}

\begin{methode}[Conduire un échange argumenté]
\begin{enumerate}
\item \textbf{Écouter} la proposition ou l'avis du camarade (identifier les mots-clés : \zh{建}, \zh{市场}, \zh{学校}, \zh{路}).
\item \textbf{Choisir sa position} : d'accord (\zh{同意}) ou pas d'accord (\zh{不同意/不太同意}).
\item \textbf{Construire sa réplique} : Formule d'avis + \zh{因为} + Argument + (éventuel \zh{可是} + nuance).
\item \textbf{Relancer} le camarade : \zh{你怎么看？} ou \zh{为什么？} pour maintenir l'interaction.
\item \textbf{Synthétiser} si vous menez le débat : \zh{总之……} (\emph{Zǒngzhī...}) « En résumé... ».
\end{enumerate}
\textbf{Critères d'évaluation (Bac) :} Pertinence du vocabulaire, justesse des tons, cohérence syntaxique (ordre des mots), politesse (\zh{不太}), interaction réelle (relance).
\end{methode}

\begin{exoresolu}[Débat guidé : Projet d'une nouvelle route]
\textbf{Énoncé.} Complétez le dialogue ci-dessous. A propose de goudronner la rue principale. B est favorable (accès hôpital). A émet une réserve (bruit, poussière). Utilisez les formules de la boîte à outils.
\begin{verbatim}
A : 我们要修路。你怎么看？
B : \_\_\_\_\_\_\_\_\_\_\_\_ (Avis + parce que)
A : \_\_\_\_\_\_\_\_\_\_\_\_ (Désaccord poli + argument)
B : \_\_\_\_\_\_\_\_\_\_\_\_ (Réaction + relance)
\end{verbatim}
\tcblower
\textbf{Proposition de complétion (plusieurs réponses possibles).}

A : \zh{我们要修路。你怎么看？} (\emph{Wǒmen yào xiū lù. Nǐ zěnme kàn?}) « Nous allons refaire la route. Qu'en penses-tu ? »

B : \zh{我同意，因为现在的路太坏了，去医院很难。} (\emph{Wǒ tóngyì, yīnwèi xiànzài de lù tài huài le, qù yīyuàn hěn nán.}) « Je suis d'accord, car la route actuelle est trop abîmée, aller à l'hôpital est difficile. »

A : \zh{我不太同意，因为施工会很吵，灰尘也很大。} (\emph{Wǒ bú tài tóngyì, yīnwèi shīgōng huì hěn chǎo, huīchén yě hěn dà.}) « Je ne suis pas tout à fait d'accord, car les travaux seront très bruyants et il y aura beaucoup de poussière. »

B : \zh{可是路修好以后就方便了。你觉得呢？} (\emph{Kěshì lù xiū hǎo yǐhòu jiù fāngbiàn le. Nǐ juéde ne?}) « Mais une fois la route refaite, ce sera pratique. Et toi, qu'en penses-tu ? »

\textbf{Analyse.} \zh{修路} (xiū lù) = construire/réparer une route. \zh{坏} (huài) = abîmé. \zh{施工} (shīgōng) = travaux. \zh{吵} (chǎo) = bruyant. \zh{灰尘} (huīchén) = poussière. \zh{以后} (yǐhòu) = après/après que. Structure \zh{因为……所以……} implicite chez B (phrase 1), explicite chez A (phrase 2).
\end{exoresolu}

\courssub{Travail phonétique : les tons du thème}

\begin{methode}[Maîtriser les tons du lexique « collectivités »]
\begin{enumerate}
\item \textbf{Mémoriser le ton de chaque caractère isolé} (dictionnaire) :
\begin{itemize}
\item \zh{城} (chéng, 2) \zh{市} (shì, 4) $\rightarrow$ \zh{城市} (\emph{chéngshì})
\item \zh{市} (shì, 4) \zh{场} (chǎng, 3) $\rightarrow$ \zh{市场} (\emph{shìchǎng})
\item \zh{干} (gān, 1) \zh{净} (jìng, 4) $\rightarrow$ \zh{干净} (\emph{gānjìng})
\item \zh{同} (tóng, 2) \zh{意} (yì, 4) $\rightarrow$ \zh{同意} (\emph{tóngyì})
\item \zh{买} (mǎi, 3) vs \zh{卖} (mài, 4) — \textbf{paire minimale cruciale}.
\end{itemize}
\item \textbf{Règle des deux 3èmes tons (Sandhi)} : Quand deux syllabes de 3ème ton se suivent, la \textbf{première} se prononce comme un 2ème ton (montant).
\begin{itemize}
\item \zh{你好} : \emph{nǐ} + \emph{hǎo} $\rightarrow$ \textbf{\emph{ní hǎo}}.
\item \zh{很好} : \emph{hěn} + \emph{hǎo} $\rightarrow$ \textbf{\emph{hén hǎo}}.
\item \zh{马上} : \emph{mǎ} + \emph{shàng} $\rightarrow$ \textbf{\emph{má shàng}}.
\end{itemize}
\textbf{Attention :} \zh{市长} (\emph{shìzhǎng}) n'est \textbf{PAS} concerné car \zh{市} est 4ème ton (\emph{shì}). On prononce bien \emph{shìzhǎng}.
\item \textbf{Règle du ton de \zh{不} (bù)} :
\begin{itemize}
\item Devant ton 1, 2, 3 : reste 4ème ton (\emph{bù}). Ex: \zh{不去} (\emph{bù qù}), \zh{不来} (\emph{bù lái}), \zh{不好} (\emph{bù hǎo}).
\item Devant ton 4 : passe au 2ème ton (\emph{bú}). Ex: \zh{不是} (\emph{bú shì}), \zh{不太} (\emph{bú tài}), \zh{不要} (\emph{bú yào}).
\end{itemize}
\item \textbf{Entraînement quotidien} : Lire le dialogue modèle 1 à voix haute 3 fois/jour en exagérant les contours tonaux. S'enregistrer.
\end{enumerate}
\end{methode}

\begin{exoresolu}[Phonétique : identification et correction]
\textbf{Énoncé.} 1) Donnez la prononciation pinyin (tons sur voyelles) de : \zh{城市}, \zh{市场}, \zh{干净}, \zh{同意}, \zh{不太同意}. 2) Un élève lit \zh{市场} \emph{shīchǎng} et \zh{买菜} \emph{wǒ mài cài}. Corrigez.
\tcblower
\textbf{1) Prononciations exactes.}
\begin{itemize}
\item \zh{城市} : \emph{chéngshì} (2 + 4)
\item \zh{市场} : \emph{shìchǎng} (4 + 3)
\item \zh{干净} : \emph{gānjìng} (1 + 4)
\item \zh{同意} : \emph{tóngyì} (2 + 4)
\item \zh{不太同意} : \emph{bú tài tóngyì} (\zh{不} $\rightarrow$ 2ème ton devant \zh{太} 4ème ton).
\end{itemize}

\textbf{2) Corrections.}
\begin{itemize}
\item \zh{市场} : L'élève dit \emph{shīchǎng} (1er ton sur \zh{市}). \textbf{Erreur.} \zh{市} se prononce \textbf{4ème ton} (\emph{shì}) dans ce mot. Correct : \emph{shìchǎng}.
\item \zh{买菜} : L'élève dit \emph{wǒ mài cài} (4ème ton sur \zh{买}). \textbf{Erreur de sens.} \zh{卖} (\emph{mài}, 4ème ton) = vendre. \zh{买} (\emph{mǎi}, 3ème ton) = acheter. Phrase correcte : \zh{我买菜。} (\emph{Wǒ mǎi cài.}) « J'achète des légumes. »
\end{itemize}

\textbf{Conclusion.} Une erreur de ton sur \zh{市} (shì vs shī) ou \zh{买/卖} change le sens ou rend le mot incompréhensible. Au Bac, la justesse tonale du lexique thématique est notée sur 5 points (prononciation).
\end{exoresolu}

\courssub{Jeux de rôle et débat guidé}

\begin{experience}[Activité : Conseil municipal — Projet « École à Buea »]
\textbf{Objectif :} Produire un échange oral interactif de 2-3 minutes en groupes de 4.
\textbf{Rôles :} 1 Maire (\zh{市长}), 2 Habitants (\zh{居民} : un favorable, un défavorable), 1 Observateur (note le vocabulaire, les tons, les formules de relance).
\textbf{Scénario :} Le maire annonce la construction d'une nouvelle école primaire à Buea (\zh{布埃亚} Bù'āiyà). Débat : emplacement, budget, nombre de classes.
\textbf{Contraintes linguistiques :}
\begin{itemize}
\item Maire : Utiliser \zh{我们要建设……} / \zh{请大家提意见。}
\item Habitants : Utiliser \zh{我觉得……因为……} / \zh{我不太同意……可是……} / \zh{你怎么看？}
\item Observateur : Remplit une grille (Vocabulaire /5, Tons /5, Interaction /5, Structure /5).
\end{itemize}
\textbf{Déroulement :} 5 min préparation (écriture des répliques clés avec pinyin) $\rightarrow$ 3 min jeu $\rightarrow$ 2 min feedback observateur $\rightarrow$ Rotation des rôles.
\end{experience}

\begin{exoresolu}[Jeu de rôle écrit : Modèle pour répétition]
\textbf{Énoncé.} Rédigez le script complet du jeu de rôle « École à Buea » (4 répliques) avec pinyin et traduction, en réinvestissant les structures vues.
\tcblower
\textbf{Script du jeu de rôle.}

\textbf{Maire (\zh{市长}) :} \zh{各位居民好！我们计划在布埃亚建一所新小学。请大家提意见。} \\
\emph{Gèwèi jūmín hǎo! Wǒmen jìhuà zài Bù'āiyà jiàn yì suǒ xīn xiǎoxué. Qǐng dàjiā tí yìjiàn.} \\
« Bonjour chers habitants ! Nous prévoyons de construire une nouvelle école primaire à Buea. Merci de donner vos avis. »

\textbf{Habitant 1 (\zh{居民甲} — favorable) :} \zh{我完全同意，因为我们社区的孩子很多，现在的学校太远了。} \\
\emph{Wǒ wánquán tóngyì, yīnwèi wǒmen shèqū de háizi hěn duō, xiànzài de xuéxiào tài yuǎn le.} \\
« Je suis tout à fait d'accord, car il y a beaucoup d'enfants dans notre quartier, l'école actuelle est trop loin. »

\textbf{Habitant 2 (\zh{居民乙} — réservé) :} \zh{我不太同意选址，因为那个地方太吵，不适合读书。你们考虑过别的地方吗？} \\
\emph{Wǒ bú tài tóngyì xuǎnzhǐ, yīnwèi nàge dìfang tài chǎo, bù shìhé dúshū. Nǐmen kǎolǜ guò bié de dìfang ma?} \\
« Je ne suis pas tout à fait d'accord sur l'emplacement, car cet endroit est trop bruyant, pas adapté aux études. Avez-vous envisagé d'autres lieux ? »

\textbf{Maire (\zh{市长}) — conclusion :} \zh{谢谢您的建议。我们会认真考虑选址问题，争取方便大家。散会！} \\
\emph{Xièxie nín de jiànyì. Wǒmen huì rènzhēn kǎolǜ xuǎnzhǐ wèntí, zhēngqǔ fāngbiàn dàjiā. Sànhuì!} \\
« Merci pour votre suggestion. Nous étudierons sérieusement la question de l'emplacement pour arranger tout le monde. La séance est levée ! »

\textbf{Points grammaticaux clés.} \zh{计划} (jìhuà) verbe/nom « planifier/plan ». Classificateur \zh{所} (suǒ) pour \zh{学校}. \zh{选址} (xuǎnzhǐ) « choix de l'emplacement ». \zh{争取} (zhēngqǔ) « s'efforcer d'obtenir ». \zh{散会} (sànhuì) « lever la séance » (formel).
\end{exoresolu}

\begin{attention}[Les 5 erreurs fatales à l'oral (Bac) — Check-list de révision]
\begin{enumerate}
\item \textbf{Tons instables sur le lexique thématique} : \zh{市场} lu \emph{shīchǎng} (1er ton) au lieu de \emph{shìchǎng} (4+3) ; \zh{买} (\emph{mǎi}) lu \emph{mài} (vendre). \textbf{Action :} Répéter les paires minimales \zh{买/卖}, \zh{市/是} (\emph{shì}) chaque jour.
\item \textbf{Confusion de caractères proches} : \zh{市} (ville/marché) vs \zh{币} (bì, monnaie) ; \zh{同} (tóng, même) dans \zh{同意} vs \zh{问} (wèn, demander) ; \zh{干} (gān/gàn) vs \zh{千} (qiān). \textbf{Action :} Écrire 5 fois chaque caractère en respectant l'ordre des traits (clé \zh{巾} pour \zh{市}, \zh{一} pour \zh{干}).
\item \textbf{Ordre des mots « à la française » (Lieu après verbe)} : Dire *\zh{我住在杜阿拉} (Structure Verbe + 在 + Lieu \textbf{correcte pour 住}) mais *\zh{我在杜阿拉住} (Structure Sujet + 在 + Lieu + Verbe \textbf{possible mais marqué « Thématisation du lieu »}). \textbf{Règle pédagogique :} Pour \zh{住}, \zh{工作}, \zh{学习} $\rightarrow$ \textbf{Sujet + Verbe + 在 + Lieu} (\zh{我住在雅温得}). Pour verbes d'action/déplacement (\zh{去}, \zh{来}, \zh{建}, \zh{开会}) $\rightarrow$ \textbf{Sujet + 在 + Lieu + Verbe} (\zh{我们在雅温得开会}). Ne pas mélanger.
\item \textbf{Oubli systématique des classificateurs} : *\zh{一学校}, *\zh{两市场}, *\zh{三医院}. \textbf{Obligatoire :} \zh{一所学校}, \zh{两个市场}, \zh{三所医院}. Même pour « un » (\zh{一}) : \zh{一个市长}, \zh{一位市长} (honorifique).
\item \textbf{Désaccord brutal / Absence de justification} : Dire \zh{我不同意！} point final. \textbf{Attendu :} \zh{我不太同意，因为……} + argument. L'interaction notée au Bac exige la \textbf{relance} (\zh{你怎么看？}). Un monologue sans relance = note plafonnée.
\end{enumerate}
\end{attention}

\newpage

\coursec{Exercices}

\courssub{Je vérifie mes connaissances}

\begin{exercice}[QCM : le vocabulaire des collectivités locales]
Pour chaque question, entoure la bonne réponse.

\begin{enumerate}
\item \zh{市长} (shìzhǎng) signifie :
\begin{enumerate}
\item le chef de quartier \item le maire \item le ministre
\end{enumerate}
\item Pour dire « construire une route », on dit :
\begin{enumerate}
\item \zh{修路} (xiū lù) \item \zh{买路} (mǎi lù) \item \zh{看路} (kàn lù)
\end{enumerate}
\item \zh{居民} (jūmín) désigne :
\begin{enumerate}
\item les visiteurs \item les habitants d'une commune \item les fonctionnaires
\end{enumerate}
\item « Notre commune a un nouveau marché » se traduit par :
\begin{enumerate}
\item \zh{我们的城市有一个新市场。}
\item \zh{我们的区有一个新市场。}
\item \zh{我们的国家有一个新市场。}
\end{enumerate}
\end{enumerate}
\end{exercice}

\begin{exercice}[Vrai ou faux ? Justifie ta réponse]
Dis si les affirmations suivantes sont vraies ou fausses, puis justifie en une phrase en français.

\begin{enumerate}
\item \zh{应该} (yīnggāi) exprime une obligation ou un conseil (« devoir, il faudrait »).
\item Dans un débat, pour donner un premier argument puis un deuxième, on peut utiliser \zh{因为……而且……}.
\item \zh{还是} (háishi) s'emploie dans les questions alternatives, comme « ou » dans « tu préfères le thé ou le café ? ».
\item Pour relancer un camarade dans un débat, on peut dire \zh{你呢？} (Nǐ ne ?).
\end{enumerate}
\end{exercice}

\begin{exercice}[Complète le vocabulaire]
Complète le tableau suivant avec le caractère, le pinyin ou la traduction manquante.

\begin{center}
\begin{tabularx}{\linewidth}{|X|X|X|X|}\hline
\rowcolor{popblueL} Caractères & Pinyin & Nature & Traduction \\ \hline
社区 & \ldots\ldots\ldots & nom & quartier, communauté \\ \hline
\ldots\ldots\ldots & yīyuàn & nom & hôpital \\ \hline
建设 & jiànshè & \ldots\ldots\ldots & construire, développer \\ \hline
\ldots\ldots\ldots & \ldots\ldots\ldots & nom & projet \\ \hline
意见 & yìjiàn & nom & \ldots\ldots\ldots \\ \hline
\end{tabularx}
\end{center}
\end{exercice}

\begin{exercice}[Associe le pinyin aux caractères]
Relie chaque mot en caractères à sa transcription pinyin correcte. Attention aux tons !

\begin{center}
\begin{tabularx}{\linewidth}{|X|X|}\hline
\rowcolor{popblueL} Caractères & Pinyin à associer \\ \hline
1. 市长 & a. shìchǎng \\ \hline
2. 市场 & b. shìzhǎng \\ \hline
3. 所以 & c. kěyǐ \\ \hline
4. 可以 & d. suǒyǐ \\ \hline
5. 政府 & e. gōngyuán \\ \hline
6. 公园 & f. zhèngfǔ \\ \hline
\end{tabularx}
\end{center}
\end{exercice}

\courssub{Je m'entraîne}

\begin{exercice}[Traduis en chinois]
Traduis les phrases suivantes en chinois (caractères), avec le pinyin.

\begin{enumerate}
\item Je voudrais présenter ma commune. (utiliser \zh{介绍})
\item Il y a beaucoup d'habitants dans notre quartier.
\item Je pense qu'il faut d'abord construire un hôpital. (utiliser \zh{觉得} et \zh{应该})
\item Parce que les gens sont souvent malades, et en plus l'hôpital est très loin. (utiliser \zh{因为……而且……})
\item Es-tu d'accord avec moi ? (utiliser \zh{同意})
\end{enumerate}
\end{exercice}

\begin{exercice}[Questions alternatives avec 还是]
Transforme les cinq phrases affirmatives suivantes en questions alternatives avec \zh{还是}, en ajoutant l'option indiquée entre parenthèses.

\begin{enumerate}
\item \zh{我们应该修路。} (construire un parc)
\item \zh{市长想建一个新市场。} (un nouvel hôpital)
\item \zh{我们应该先建学校。} (d'abord construire un centre de santé)
\item \zh{这个项目对年轻人有好处。} (pour les personnes âgées)
\item \zh{我想谈谈交通问题。} (le problème de l'eau potable)
\end{enumerate}
\end{exercice}

\begin{exercice}[Texte à trous]
Complète le texte suivant avec les mots de la liste : \zh{居民、项目、应该、因为、意见、市场}.

\begin{textebox}[Un projet pour notre commune]
我们的区有一个新\ldots\ldots\ldots ：建一个大市场。\\
Wǒmen de qū yǒu yí ge xīn \ldots\ldots\ldots : jiàn yí ge dà shìchǎng.\\
\ldots\ldots\ldots 旧的\ldots\ldots\ldots 太小，所以很多\ldots\ldots\ldots 不满意。\\
\ldots\ldots\ldots jiù de \ldots\ldots\ldots tài xiǎo, suǒyǐ hěn duō \ldots\ldots\ldots bù mǎnyì.\\
我觉得我们\ldots\ldots\ldots 听听大家的\ldots\ldots\ldots 。\\
Wǒ juéde wǒmen \ldots\ldots\ldots tīngting dàjiā de \ldots\ldots\ldots .
\end{textebox}
\end{exercice}

\begin{exercice}[Remets le dialogue dans l'ordre]
Voici les répliques d'un débat sur le « nouveau marché ». Numérote-les de 1 à 6 pour reconstituer le dialogue, puis joue-le avec un camarade sans regarder le support.

\begin{enumerate}
\item[\textbullet] \zh{我不同意。因为病人很多，而且医院太远了。}\\
(Wǒ bù tóngyì. Yīnwèi bìngrén hěn duō, érqiě yīyuàn tài yuǎn le.)
\item[\textbullet] \zh{我觉得应该先建一个新市场。}\\
(Wǒ juéde yīnggāi xiān jiàn yí ge xīn shìchǎng.)
\item[\textbullet] \zh{你说得有道理，但是市场也很重要。}\\
(Nǐ shuō de yǒu dàolǐ, dànshì shìchǎng yě hěn zhòngyào.)
\item[\textbullet] \zh{为什么呢？}\\
(Wèishénme ne ?)
\item[\textbullet] \zh{因为旧市场太小，而且不干净。你呢？}\\
(Yīnwèi jiù shìchǎng tài xiǎo, érqiě bù gānjìng. Nǐ ne ?)
\item[\textbullet] \zh{好吧，我们可以先建医院，再建市场。}\\
(Hǎo ba, wǒmen kěyǐ xiān jiàn yīyuàn, zài jiàn shìchǎng.)
\end{enumerate}
\end{exercice}

\courssub{Je me prépare au Bac}

\begin{exercice}[Compréhension de l'écrit et questions de langue]
Lis le texte suivant, puis réponds aux questions.

\begin{textebox}[Notre commune change]
我住在雅温得的一个区。以前，我们区没有大市场，居民们要坐一个小时的车去市中心买东西，很不方便。\\
Wǒ zhù zài Yǎwēndé de yí ge qū. Yǐqián, wǒmen qū méiyǒu dà shìchǎng, jūmínmen yào zuò yí ge xiǎoshí de chē qù shì zhōngxīn mǎi dōngxi, hěn bù fāngbiàn.\\
去年，市长和区政府决定建一个新市场，还要修两条路。大家都很高兴，因为这个项目对商人和居民都有好处。\\
Qùnián, shìzhǎng hé qū zhèngfǔ juédìng jiàn yí ge xīn shìchǎng, hái yào xiū liǎng tiáo lù. Dàjiā dōu hěn gāoxìng, yīnwèi zhè ge xiàngmù duì shāngrén hé jūmín dōu yǒu hǎochù.\\
但是，也有人觉得应该先建医院。他们说：「市场可以等，病人的健康不能等。」\\
Dànshì, yě yǒu rén juéde yīnggāi xiān jiàn yīyuàn. Tāmen shuō : « Shìchǎng kěyǐ děng, bìngrén de jiànkāng bù néng děng. »\\
下个月，区政府要开一个会，听听居民们的意见。\\
Xià ge yuè, qū zhèngfǔ yào kāi yí ge huì, tīngting jūmínmen de yìjiàn.
\end{textebox}

\textbf{Questions de compréhension} (réponds en français) :
\begin{enumerate}
\item Pourquoi les habitants du quartier n'étaient-ils pas satisfaits avant ?
\item Quelles sont les deux réalisations décidées par le maire et le gouvernement de l'arrondissement ?
\item Quel argument avancent les partisans de la construction prioritaire d'un hôpital ? Cite la phrase du texte en chinois.
\item Que va faire le gouvernement de l'arrondissement le mois prochain ?
\end{enumerate}

\textbf{Questions de langue} :
\begin{enumerate}
\item Relève dans le texte deux phrases contenant \zh{因为} et traduis-les en français.
\item Trouve dans le texte un synonyme de \zh{高兴}.
\item Conjugue la structure : transforme \zh{居民们要坐车去市中心} en phrase au passé avec \zh{了}.
\end{enumerate}
\end{exercice}

\begin{exercice}[Thème et version]
\textbf{Thème.} Traduis en chinois (caractères et pinyin) :
\begin{enumerate}
\item Notre commune a beaucoup de projets : un marché, une école et un hôpital.
\item Je pense qu'il faut d'abord réparer les routes, parce qu'elles sont en mauvais état.
\item Le maire veut écouter l'avis des habitants.
\end{enumerate}

\textbf{Version.} Traduis en français le texte suivant :

\begin{textebox}[Le débat du quartier]
昨天，我们区开了一个会。有的居民说应该先建市场，有的居民说应该先建学校。市长说：「大家的意见都很重要，我们一起想办法。」\\
Zuótiān, wǒmen qū kāi le yí ge huì. Yǒu de jūmín shuō yīnggāi xiān jiàn shìchǎng, yǒu de jūmín shuō yīnggāi xiān jiàn xuéxiào. Shìzhǎng shuō : « Dàjiā de yìjiàn dōu hěn zhòngyào, wǒmen yìqǐ xiǎng bànfǎ. »
\end{textebox}
\end{exercice}

\begin{exercice}[Expression écrite et orale type Bac]
\textbf{Sujet d'expression écrite.} Tu es délégué(e) de ta classe. Rédige en chinois un court texte de présentation de ta commune (environ 60 à 80 caractères) en suivant ce plan : introduction (nom de la commune, situation), trois informations (habitants, marché, école ou hôpital), un projet pour l'avenir avec \zh{应该}, conclusion. Utilise au moins une fois \zh{因为} et une fois \zh{而且}.

\textbf{Sujet d'expression orale (exposé).} Présente oralement ta commune en 8 à 10 phrases, sans lire : introduction, trois informations, un projet, conclusion. Tu seras évalué(e) sur la justesse des tons, la correction des structures et la fluidité.

\textbf{Débat en binôme.} Avec un camarade, débats de la question : \zh{应该先建医院还是先建市场？} (Yīnggāi xiān jiàn yīyuàn háishi xiān jiàn shìchǎng ?). Chaque élève doit donner deux arguments avec \zh{因为} et \zh{而且}, puis réagir à l'avis de son camarade avec \zh{我同意} ou \zh{我不同意}.

\textbf{Jeu de rôle d'examen blanc.} Un élève joue un journaliste chinois, l'autre joue le délégué du maire de ta commune. Le journaliste pose au moins quatre questions (avec \zh{为什么}, \zh{什么时候}, \zh{还是}) ; le délégué répond en phrases complètes. L'enseignant évalue avec la grille officielle : prononciation et tons, grammaire, vocabulaire du thème, interaction.
\end{exercice}

\newpage

\coursec{Corrigés des exercices}

\begin{corrige}[Exercice 1 — QCM : le vocabulaire des collectivités locales]
\begin{enumerate}
\item Bonne réponse : \textbf{b}. \zh{市长} (shìzhǎng) = le maire. Le caractère \zh{市} signifie « ville, municipal » et \zh{长} « chef, responsable ». Attention à ne pas confondre avec \zh{部长} (bùzhǎng, ministre) ni avec \zh{校长} (xiàozhǎng, directeur d'école).
\item Bonne réponse : \textbf{a}. \zh{修路} (xiū lù) = construire / réparer une route. \zh{修} exprime la construction ou la réparation ; \zh{买} = acheter, \zh{看} = regarder.
\item Bonne réponse : \textbf{b}. \zh{居民} (jūmín) = les habitants d'une commune, d'un quartier. \zh{居} renvoie à « résider » (comme dans \zh{住} zhù, habiter).
\item Bonne réponse : \textbf{b}. \zh{我们的区有一个新市场。} (Wǒmen de qū yǒu yí ge xīn shìchǎng.) — « commune / arrondissement » se dit \zh{区} (qū) ; \zh{城市} (chéngshì) = ville, \zh{国家} (guójiā) = pays.
\end{enumerate}
\end{corrige}

\begin{corrige}[Exercice 2 — Vrai ou faux ?]
\begin{enumerate}
\item \textbf{Vrai.} \zh{应该} (yīnggāi) est un verbe modal qui exprime le devoir raisonnable ou le conseil : \zh{我们应该先建医院。} (Wǒmen yīnggāi xiān jiàn yīyuàn.) « Nous devrions d'abord construire un hôpital. » Il se place toujours avant le verbe principal.
\item \textbf{Faux.} \zh{因为……而且……} combine une cause (\zh{因为}) et un ajout (\zh{而且}) : c'est correct pour \emph{justifier} un avis, mais pour \emph{énumérer} plusieurs arguments dans un débat, on utilise plutôt \zh{第一……第二……} ou \zh{首先……然后……}. \zh{因为} introduit une justification, pas une énumération.
\item \textbf{Vrai.} \zh{还是} (háishi) introduit l'alternative dans une question : \zh{你喝茶还是喝咖啡？} (Nǐ hē chá háishi hē kāfēi ?). Ne pas confondre avec \zh{或者} (huòzhě), réservé aux phrases affirmatives : \zh{我想喝茶或者咖啡。}
\item \textbf{Vrai.} \zh{你呢？} (Nǐ ne ?) est la formule de relance la plus simple : « Et toi ? ». On peut aussi dire \zh{你觉得呢？} (Nǐ juéde ne ?) « Qu'en penses-tu ? » ou \zh{你的意见呢？} (Nǐ de yìjiàn ne ?) « Et ton avis ? ».
\end{enumerate}
\end{corrige}

\begin{corrige}[Exercice 3 — Complète le vocabulaire]
\begin{center}
\begin{tabularx}{\linewidth}{|X|X|X|X|}\hline
\rowcolor{popblueL} Caractères & Pinyin & Nature & Traduction \\ \hline
社区 & shèqū & nom & quartier, communauté \\ \hline
医院 & yīyuàn & nom & hôpital \\ \hline
建设 & jiànshè & verbe & construire, développer \\ \hline
项目 & xiàngmù & nom & projet \\ \hline
意见 & yìjiàn & nom & avis, opinion \\ \hline
\end{tabularx}
\end{center}
Points de vigilance : \zh{医} s'écrit avec la clé \zh{匚} (boîte) autour de \zh{矢} ; \zh{院} porte la clé de la colline \zh{阝} à gauche. \zh{建设} est un verbe (on peut dire \zh{建设我们的城市} jiànshè wǒmen de chéngshì, « développer notre ville »), tandis que \zh{项目} est un nom compté avec le classificateur \zh{个} : \zh{一个项目} (yí ge xiàngmù).
\end{corrige}

\begin{corrige}[Exercice 4 — Associe le pinyin aux caractères]
Correspondances : 1-b, 2-a, 3-d, 4-c, 5-f, 6-e.
\begin{itemize}
\item \zh{市长} shìzhǎng (le maire) ≠ \zh{市场} shìchǎng (le marché) : même premier caractère et même ton (3\up{e} ton) sur la deuxième syllabe, mais l'initiale change : \textbf{zh} (rétroflexe, langue recourbée) dans \zh{长} zhǎng contre \textbf{ch} (rétroflexe aspirée) dans \zh{场} chǎng. C'est le piège classique de la leçon : à l'oral, une mauvaise aspiration transforme le maire en marché !
\item \zh{所以} suǒyǐ (donc) : deux 3\up{e} tons ; à l'oral, par sandhi, le premier devient un 2\up{e} ton (on prononce suóyǐ).
\item \zh{可以} kěyǐ (pouvoir) : même règle de sandhi des 3\up{e} tons (on prononce kéyǐ).
\item \zh{政府} zhèngfǔ (gouvernement) : 4\up{e} ton + 3\up{e} ton.
\item \zh{公园} gōngyuán (parc) : 1\up{er} ton + 2\up{e} ton.
\end{itemize}
\end{corrige}

\begin{corrige}[Exercice 5 — Traduis en chinois]
\begin{enumerate}
\item \zh{我想介绍一下我的区。}\\(Wǒ xiǎng jièshào yíxià wǒ de qū.) — \zh{一下} adoucit le verbe ; « commune / arrondissement » = \zh{区}.
\item \zh{我们的区有很多居民。}\\(Wǒmen de qū yǒu hěn duō jūmín.) — Structure d'existence : lieu + \zh{有} + quantité + nom.
\item \zh{我觉得我们应该先建一个医院。}\\(Wǒ juéde wǒmen yīnggāi xiān jiàn yí ge yīyuàn.) — Ordre : \zh{觉得} + proposition complète ; \zh{先} (d'abord) se place juste avant le verbe.
\item \zh{因为人们常常生病，而且医院很远。}\\(Yīnwèi rénmen chángcháng shēngbìng, érqiě yīyuàn hěn yuǎn.) — \zh{而且} relie la deuxième cause à la première.
\item \zh{你同意我的意见吗？}\\(Nǐ tóngyì wǒ de yìjiàn ma ?) — Question en \zh{吗} ; on peut aussi dire plus simplement \zh{你同意吗？}
\end{enumerate}
\end{corrige}

\begin{corrige}[Exercice 6 — Questions alternatives avec 还是]
\begin{enumerate}
\item \zh{我们应该修路还是建公园？}\\(Wǒmen yīnggāi xiū lù háishi jiàn gōngyuán ?)
\item \zh{市长想建一个新市场还是一个新医院？}\\(Shìzhǎng xiǎng jiàn yí ge xīn shìchǎng háishi yí ge xīn yīyuàn ?)
\item \zh{我们应该先建学校还是先建一个医疗中心？}\\(Wǒmen yīnggāi xiān jiàn xuéxiào háishi xiān jiàn yí ge yīliáo zhōngxīn ?)
\item \zh{这个项目对年轻人有好处还是对老年人有好处？}\\(Zhè ge xiàngmù duì niánqīngrén yǒu hǎochù háishi duì lǎoniánrén yǒu hǎochù ?)
\item \zh{你想谈谈交通问题还是饮用水的问题？}\\(Nǐ xiǎng tántan jiāotōng wèntí háishi yǐnyòngshuǐ de wèntí ?)
\end{enumerate}
Règle : la question alternative se construit \cle{option A + 还是 + option B}, \textbf{sans} \zh{吗} final. Les deux options doivent être de même nature grammaticale (verbe + complément des deux côtés). Rappel : dans une phrase affirmative, « ou » se dit \zh{或者} : \zh{我们修路或者建公园。}
\end{corrige}

\begin{corrige}[Exercice 7 — Texte à trous]
Texte complété :
\begin{itemize}
\item \zh{我们的区有一个新\textbf{项目}：建一个大市场。} (xiàngmù)
\item \zh{\textbf{因为}旧的\textbf{市场}太小，所以很多\textbf{居民}不满意。} (yīnwèi / shìchǎng / jūmín)
\item \zh{我觉得我们\textbf{应该}听听大家的\textbf{意见}。} (yīnggāi / yìjiàn)
\end{itemize}
Traduction complète : « Notre arrondissement a un nouveau projet : construire un grand marché. Parce que l'ancien marché est trop petit, beaucoup d'habitants ne sont pas satisfaits. Je pense que nous devrions écouter l'avis de tout le monde. »
Justification : \zh{因为} ouvre la subordonnée de cause (réponse à « pourquoi ? ») ; \zh{应该} précède le verbe \zh{听} ; \zh{意见} est le complément de \zh{听听} (verbe redoublé = action brève, informelle : « écouter un peu, recueillir »).
\end{corrige}

\begin{corrige}[Exercice 8 — Remets le dialogue dans l'ordre]
Ordre correct : \textbf{2 — 4 — 5 — 1 — 3 — 6}.
\begin{enumerate}
\item \zh{我觉得应该先建一个新市场。} (Wǒ juéde yīnggāi xiān jiàn yí ge xīn shìchǎng.) — ouverture : avis avec \zh{觉得} + \zh{应该}.
\item \zh{为什么呢？} (Wèishénme ne ?) — demande de justification.
\item \zh{因为旧市场太小，而且不干净。你呢？} (Yīnwèi jiù shìchǎng tài xiǎo, érqiě bù gānjìng. Nǐ ne ?) — justification + relance.
\item \zh{我不同意。因为病人很多，而且医院太远了。} (Wǒ bù tóngyì. Yīnwèi bìngrén hěn duō, érqiě yīyuàn tài yuǎn le.) — désaccord argumenté.
\item \zh{你说得有道理，但是市场也很重要。} (Nǐ shuō de yǒu dàolǐ, dànshì shìchǎng yě hěn zhòngyào.) — concession polie : \zh{有道理} + \zh{但是}.
\item \zh{好吧，我们可以先建医院，再建市场。} (Hǎo ba, wǒmen kěyǐ xiān jiàn yīyuàn, zài jiàn shìchǎng.) — compromis final : \zh{先……再……} = d'abord… ensuite…
\end{enumerate}
Logique du dialogue : avis → pourquoi → argument → contre-argument → concession → compromis. C'est exactement le schéma à reproduire dans le débat de l'examen.
\end{corrige}

\begin{corrige}[Exercice 9 — Compréhension de l'écrit et questions de langue]
\textbf{Questions de compréhension :}
\begin{enumerate}
\item Parce qu'il n'y avait pas de grand marché dans l'arrondissement : les habitants devaient prendre une heure de voiture jusqu'au centre-ville pour faire leurs courses, ce qui était très peu pratique (\zh{很不方便} hěn bù fāngbiàn).
\item Ils ont décidé de construire un nouveau marché et de réparer deux routes (\zh{建一个新市场，修两条路}).
\item Leur argument est que le marché peut attendre, mais pas la santé des malades : \zh{市场可以等，病人的健康不能等。} (Shìchǎng kěyǐ děng, bìngrén de jiànkāng bù néng děng.)
\item Le mois prochain, il va organiser une réunion pour écouter les avis des habitants (\zh{听听居民们的意见} tīngting jūmínmen de yìjiàn).
\end{enumerate}
\textbf{Questions de langue :}
\begin{enumerate}
\item Phrase avec \zh{因为} relevée dans le texte : \zh{大家都很高兴，因为这个项目对商人和居民都有好处。} (Dàjiā dōu hěn gāoxìng, yīnwèi zhè ge xiàngmù duì shāngrén hé jūmín dōu yǒu hǎochù.) « Tout le monde était très content, parce que ce projet est profitable aux commerçants comme aux habitants. » Structure : \cle{因为 + cause} placée après la proposition principale, ce qui est possible en chinois comme en français.
\item Un mot de sens proche de \zh{高兴} (gāoxìng, content) : \zh{满意} (mǎnyì, satisfait) ; on accepte aussi \zh{快乐} (kuàilè, joyeux).
\item Phrase transformée au passé avec \zh{了} : \zh{居民们坐了一个小时的车去市中心。} (Jūmínmen zuò le yí ge xiǎoshí de chē qù shì zhōngxīn.) « Les habitants ont pris une heure de voiture pour aller au centre-ville. » — \zh{了} se place juste après le verbe \zh{坐} pour marquer l'action accomplie ; la durée \zh{一个小时} reste entre le verbe et son complément \zh{车}.
\end{enumerate}
Barème indicatif (sur 20) : compréhension 8 pts (2 par question), questions de langue 6 pts (2 par question), qualité de la traduction et de l'expression 6 pts.
\end{corrige}

\begin{corrige}[Exercice 10 — Thème et version]
\textbf{Thème :}
\begin{enumerate}
\item \zh{我们的区有很多项目：一个市场、一所学校和一个医院。}\\(Wǒmen de qū yǒu hěn duō xiàngmù : yí ge shìchǎng, yì suǒ xuéxiào hé yí ge yīyuàn.) — Classificateur des écoles et institutions : \zh{所} ; énumération avec la virgule chinoise \zh{、}.
\item \zh{我觉得应该先修路，因为路很不好。}\\(Wǒ juéde yīnggāi xiān xiū lù, yīnwèi lù hěn bù hǎo.) — « en mauvais état » se rend simplement par \zh{不好} ; \zh{先} avant le verbe \zh{修}.
\item \zh{市长想听听居民们的意见。}\\(Shìzhǎng xiǎng tīngting jūmínmen de yìjiàn.) — Redoublement \zh{听听} = « écouter un peu, recueillir » ; \zh{居民们} avec le suffixe de pluriel \zh{们}.
\end{enumerate}
\textbf{Version :} « Hier, notre arrondissement a organisé une réunion. Certains habitants ont dit qu'il fallait d'abord construire un marché, d'autres ont dit qu'il fallait d'abord construire une école. Le maire a dit : "Les avis de tous sont très importants, cherchons ensemble une solution." »
Points de vigilance : \zh{有的……有的……} (yǒude… yǒude…) = « les uns… les autres… » ; \zh{想办法} (xiǎng bànfǎ) = « trouver une solution » ; \zh{一起} (yìqǐ, ensemble) se place avant le verbe en chinois, mais se traduit en français selon la place naturelle (« cherchons ensemble »).
Barème indicatif (sur 20) : thème 10 pts (exactitude des caractères 4, structures 4, pinyin 2), version 10 pts (fidélité 6, qualité du français 4).
\end{corrige}

\begin{corrige}[Exercice 11 — Expression écrite et orale type Bac]
\textbf{Proposition de rédaction modèle (expression écrite) :}

\zh{我住在雅温得的埃索斯区。我们区有很多居民，还有一个大市场和一个新学校。以前，我们区没有医院，病人要去很远的地方看病，很不方便。我觉得我们应该先建一个医院，因为大家的健康最重要，而且医院对老人和孩子都有好处。我希望我们的区越来越好。}

(Wǒ zhù zài Yǎwēndé de Àisuǒsī qū. Wǒmen qū yǒu hěn duō jūmín, hái yǒu yí ge dà shìchǎng hé yí ge xīn xuéxiào. Yǐqián, wǒmen qū méiyǒu yīyuàn, bìngrén yào qù hěn yuǎn de dìfang kànbìng, hěn bù fāngbiàn. Wǒ juéde wǒmen yīnggāi xiān jiàn yí ge yīyuàn, yīnwèi dàjiā de jiànkāng zuì zhòngyào, érqiě yīyuàn duì lǎorén hé háizi dōu yǒu hǎochù. Wǒ xīwàng wǒmen de qū yuè lái yuè hǎo.)

Traduction : « J'habite dans l'arrondissement d'Essos, à Yaoundé. Notre arrondissement compte beaucoup d'habitants, et possède aussi un grand marché et une nouvelle école. Avant, il n'y avait pas d'hôpital : les malades devaient aller très loin pour se soigner, c'était très peu pratique. Je pense que nous devrions d'abord construire un hôpital, parce que la santé de tous est ce qu'il y a de plus important, et en plus un hôpital profite aux personnes âgées comme aux enfants. J'espère que notre arrondissement ira de mieux en mieux. »

Le texte respecte le plan imposé : introduction, trois informations, projet avec \zh{应该}, conclusion ; il contient bien \zh{因为} et \zh{而且}, ainsi que la structure \zh{越来越} (yuè lái yuè…, de plus en plus…).

\textbf{Barème indicatif (sur 20) :} respect du plan et de la consigne 4 pts ; correction des caractères 4 pts ; grammaire (ordre des mots, \zh{因为}, \zh{而且}, \zh{应该}) 6 pts ; vocabulaire du thème 4 pts ; cohérence et conclusion 2 pts.

\textbf{Points de vigilance :}
\begin{itemize}
\item Place de \zh{先} : toujours avant le verbe (\zh{先建}), jamais en fin de phrase comme le « d'abord » français.
\item \zh{因为……所以……} peuvent se combiner dans la même phrase (c'est correct en chinois), mais dans une phrase courte, un seul des deux suffit ; ne traduis jamais « parce que… donc… » mot à mot dans une phrase longue, cela alourdit.
\item Tons pièges à l'oral : \zh{市场} shìchǎng ≠ \zh{市长} shìzhǎng (initiales ch / zh) ; \zh{医院} yīyuàn (1\up{er} + 4\up{e} ton) ; \zh{意见} yìjiàn (4\up{e} + 4\up{e} ton, bien marquer la chute).
\item Pour le débat, chaque argument doit suivre le schéma : avis (\zh{我觉得}) → justification (\zh{因为}…\zh{而且}…) → réaction (\zh{我同意 / 我不同意}) → concession (\zh{你说得有道理，但是……}) → compromis (\zh{先……再……}).
\item Dans le jeu de rôle, les questions doivent être complètes : \zh{你们区什么时候建新市场？} (Nǐmen qū shénme shíhou jiàn xīn shìchǎng ?) « Quand votre arrondissement construira-t-il le nouveau marché ? » ; \zh{这个项目要花多少钱？} (Zhè ge xiàngmù yào huā duōshao qián ?) « Combien ce projet va-t-il coûter ? »
\end{itemize}
\end{corrige}

\newpage

\coursec{Fiche bilan}

\begin{popfigure}
\begin{tikzpicture}[mindmap, grow cyclic, every node/.style={concept, text=white, font=\small},
root concept/.append style={concept color=popblue, font=\bfseries},
level 1/.append style={level distance=3.2cm, sibling angle=51.4},
level 1 concept/.append style={font=\footnotesize}]
\node[root concept]{Collectivités locales}
child[concept color=poporange]{node[align=center]{Lexique\\ 市政府 shìzhèngfǔ\\ 社区 shèqū}}
child[concept color=popgreen]{node[align=center]{Présenter\\ 我住在……\\ 我们市有……}}
child[concept color=poppurple]{node[align=center]{Opiner\\ 我觉得……\\ 我认为……}}
child[concept color=poppink]{node[align=center]{Réagir\\ 我同意 / 不同意\\ 你说得对}}
child[concept color=popgold]{node[align=center]{Relancer\\ 你怎么看？\\ 还有呢？}}
child[concept color=popblue!70]{node[align=center]{Tons\\ 4 tons + neutre\\ mā má mǎ mà}}
child[concept color=popmuted]{node[align=center]{Débat\\ projet municipal\\ rôles : maire, habitant}};
\end{tikzpicture}
\legende{Carte mentale de la leçon : parler des collectivités locales en chinois.}
\end{popfigure}

\begin{bilan}
\textbf{1. Lexique clé à connaître par cœur}
\begin{center}
\begin{tabularx}{\linewidth}{|X|X|X|X|}\hline
\rowcolor{popblueL} Caractères & Pinyin & Nature & Traduction \\ \hline
市政府 & shìzhèngfǔ & n. & mairie, gouvernement municipal \\ \hline
社区 & shèqū & n. & communauté, quartier \\ \hline
居民 & jūmín & n. & habitant, résident \\ \hline
市长 & shìzhǎng & n. & maire \\ \hline
环境 & huánjìng & n. & environnement \\ \hline
建设 & jiànshè & n./v. & construction ; construire \\ \hline
意见 & yìjiàn & n. & avis, opinion \\ \hline
问题 & wèntí & n. & problème, question \\ \hline
解决 & jiějué & v. & résoudre \\ \hline
建议 & jiànyì & n./v. & suggestion ; suggérer \\ \hline
\end{tabularx}
\end{center}

\textbf{2. Structures à connaître par cœur}
\begin{center}
\begin{tabularx}{\linewidth}{|X|X|X|}\hline
\rowcolor{popblueL} Fonction & Structure + exemple & Traduction \\ \hline
Présenter & 我住在…… Wǒ zhù zài…… & J'habite à… \\ \hline
Décrire & 我们市有…… Wǒmen shì yǒu…… & Notre ville a… \\ \hline
Opiner & 我觉得…… Wǒ juéde…… & Je trouve que… \\ \hline
Opiner (fort) & 我认为…… Wǒ rènwéi…… & Je pense que… \\ \hline
Accord & 我同意你的看法。 Wǒ tóngyì nǐ de kànfǎ. & Je suis d'accord avec toi. \\ \hline
Désaccord & 我不太同意。 Wǒ bù tài tóngyì. & Je ne suis pas tout à fait d'accord. \\ \hline
Relancer & 你怎么看？ Nǐ zěnme kàn ? & Qu'en penses-tu ? \\ \hline
Suggérer & 我建议…… Wǒ jiànyì…… & Je suggère… \\ \hline
\end{tabularx}
\end{center}

\textbf{3. Méthodes express}
\begin{itemize}
\item \cle{Exposé simple} : 3 étapes — saluer (\zh{大家好！}), présenter (\zh{我给大家介绍一下……}), conclure (\zh{谢谢大家！}).
\item \cle{Débat} : donner son avis avec \zh{我觉得}, réagir avec \zh{我同意 / 我不同意}, relancer avec \zh{你怎么看？}.
\item \cle{Tons} : articuler lentement, exagérer les tons à l'entraînement ; attention au 3\textsuperscript{e} ton devant un autre 3\textsuperscript{e} ton (\zh{nǐ hǎo} se prononce \emph{ní hǎo}).
\item \cle{Prise de parole} : phrases courtes, regarder l'auditoire, ne jamais lire mot à mot.
\end{itemize}
\end{bilan}

\begin{aretenir}[Checklist avant le Bac]
\begin{itemize}
\item[$\square$] Je connais le lexique des collectivités locales (\zh{市政府、社区、居民、市长}).
\item[$\square$] Je sais me présenter et présenter ma commune en 3 phrases.
\item[$\square$] Je maîtrise \zh{我觉得} et \zh{我认为} pour donner mon avis.
\item[$\square$] Je sais exprimer l'accord (\zh{我同意}) et le désaccord poli (\zh{我不太同意}).
\item[$\square$] Je sais relancer un débat avec \zh{你怎么看？} et \zh{还有呢？}.
\item[$\square$] Je prononce correctement les quatre tons et le ton neutre.
\item[$\square$] J'applique la règle du 3\textsuperscript{e} ton devant un 3\textsuperscript{e} ton.
\item[$\square$] Je sais ouvrir un exposé par \zh{大家好！} et le clore par \zh{谢谢大家！}.
\item[$\square$] Je sais suggérer une solution avec \zh{我建议……}.
\item[$\square$] Je peux tenir un rôle (maire, habitant, journaliste) dans un jeu de rôle.
\item[$\square$] Je parle en phrases courtes et correctes, sans traduire mot à mot du français.
\item[$\square$] J'ai répété le dialogue modèle à voix haute au moins trois fois.
\end{itemize}
\end{aretenir}

\findecours

\end{document}
